% -----------------------------------------------------------
\section{1833}
% -----------------------------------------------------------
	\label{EMS-1833}
	
	% ----------------------
	\subsection{Joseph Smith}
	% ----------------------
		\label{EMS-1833-JS}
		
		% ----------------------
		\subsubsection{Introduction to 1833}
		% ----------------------
		
			sdf
			
		% -----
		\subsubsection{Revelation, 3 January 1833 [D\&C 88:127--137]}
		% -----
			\label{EMS-1833-JS-3-JAN-1833}
			% \label{EMS-1830-JS-DAY-3LETTERMONTH-YEAR}
		
			\begin{description}
				\item[Print Sources]
			\end{description}

				\begin{tabular}{p{0.5in}p{1in}p{0.5in}p{1in}}
				JSP	 	& D2:346--348		& JSR 		& 229--230 \\
				DC	 	& Section 88 		& HC 		& 1:302--312 \\
				HDDC 	& 1127--1168		& TCHC 		& 1:218--224 \\
				\end{tabular}
				% HDDC = woodford's DC diss; JSR = Marquardt's JS Revelations; TCHC = Vogel HC
			
			\begin{description}
				\item[Online Access] \href{https://www.josephsmithpapers.org/paper-summary/revelation-3-january-1833-dc-88127-137/1}{Here}
				% \item[Analysis] \hyperref[DC88]{Here}
			\end{description}
			
			\begin{description}
				\item[Background]
			\end{description}
			
				% It might also be useful to give a two sentence context of the period: e.g., "This speech was given in Nauvoo to a small congregation one month prior to the destruction of the Nauvoo Expositor. These and other tensions may have been on the prophet's mind when he gave the speech."
			
			\begin{description}
				\item[Original Source]
			\end{description}
			
				% brief description of original source material, with reference to JSP or somewhere else for more info
				% Background material: it might be helpful to have a note that says whether a given document was presented as a speech, was written in a journal, etc.
				% Also a comment here about how reliable the source might be, or what shape it is in, if there are large omissions or parts that are hard to understand, and so on.
			
			\begin{description}
				\item[Text]
			\end{description}
			
				Kirtland January 3\supers{d.} 1833. \uline{Revelation} given to \sout{organize} for a patern \&c
				
				The order of the house, of God prepared for the presedency, and instruction, in all things, that is expedient for the officers; or in other words, them who are called to the ministry in the Church, begining at the high Priests, even down to the deacon, and this shall be the order of the house, he that is appointed, to be a teacher, shall be found standing in his place, which shall be \sout{appointed} prepared, for him, in the house of God, in a place that the congregation, in the house may hear his words, correctly and distinctly; not with loud speach; and he cometh into the house of God, (for he should be first in the house, behold this is beautiful, that he may be an example) let him offer himself in prayer upon his knees, before God, in token of the everlasting covenant, and when any shall, come in after him, let the teacher arise, and with uplifted hands to heaven, yea even directly, salute his brother, or brethren, with these words saying, art thou a brother, or brethren, I salute you in the name of the Lord Jesus Christ, in tocen [token] of the everlasting covenant, in which covenant, I receive you to fellowship, in a determination, that is fixed immovable, and unchangable, to be your friend and brother, through the grace of God, in the bonds of Love, to walk in all the commandments, of God, blameless, in thanksgiving for ever, and ever; Amen. and he that cometh, in, and is a brother, or brethren shall salute the teacher with uplifted hands to heaven, with this same prayer, and covenant or by saying amen; in token of the same, Behold verily I say unto you, this [is?] a sample, unto you for a salutation, to one another, in the house of God, and ye are called to do this by prayer, and thanksgiving, as the spirit shall give utterance in all your doings, in the house of the Lord, that it may become a sanctury, a tabernacle, of the holy spirit to your, edefication Amen
				
				Given by Joseph the seer, and writen by Frederick [G. Williams] assistant scribe and councellor------
				
		% -----
		\subsubsection{Letter to Noah C. Saxton, 4 January 1833}
		% -----
			\label{EMS-1833-JS-4-JAN-1833}
			% \label{EMS-1830-JS-DAY-3LETTERMONTH-YEAR}
		
			\begin{description}
				\item[Print Sources]
			\end{description}

				\begin{tabular}{p{0.5in}p{1in}p{0.5in}p{1in}}
				JSP	 	& D2:348--355		& JSR 		& --- \\
				DC	 	& --- 				& HC 		& 1:312--316 \\
				HDDC 	& ---				& TCHC 		& 1:224--227 \\
				TPJS 	& 13--18			& 	 		&  \\
				\end{tabular}
				% HDDC = woodford's DC diss; JSR = Marquardt's JS Revelations; TCHC = Vogel HC
			
			\begin{description}
				\item[Online Access] \href{https://www.josephsmithpapers.org/paper-summary/letter-to-noah-c-saxton-4-january-1833/1}{Here}
				% \item[Analysis] \hyperref[XX]{Here}
			\end{description}
			
			\begin{description}
				\item[Background]
			\end{description}
			
				% It might also be useful to give a two sentence context of the period: e.g., "This speech was given in Nauvoo to a small congregation one month prior to the destruction of the Nauvoo Expositor. These and other tensions may have been on the prophet's mind when he gave the speech."
			
			\begin{description}
				\item[Original Source]
			\end{description}
			
				% brief description of original source material, with reference to JSP or somewhere else for more info
				% Background material: it might be helpful to have a note that says whether a given document was presented as a speech, was written in a journal, etc.
				% Also a comment here about how reliable the source might be, or what shape it is in, if there are large omissions or parts that are hard to understand, and so on.
			
			\begin{description}
				\item[Text]
			\end{description}
			
				\begin{tightcenter}
						\hypertarget{EMS-1833-JS-4-JAN-1833-a}%
					Kirtland
						\marginnote{a}
					4\supers{th.} Jan\supersu{y}\supers{.} 1833---
				\end{tightcenter}
				
				M\supersu{r}\supers{.} Editor Sir,
				
				Considering the Liberal prisciples [principles] upon which your interesting and valuable paper is published and myself being a subscriber and feeling a deep intrist in the cause of Zion and in the happiness of my brethren of mankind I cheerfully take up my pen to contribute my mite at this every [very] interesting and important period
				
					\hypertarget{EMS-1833-JS-4-JAN-1833-b}%
				For
					\marginnote{b}%
				some length of time I have been car[e]fully viewing the state of things as now appear throug[h]out our christian Land and have looked at it with feelings of the most painful anxiety while upon the one hand beholding the manifeste\sout{d} withdrawal of Gods holy Spirit and the vail of stupidity which seems to be drawn over the hearts of the people and upon the other hand beholding the Judgments of God that have swept and are still sweeping hundreds and thousands of our race (and I fear unprepared) down to the shades of death with this solemn and alarming fact before me I am led to exclaim [``]O that my head were waters and mine ey[e]s a fountain of tears that I might weep day and night \&c,''
					\hypertarget{EMS-1833-JS-4-JAN-1833-c}%
				I
					\marginnote{c}%
				think that it is high time for a christian world to awake out of sleep and cry mightely to that God day and night whose anger we have Justly incured. Are not these things a sufficient stimulant to arouse the faculties and call forth the energies of evry man woman and child that poseses feeling of sympathy for his fellow[s] or that is in any degree endeared to the bud[d]ing cause of our glorious Lord; I leave an inteligent community to answer this important question with a confession that this is what has caused me to overlook my own inability and expose my weakness to a learned world but trusting in that God. who has said these things are hid from the wise and prudent and reve[a]led unto babes I step forth into the field to tell you what the Lord is doing and what you must do to enjoy the smiles of your saviour in these last day[s]------
					\hypertarget{EMS-1833-JS-4-JAN-1833-d}%
				The
					\marginnote{d}%
				time has at last \sout{come} arived when the God\sout{s} of Abraham of Isaac and of Jacob has set his hand again the seccond time to recover the remnants of his people which have been left from Assyria, and from Egypt and from Pathros \&.c. and from the Islands of the sea and with them to bring in the fulness of the Gentiles and establish that covenant with them which was promised when their sins should be taken away. see Romans 11, 25, 26, \& 27 and also Jeremiah 31. 31, 32, \& 33,
					\hypertarget{EMS-1833-JS-4-JAN-1833-e}%
				This
					\marginnote{e}%
				covenant has never been established with the house of Isreal nor with th[e] house of Judah for it requires two parties to make a covenant and those two parties must be agreed or no covenant can be made. Christ in the days of his flesh proposed to make a covenant with them but they rejected him and his proposals and in consequence thereof they were broken off and no covenant was made with them at that time but their unbelief has not rendered the promise of God of none effect;
					\hypertarget{EMS-1833-JS-4-JAN-1833-f}%
				no,
					\marginnote{f}%
				for there was another day limited in David which was the day of his power and then his people \uline{Isreal}, should be a willing people and he would write his laws in their hearts and print them in their thoughts their sins and their eniquities he would remember no more, Thus after this chosen family had rejected Christ and his proposals the heralds of salvation said to them. ``lo \sout{I} we turn <un>to the gentiles,'' and the gentiles received the covenant and were grafted in from whence the chosen family were broken off
					\hypertarget{EMS-1833-JS-4-JAN-1833-g}%
				but
					\marginnote{g}%
				the Gentiles have not continued in the goodness of God but have departed from the faith that was once delivered to the saints and have broken the \sout{everlasting} covenant in which their fathers were established see Isaiah 24\supers{th} 5\supers{th.} and have become high minded and have not feared therefore but few of them will be gathered with the \uline{chosen} \uline{family} Has not the pride highmindedness and unbelief of the Gentiles provoked the holy one of Israel to withdraw his holy spirit from them and send forth his Judgments to scourge them for their wickedness; this is certianly the case, Christ said to his deciples Mark 16, 17 \& 18 that these signs should follow them that believe;
					\hypertarget{EMS-1833-JS-4-JAN-1833-h}%
				In
					\marginnote{h}%
				my name shall they cast out Devils they shall speap [speak] with new tongues they shall take up serpants and if they drink any deadly thing it shall not hurt them they shall lay hands on the sick and they shall recover, and also in connection with this read 1 Corinthians 12 Chapt, By the foregoing testamonies or through the glass of the foregoing testamonies we may look at the Christian world and see the apostacy there has been from the Apostolic platform, and who can look at this, \sout{and} and not exclaim in the language of Isaiah, [``]the earth is defiled under the inhabitants thereof because they have transgressed the Laws; changed the ordinances and broken the everlasting covenant''
				
					\hypertarget{EMS-1833-JS-4-JAN-1833-i}%
				The
					\marginnote{i}%
				plain fact is this, the power of God begins to fall upon the Nations, and the light of the latter day glory begins to break forth through the dark atmosphere of sectarian wickedness and their iniquity rools [rolls] up into view and the Nations of the Gentiles are like the waves of the sea casting up mire and dirt or all in commotion and they hastily are preparing to act the part allotted them when the Lord rebukes the nations, when he shall rule them with a rod of iron \& break them in peaces like a potters vessel,
					\hypertarget{EMS-1833-JS-4-JAN-1833-j}%
				The
					\marginnote{j}%
				Lord has declared to his servants some Eighteen months since that he was then withdrawing his spirit from the earth, and we can see that such is the fact for not only the churches are dwindling away, but there are no convers[i]ons, or but very few, and this is not all, the governments of the earth are thrown into confusion \& division, and distruction to the eye of the spiritual beholder seemes to be writen by the finger of an invisable hand in Large capitals upon almost evry thing we behold------
				
					\hypertarget{EMS-1833-JS-4-JAN-1833-k}%
				And
					\marginnote{k}%
				now what remains to be done under circumstan[c]es like these, I will proce[e]d to tell you what the Lord requires of all people high and Low, rich and poor, male and female, ministers \& people professors of religeon, and nonproffessors in order that they may enjoy the holy spirit of God to a fulness, and escape the Judgments of God which are almost ready to burst upon the nations of the earth--- Repent of all your sins and be baptized in water for the remission of them, in the name of the father, and of the son, and of the Holy Ghost, and receive the ordinance of the laying on of the hands of him who is ordained and sealed unto this power, that ye may receive the holy spirit of God, and this according to the holy scriptures, and of the Book of Mormon; and the only way that man can enter into the Celestial kingdom.
					\hypertarget{EMS-1833-JS-4-JAN-1833-l}%
				These
					\marginnote{l}%
				are the requesitions of the new Covenant or first principles of of the Gospel of Christ; then add to you[r] faith virtue and to virtue knowledge and to knowledge temperance, and to temperance patience, and to patience, brotherly kindness and to brotherly kindness charity (or Love) and if these things be in you and abound, they make you to be neither baran nor unfruitful in the knowledge of our Lord Jesus Christ------
				
					\hypertarget{EMS-1833-JS-4-JAN-1833-m}%
				The
					\marginnote{m}%
				Book of Mormon is a reccord of the forefathers of our western Tribes of Indians, having been found through the ministration of an holy Angel translated into our own Language by the gift and power of God, after having been hid up in the earth for the last fourteen hundred years containing the word of God, which was delivered unto them, By it we learn that our western tribes of Indians are desendants from that Joseph that was sold into Egypt, and that the Land of America is a promised land unto them, and unto it all the tribes of Israel will come. with as many of the gentiles as shall comply with the requesitions of the new co[v]enant.
					\hypertarget{EMS-1833-JS-4-JAN-1833-n}%
				But
					\marginnote{n}%
				the tribe of Judah will return to old Jerusalem, The City, of Zion, spoken of by David in the 102 Psalm will be built upon the Land of America and the ransomed of the Lord shall return and come to it with songs and everlasting joy upon their heads, and then they will be delivered from the overflowing scourge that shall pass through the Land But Juda\sout{h} shall obtain deliverence at Jerusalem see Joel 2. 32. Isaiah 26, 20 \& 21, Jer. 31, 12, Psalm 50. 5, Ezekiel 34, 11, 12 \& 13, These are testamonies that the good Shepherd will put forth his own sheep and Lead them out from all nations where they have been scattered in a cloudy and dark day, to Zion and to Jerusalem beside many more testamonies which might be brought------
				
					\hypertarget{EMS-1833-JS-4-JAN-1833-o}%
				And
					\marginnote{o}%
				now I am prepared to say by the authority of Jesus Christ, that not many years shall pass away before the United States shall present such a scene of \uline{bloodshed} as has not a parallel in the hystory of our nation pestalence hail famine and earthquake will sweep the wicked of\sout{f} this generation from off the face of this Land to open and prepare the way for the return of the lost tribes of Israel from the north country--- The people of the Lord, those who have complied with the requsitions of the new covenant have already commenced gathering togethe[r] to Zion which is in the State of Missouri.
					\hypertarget{EMS-1833-JS-4-JAN-1833-p}%
				Therefore
					\marginnote{p}%
				I declare unto you the warning which the lord has commanded me to declare unto this generation, rembring [remembering] that the eyes of my maker are upon me and that to him I am accountabl for evry word I say wishing nothing worse to my fellow men then their eternal salvation therefore fear God, and give glory to him for the hour of his Judgment is come,
					\hypertarget{EMS-1833-JS-4-JAN-1833-q}%
				<Repent
					\marginnote{q}%
				ye> Repent, ye and imbrace the everlasting Covenant and flee to Zion before the overflowing scourge overtake you, For there are those now living upon the earth whose eyes shall not be closed in death until they see all these things which I have spoken fulfilled \uline{Rem}[\uline{em}]\uline{ber} these things, call upon the Lord while he is near and seek him while he may be found is the exhortation of your unworthy servant
				
				\begin{tightright}
				Joseph Smith J\uline{r}
				\end{tightright}
				
				<To N. E. Sextan [Noah C. Saxton] Rochester N Y.>
				
		% -----
		\subsubsection{Revelation, 5 January 1833}
		% -----
			\label{EMS-1833-JS-5-JAN-1833}
			% \label{EMS-1830-JS-DAY-3LETTERMONTH-YEAR}
		
			\begin{description}
				\item[Print Sources]
			\end{description}

				\begin{tabular}{p{0.5in}p{1in}p{0.5in}p{1in}}
				JSP	 	& D2:356--361		& JSR 		& 231 \\
				DC	 	& --- 				& HC 		& --- \\
				HDDC 	& ---				& TCHC 		& --- \\
				\end{tabular}
				% HDDC = woodford's DC diss; JSR = Marquardt's JS Revelations; TCHC = Vogel HC
			
			\begin{description}
				\item[Online Access] \href{https://www.josephsmithpapers.org/paper-summary/revelation-5-january-1833/1}{Here}
				% \item[Analysis] \hyperref[XX]{Here}
			\end{description}
			
			\begin{description}
				\item[Background]
			\end{description}
			
				% It might also be useful to give a two sentence context of the period: e.g., "This speech was given in Nauvoo to a small congregation one month prior to the destruction of the Nauvoo Expositor. These and other tensions may have been on the prophet's mind when he gave the speech."
			
			\begin{description}
				\item[Original Source]
			\end{description}
			
				% brief description of original source material, with reference to JSP or somewhere else for more info
				% Background material: it might be helpful to have a note that says whether a given document was presented as a speech, was written in a journal, etc.
				% Also a comment here about how reliable the source might be, or what shape it is in, if there are large omissions or parts that are hard to understand, and so on.
			
			\begin{description}
				\item[Text]
			\end{description}
				
				\begin{tightcenter}
				Januery 5
				\end{tightcenter}
				
				Behold I Say unto you my Servent Frederick [G. Williams], Listen to the word of Jesu Chrit [Jesus Christ] your Lord and your Redeemer thou hast desired of me to know which would be the most worth unto you, behold blessed art tho[u] for this thing. now I say unto you, my Servant Joseph is called to do a great work and hath need that he may do the work of translation for the salvation of souls--- Verily verily I say unto you thou art calld to be a Councillor \& scribe unto my servant Joseph Let thy farm be consecratd f[o]r bringing forth of the revelations and tho[u] shalt be blessed and lifted up at the last day even so Amen
				
		% -----
		\subsubsection{Letter to William W. Phelps, 11 January 1833}
		% -----
			\label{EMS-1833-JS-11-JAN-1833}
			% \label{EMS-1830-JS-DAY-3LETTERMONTH-YEAR}
		
			\begin{description}
				\item[Print Sources]
			\end{description}

				\begin{tabular}{p{0.5in}p{1in}p{0.5in}p{1in}}
				JSP	 	& D2:364--368		& JSR 		& --- \\
				DC	 	& --- 				& HC 		& 1:316--317 \\
				HDDC 	& ---				& TCHC 		& 1:227--229 \\
				TPJS 	& 18--19			& 	 		&  \\
				\end{tabular}
				% HDDC = woodford's DC diss; JSR = Marquardt's JS Revelations; TCHC = Vogel HC
			
			\begin{description}
				\item[Online Access] \href{https://www.josephsmithpapers.org/paper-summary/letter-to-william-w-phelps-11-january-1833/1}{Here}
				% \item[Analysis] \hyperref[XX]{Here}
			\end{description}
			
			\begin{description}
				\item[Background]
			\end{description}
			
				% It might also be useful to give a two sentence context of the period: e.g., "This speech was given in Nauvoo to a small congregation one month prior to the destruction of the Nauvoo Expositor. These and other tensions may have been on the prophet's mind when he gave the speech."
			
			\begin{description}
				\item[Original Source]
			\end{description}
			
				% brief description of original source material, with reference to JSP or somewhere else for more info
				% Background material: it might be helpful to have a note that says whether a given document was presented as a speech, was written in a journal, etc.
				% Also a comment here about how reliable the source might be, or what shape it is in, if there are large omissions or parts that are hard to understand, and so on.
			
			\begin{description}
				\item[Text]
			\end{description}
				
				\begin{tightcenter}
						\hypertarget{EMS-1833-JS-11-JAN-1833-a}%
					Kirtland
						\marginnote{a}%
					1833 January 11---%
						\footnote{%
							\label{EMS-1833-JS-11-JAN-1833-sanction}%
							Note that this letter was ``sanctioned'' by a conference of the church, in JSP Minutes, 13–14 January 1833.
							}
					
					Brother William [W. Phelps]
				\end{tightcenter}
				
				I Send you the Olieve leaf which we have plucked from the tree of Paradise, the Lords message of peace to us, for though our Brethren in Zion, indulge in feelings towards us, which \sout{is} are not according to the requirements of the new covenant yet we have the satisfaction of knowing that the Lord approves of us \& has accepted us, \& established his name in kirtland for the salvation of the nations, for the Lord will have a place \sout{for} from whence his word will go forth in these last days in purity, for if Zion, will not purify herself so as to be approved of in all things in his sight he will seek--- another people for his work will go on untill Isreal is gathered \& they who will not hear his voice must expect to feel his wrath,
					\hypertarget{EMS-1833-JS-11-JAN-1833-b}%
				Let
					\marginnote{b}%
				me say \sout{un}to you, seek to purefy yourselves, \& also all the inhabitants of Zion lest the Lords anger be kindled to fierceness, repent, repent, is the voice of God, to Zion, \& yet strange as it may appear, yet it is true mankind will presist in self Justification until all their eniquity is exposed \& their character past being redeemed, \& that which is treasured up in their hearts be exposed to the gaze of mankind, I say to you--- (\& what I say to you, I say to all) hear the <warning.> voice of God lest Zion fall, \& the Lord \sout{swore} swear in his wrath the inhabiteints of Zion shall not enter into my rest,
					\hypertarget{EMS-1833-JS-11-JAN-1833-c}%
				The
					\marginnote{c}%
				Brethren in Kirtland pray for you unceasingly, for knowing the terrors of the Lord, they greatly fear for you; you will see that the Lord commanded us in Kirtland to build an house of God, \& establish a school for the Prophets, this is the word of the Lord to us, \& we must--- yea the Lord helping us we will obey, as on conditions of our obedience, he has promised <us> great things, yea <even> a visit from the heavens to honor us with his own presence,
					\hypertarget{EMS-1833-JS-11-JAN-1833-d}%
				we
					\marginnote{d}%
				greatly fear before the Lord lest we should fail of this great honor which our master proposes to confer on us, we are seeking for humility \& great faith lest we be ashamed in his presence, our hearts are greatly greaved at the spirit which is breathed both in your letter \& that of Bro G---s [Sidney Gilbert] the wery spirit which is wasting the strength of Zion like a pestalence, and if it is not detected \& driven from you it will ripen Zion for the threatened Judgments of God, remember God sees the secret springs of human action, \& knows the hearts of all living,
					\hypertarget{EMS-1833-JS-11-JAN-1833-e}%
				Br
					\marginnote{e}%
				suffer us to speak plainly for God has respect to the feelings of his saints \& he will not let them be tantalized with impunity tell Br. G---t [Gilbert] that low insinuations God hates, but he rejoices in an honest heart and knows better who is guilty than he does, we send him this worning [warning] voice \& let him fear greatly for himself lest a worse thing overtake him, all we can say by way of conclusion is, if the fountain of our tears are not dried up we will <still> weep for Zion, this from your brother who trembles greatly for Zion,--- and for the wrath of heaven which awaits her if she repent not,------
					\hypertarget{EMS-1833-JS-11-JAN-1833-f}%
				PS.
					\marginnote{f}%
				I am not in the habit of crying peace, when there is no peace, and knowing the th[r]eatened Judgments of God, I say Wo,--- unto them who are at ease in Zion fearfulnes will speedily lay hold of the hypocrite, I did not expect that you had lost Commandments, but thought from your letters you had neglected to read them, otherwise you would not have writen as you did, it is in vain to try to hide a bad spirit from the eyes of them who are spiritual for it will shewe itself in speaking \& in writing as well as all our other conduct,
					\hypertarget{EMS-1833-JS-11-JAN-1833-g}%
				it
					\marginnote{g}%
				is also useless to mak[e] great pretentions when the heart is not right before God, for God looks at the heart, and where the heart is not right the Lord will expose it to the view of his faithful saints, we wish you to render the Star as interesting as possable by setting forth the rise progress and faith of the church, as well as the doctrine for if you do not render it more interesting than at present it will fall, and the church suffer a great Loss thereby------
				
				\begin{tightright}
				Joseph Smith Jr
				\end{tightright}
				
		% -----
		\subsubsection{Letter to Edward Partridge and Others, 14 January 1833}
		% -----
			\label{EMS-1833-JS-14-JAN-1833}
			% \label{EMS-1830-JS-DAY-3LETTERMONTH-YEAR}
		
			\begin{description}
				\item[Print Sources]
			\end{description}

				\begin{tabular}{p{0.5in}p{1in}p{0.5in}p{1in}}
				JSP	 	& D2:371--378		& JSR 		& --- \\
				DC	 	& --- 				& HC 		& 1:317--321 \\
				HDDC 	& ---				& TCHC 		& 1:229--233 \\
				\end{tabular}
				% HDDC = woodford's DC diss; JSR = Marquardt's JS Revelations; TCHC = Vogel HC
			
			\begin{description}
				\item[Online Access] \href{https://www.josephsmithpapers.org/paper-summary/letter-to-edward-partridge-and-others-14-january-1833/1}{Here}
				% \item[Analysis] \hyperref[XX]{Here}
			\end{description}
			
			\begin{description}
				\item[Background]
			\end{description}
			
				% It might also be useful to give a two sentence context of the period: e.g., "This speech was given in Nauvoo to a small congregation one month prior to the destruction of the Nauvoo Expositor. These and other tensions may have been on the prophet's mind when he gave the speech."
			
			\begin{description}
				\item[Original Source]
			\end{description}
			
				Note that this letter was written by Orson Hyde and Hyrum Smith, but was approved by a conference of which JS was a part (the end of the letter mention this).
			
				% brief description of original source material, with reference to JSP or somewhere else for more info
				% Background material: it might be helpful to have a note that says whether a given document was presented as a speech, was written in a journal, etc.
				% Also a comment here about how reliable the source might be, or what shape it is in, if there are large omissions or parts that are hard to understand, and so on.
			
			\begin{description}
				\item[Text]
			\end{description}
			
				\begin{tightright}
						\hypertarget{EMS-1833-JS-14-JAN-1833-a}%
					Kirtland
						\marginnote{a}
					Mills Geauga Co Ohio Jan\supersu{y}\supers{.} 14--- 1833
				\end{tightright}
				
				From a conference of 12 High Priests to the Bishop his councel and the inhabitants of Zion, Orson Hyde and Hyram [Hyrum] Smith being appointed by the said conference to write this Epistle in obedience to the commandment given the 22 \& 23 of Sept last which says ``But verily I say unto \sout{you} all those to whom the kingdom has been given from you it must be preached unto \sout{all} them that they shall repent of their former evil works for they are to be upbraided for their evil hearts of unbelief and your brethren in Zion for their rebellion against you at the time I sent you.''
					\hypertarget{EMS-1833-JS-14-JAN-1833-b}%
				Bro
					\marginnote{b}%
				Joseph \& certain <others> have writen to you upon this all important subject, but you have <never> been apprised of these things by the united voice of a conference of those high priests that were present at the time this commandment was given, we therefore Orson \& Hyram the committe appointed by said conference to write this Epistle having received the prayers of s\supers{d} conference that we might be enabled to write the mind \& will of God upon this subject,
					\hypertarget{EMS-1833-JS-14-JAN-1833-c}%
				now
					\marginnote{c}%
				take up our pen to address you in the name of the conference, relying upon the arm of the \uline{Great} \uline{Head} of the Church, in the commandment above alluded to the Children of Zion were all, yea even evry one under condemnation, and were to remain in that state until they repented and remembred the new covenant even the Book of Mormon and the former commandments which the Lord had given them, not only to say but to do them and bring forth fruit meet for ther Fathers Kingdom otherwise.
					\hypertarget{EMS-1833-JS-14-JAN-1833-d}%
				there
					\marginnote{d}%
				remaineth a scourge \& a judgment to be poured out upon the children of Zion, for shall the children of the Kingdom pollute my holy land I say unto you nay--- The answers received from those letters which have been sent to you upon this subject have failed to bring to us that satisfactory confession and acknowledgment which the spirit of our Master requires--- we therefore feeling a deep intrest for Zion \& knowing the judgments of God that will come upon her except she repent; resort to these \uline{last}, and most efficient means in our power to bring her to a sense of her standing before the most High,
					\hypertarget{EMS-1833-JS-14-JAN-1833-e}%
				At
					\marginnote{e}%
				the time Joseph Sidney [Rigdon] \& Newel [K. Whitney] left Zion, all matters of hardness and misunderstanding were settled and buried (as they supposed) and you gave them the hand of fellowship but afterwards you brought up all these things again in a sensorious spirit accusing Bro Joseph in rather an indirect way \sout{in} of seeking after Monarchal power and authority, this came to us in Bro Carrolls [John Corrill's] letter of June 2\supers{d.},
					\hypertarget{EMS-1833-JS-14-JAN-1833-f}%
				We
					\marginnote{f}%
				are sensable that this is not the thing Bro J is seeking after, but to magnify the high office and calling whereunto he has been called and appointed by the command of God and the united voice of this Church, It might not be a miss for you to call to mind the circumstances of the Nephites and the Children of Israel rising up against their prophets and accusing them of seeking after Kingly power \&c--- 
					\hypertarget{EMS-1833-JS-14-JAN-1833-g}%
				and
					\marginnote{g}%
				see what befel them and take warning before it is to[o] late Bro [Sidney] Gilberts letter of Dec 10\supers{th.} has been received and read attentively, and the low, dark, \& blind insinuations which were in it were not received by us as from the fountain of light, though his claims and pretentions to holiness were great, we are not unwilling to be chastened or rebuked for our faults but we want to receive it in Language that we can understand, as Nathan said to David, Thou art the man
					\hypertarget{EMS-1833-JS-14-JAN-1833-h}%
				We
					\marginnote{h}%
				are aware that Bro Gs is doing much, and a multitude of business on hand but let him purge out all the old leaven and do his business in the spirit of the Lord. and then the Lord will bless him <otherwise the frown of the Lord will remain upon him---> There is ma[n]ifestly an uneassness in Bro G, and a fearfulness that God will not provide for his saints in their last days and these fears lead him on to covitousness, This ought not so to be, but let him do just as the Lord has commanded him and then the Lord will open his coffers, and his wants will be liberally supplied, 
					\hypertarget{EMS-1833-JS-14-JAN-1833-i}%
				But
					\marginnote{i}%
				if this uneasy covetous disposition be cherished by him the Lord will bring him to poverty shame, and disgrace, Bro [William W.] Phelps letter is also received of \uline{Dec}15\supersu{th} and carefully read, and it betrays a \uline{lightness} of spirit that ill becomes a man plased in the important and responsable station that he is placed in;--- If you have fat beef and potatoes eat them in singleness of heart. and boast not yourselves in these things think not Brethren that we make a man an offender for a word this is not the case, but we want to see a spirit in Zion by which the Lord will build it up, that is the plain solem and pure spirit of Christ;
					\hypertarget{EMS-1833-JS-14-JAN-1833-j}%
				Bro
					\marginnote{j}%
				Phelps requested in his last letter that Bro Joseph should come to Zion, but we say that Br J. will not settle in Zion until she repent and purify herself \& abide by the new covenant, and remember the commandments which have been given her, to do them as well as say them You may think it strange that we manifest no cheerfulness of heart upon the reception of your letters, you may think that our minds are pregudiced so much that we can see no good that comes from you, but rest assured brethren that this is not the case,
					\hypertarget{EMS-1833-JS-14-JAN-1833-k}%
				We
					\marginnote{k}%
				have the best of feelings, and feelings of the greatest anxiety for the welfare of Zion we feel more like weeping over Zion than we do like rejoicing over her, for we know the judgments of God that hang over her and will fall upon her except she repent and purify herself before the Lord and put away from her evry foul spirit---
					\hypertarget{EMS-1833-JS-14-JAN-1833-l}%
				we
					\marginnote{l}%
				now say to Zion \uline{this} \uline{once} in the name of the Lord, Repent! Repent! Awake! Awake! put on thy beautiful garments before you are made to feel the chastning rod of him whose anger is kind[l]ed against you, let not satan tempt you to think we want to make you bow to us to domaneur over you for God knows this is not the case our eyes are watered with tears and our hearts are poured out to God in prayer for you that he will spare you and turn away his anger from you,
					\hypertarget{EMS-1833-JS-14-JAN-1833-m}%
				There
					\marginnote{m}
				\phantom{ }\sout{is} are many things in the last letters from Brs G and P. that are good and we esteem them much, The Idea of having certain ones appointed to regulate Zion and travling Elders have nothing to do with this part of the matter is something that we highly approbate, and you will doubtless know before this reaches you why Wm McLelin [William E. McLellin] opposed you in this move---
					\hypertarget{EMS-1833-JS-14-JAN-1833-n}%
				We
					\marginnote{n}%
				fear there was something in Bro Gilberts when he returned to this place from N York last fall in relation to his Bro William that was not right, For Bro G was asked 2 or 3 times about his Bro W, but gave evasive answers, and at the same time he knew that William was in Cleaveland, but the Lord has taken him we merely mention this that all may take warning to work in the light for God will bring evry secret thing to Light,
					\hypertarget{EMS-1833-JS-14-JAN-1833-o}%
				we
					\marginnote{o}%
				now close our Epistle by saying unto you the Lord has commanded us to purify ourselves, to wash our hands and our feet that he might testify to his father and our Father to his God \& our God that we are clean from the blood of this Generations and before we could wash our our hands and our feet we were constrained to write this letter, Therefore with the feelings of inexpressable anxiety for your welfare we say again Repent Rep[en]t or Zion must suffer, for the scourge and judgments must come upon her, Let the Bishop read this to the Elders that they may warn the members of the scourge that is coming exept they repent.
					\hypertarget{EMS-1833-JS-14-JAN-1833-p}%
				Tell
					\marginnote{p}%
				them to read the Book of \sout{Commandments} Mormon and obey it, read the commandments that are printed and obey them. yea humble yourselves under the mighty hand of God, that peradventure he may turn away his anger from you, tell them that they have not come up to Zion to sit down in idleness neglecting the things of God, but they are to be dilligent and faithful in obeying the new covenant,
					\hypertarget{EMS-1833-JS-14-JAN-1833-q}%
				There
					\marginnote{q}%
				is one clause in Br Joseph Letter which you may not understand. that is this if the people of Zion did not repent the Lord would seek another place, and another people, Zion is the place where the Temple will be built, and the people gathered but all people upon that holy Land being under condemnation, the lord will cut off if they repent not and bring another race upon it that will serve him, The Lord will seek another place to bring forth and prepare his word to go forth to the nations, and as we said before so we say again Bro Joseph will not settle in Zion except she repent and serve God and obey the new covenant with this explanation the conference sanctions Bro Joseph letter---
					\hypertarget{EMS-1833-JS-14-JAN-1833-r}%
				Brethren
					\marginnote{r}%
				the conference meets again this evening to hear this Letter read and if it meets their minds we have all agreed to kneel down before the Lord and cry unto him with all our hearts that this epistle \& Bro J\supers{s.} and the revelations also may have there desired effect and accomplish the things whereunto they are sent and that they may stimulate you to cleans Zion that she mourn not, Therefore when you get this know ye that a conference of 12 High Priests have cried unto the Lord for you and are still crying saying spare thy peopl[e] O! Lord and give not thy heritage to reproach,
					\hypertarget{EMS-1833-JS-14-JAN-1833-s}%
				We
					\marginnote{s}%
				now feel that our garments are clean from you, and all men, when we have washed our \sout{hand \& our} feet \& hands according to the commandment We have written plain at this time but we believe not harsh, \sout{Plainly} Plainness is what the Lord requires and we should not feel ourselves clean unless we had done so, and if the things we have told you be not attended to you will not long have occasion to say or to think rather that we may be rong in what we have stated Your unworthy brethren are determined to pray unto the Lord for Zion \sout{as} so long as we can shed \sout{a tear} the sympathising tear, or feel any spirit \sout{of} to supplicate a throne of Grace in her behalf,
					\hypertarget{EMS-1833-JS-14-JAN-1833-t}%
				The
					\marginnote{t}%
				School of the Prophets will commence if the Lord will in 2 or 3 days, a general time of health with us the cause of God seems to be rapidly advancing in the eastern \sout{country} countries, the gifts begining to break forth so as to astonish the world, and even believers marvel at the power and goodness of God, Thanks be rendered to his holy name for what he is doing, We are your unworthy brethren in the Lord \& may the Lord help us all to do his will that we may at last be saved in his kingdom------
				
				\begin{tightright}
					Orson Hyde
					
					Hyram Smith
				\end{tightright}
				
					\hypertarget{EMS-1833-JS-14-JAN-1833-s}%
				N. B.
					\marginnote{s}%
				we stated that Bro Gilbert knew that William was in Cleaveland last fall, when he was in Kirtland, we wrote this upon the strength of hearsay, but Wm being left at St Louis strengthened our suppositions that such was the fact, we stated farther respecting this matter \uline{or} \uline{this} \uline{itim} than the testamony will warrant us--- with this exception the Conference sanctions this letter
				
		% -----
		\subsubsection{Minutes, 22--23 January 1833}
		% -----
			\label{EMS-1833-JS-22-23-JAN-1833}
			% \label{EMS-1830-JS-DAY-3LETTERMONTH-YEAR}
		
			\begin{description}
				\item[Print Sources]
			\end{description}

				\begin{tabular}{p{0.5in}p{1in}p{0.5in}p{1in}}
				JSP	 	& D2:378--382		& JSR 		& --- \\
				DC	 	& --- 				& HC 		& 1:322--323 \\
				HDDC 	& ---				& TCHC 		& 1:235--236 \\
				\end{tabular}
				% HDDC = woodford's DC diss; JSR = Marquardt's JS Revelations; TCHC = Vogel HC
			
			\begin{description}
				\item[Online Access] \href{https://www.josephsmithpapers.org/paper-summary/minutes-22-23-january-1833/1}{Here}
				% \item[Analysis] \hyperref[XX]{Here}
			\end{description}
			
			\begin{description}
				\item[Background]
			\end{description}
			
				% It might also be useful to give a two sentence context of the period: e.g., "This speech was given in Nauvoo to a small congregation one month prior to the destruction of the Nauvoo Expositor. These and other tensions may have been on the prophet's mind when he gave the speech."
			
			\begin{description}
				\item[Original Source]
			\end{description}
			
				% brief description of original source material, with reference to JSP or somewhere else for more info
				% Background material: it might be helpful to have a note that says whether a given document was presented as a speech, was written in a journal, etc.
				% Also a comment here about how reliable the source might be, or what shape it is in, if there are large omissions or parts that are hard to understand, and so on.
			
			\begin{description}
				\item[Text]
			\end{description}
			
				\begin{tightright}
						\hypertarget{EMS-1833-JS-22-23-JAN-1833-a}%
					Kirtland 
						\marginnote{a}
					\sout{January 2$\diamond$}\supers{\sout{th}} 22\supers{d} 1833
				\end{tightright}
				
				conference of high pr[i]ests convened \sout{at} <in> the coucel [council] room high pr[i]ests p[r]esent Joseph Smith J\supersu{un} President Sidney Rigdon cheif scribe and high counceler Frederick G. Williams assistant scribe and counceler Newel K. Whitney Bishop Hiram [Hyrum] Smith Bishops counceler Zebedee Coltrin Joseph Smith \supers{Sr.} Samuel H Smith John Murdock Lyman Johnson Orson Hyde Ezra Thair [Thayer]; Elders, Levy [Levi] Hancock William Smith------
				
					\hypertarget{EMS-1833-JS-22-23-JAN-1833-b}%
				Conference
					\marginnote{b}%
				opened with prayer by the President, after prayer the President spake in an unknown Tongue he was followed by Br Zebede Coltrin and he by Bro William Smith after this the gift was poured out in a miraculous manner until all the Elders obtained the gift together with several of the members of the Church both male \& female Great and glorious were the divine manifestation of the Holy Spirit, Praises were sang to God \& the Lamb besides much speaking \& praying all in tongues. The conference adjourned at a late hour in the night to meet next morning at 9 oclock closed \sout{by} with prayer by the President
				
					\hypertarget{EMS-1833-JS-22-23-JAN-1833-c}%
				Wednesday
					\marginnote{c}%
				Jan\supersu{ry} 23\supers{d} Meet agreeable to adjournment. Conference opened with Prayer by the President and after much speaking praying and singing, all done in Tongues proceded to washing hands faces \& feet in the name of the Lord as commanded \sout{by} of God each one washing his own after which the president guirded himself with a towel and again washed the feet of all the Elders wiping them with the towel, his father presenting himself the President asked of him a blessing before he would wash his feet which he obtained by the laying on of his fathers hands, pronouncing upon his head that he should continue in his Priests office untill Christ come------
				
					\hypertarget{EMS-1833-JS-22-23-JAN-1833-d}%
				at
					\marginnote{d}%
				the close of which scene Br F G Williams being moved upon by the Holy Ghost washed the feet of the President as a token of his fixed determination to be with him in suff[er]ing or in rejoicing, in life or in death and to be continually on his right hand in which thing he was accepted, The President said after he had washed the feet of the Elders, as I have done so do ye wash ye therefee [therefore] one anothers feet pronouncing at the same time through the power of the Holy Ghost that the Elders were all clean from the blood of this generation
					\hypertarget{EMS-1833-JS-22-23-JAN-1833-e}%
				but
					\marginnote{e}%
				that those \sout{who} among them who should sin wilfully after they were thus cleansed and sealed up unto eternal life should be given over unto the buffettings of Satan until the day of redemption Having continued all day in fasting \& prayer before the Lord at the close they partook of the Lords supper which was blessed by the president in the name of the Lord all eat and drank and were filled then sang an hymn and went out---
				
				\begin{tightright}
				\uline{F. G. Williams Clk---}
				\end{tightright}
				
		% -----
		\subsubsection{Letter to the Church in Thompson, Ohio, 6 February 1833}
		% -----
			\label{EMS-1833-JS-6-FEB-1833}
			% \label{EMS-1830-JS-DAY-3LETTERMONTH-YEAR}
		
			\begin{description}
				\item[Print Sources]
			\end{description}

				\begin{tabular}{p{0.5in}p{1in}p{0.5in}p{1in}}
				JSP	 	& D3:3--6		& JSR 		& --- \\
				DC	 	& --- 			& HC 		& 1:324--325 \\
				HDDC 	& ---			& TCHC 		& 1:236--237 \\
				TPJS 	& 19--20		& 	 		&  \\
				\end{tabular}
				% HDDC = woodford's DC diss; JSR = Marquardt's JS Revelations; TCHC = Vogel HC
			
			\begin{description}
				\item[Online Access] \href{https://www.josephsmithpapers.org/paper-summary/letter-to-the-church-in-thompson-ohio-6-february-1833/1}{Here}
				% \item[Analysis] \hyperref[XX]{Here}
			\end{description}
			
			\begin{description}
				\item[Background]
			\end{description}
			
				% It might also be useful to give a two sentence context of the period: e.g., "This speech was given in Nauvoo to a small congregation one month prior to the destruction of the Nauvoo Expositor. These and other tensions may have been on the prophet's mind when he gave the speech."
			
			\begin{description}
				\item[Original Source]
			\end{description}
			
				% brief description of original source material, with reference to JSP or somewhere else for more info
				% Background material: it might be helpful to have a note that says whether a given document was presented as a speech, was written in a journal, etc.
				% Also a comment here about how reliable the source might be, or what shape it is in, if there are large omissions or parts that are hard to understand, and so on.
			
			\begin{description}
				\item[Text]
			\end{description}
			
				\begin{tightright}
						\hypertarget{EMS-1833-JS-6-FEB-1833-a}%
					An Epistle to the Church in Thompson
						\marginnote{a}
					
					Kirtland February 6th 1833------
					
					To the Church of Christ in Thompson Geauga County Ohio
				\end{tightright}
				
				Dear Brethren
				
				We salute you by this our Epistle in the bonds of Love rejoicing in your steadfastness in the faith which is in Christ Jesus our Lord and we desire your prosperity in the ways of truth and righteousness in the bowels of Jesus Christ praying for you continually that your faith fail not and that you may overcome all the evils with which you are surrounded and become pure and holy before God even our father to whom be glory for \sout{and} ever and ever Amen
				
					\hypertarget{EMS-1833-JS-6-FEB-1833-b}%
				It has seemed good unto the holy spirit
					\marginnote{b}%
				and unto us to send this our epistle to you by the hand of our beloved broth[e]r Salmon [Gee] your messinger who has been ordained by us in obedience to the commandment of God to the office of Elder to preside over the Church in Thompson taking the oversight thereof to Lead you \& to teach the thing[s] which are according to Godliness in whom we have great confidence as we presume also you we therefore say to you yea not us only but the Lord also receive you him as such knowing that the lord has appointed him to this office for your good holding him up by your prayers praying for him continually that he may be endued with wisdom and understanding in the knowledge of the Lord that through him you may be kept from evil spirits and all strifes and discensions and grow in grace and in the know[l]edge of our Lord and Saviour Jesus Christ
				
					\hypertarget{EMS-1833-JS-6-FEB-1833-c}%
				Brethren beloved continue in brotherly
					\marginnote{c}%
				Love walk\sout{ing} in meekness watching unto prayer that you be not overcome follow after peace as said our beloved brothe[r] Paul that you may be the children of our heavenly Father and <not> give occasion <of> \sout{for} stumbling to saint or sinner--- finely [finally] brethren pray for us that we may be enabled to do the work whereunto we are called that you may enjoy the mysteries of God even a fulness and may the grace of our Lord Jesus Christ be with you all amen.
				
				\begin{tightright}
					Joseph Smith Jr---
					
					Sidney Rigdon
					
					F. G. William [Frederick G. Williams]------
				\end{tightright}
				
		% -----
		\subsubsection{Letter to Noah C. Saxton, 12 February 1833}
		% -----
			\label{EMS-1833-JS-12-FEB-1833}
			% \label{EMS-1830-JS-DAY-3LETTERMONTH-YEAR}
		
			\begin{description}
				\item[Print Sources]
			\end{description}

				\begin{tabular}{p{0.5in}p{1in}p{0.5in}p{1in}}
				JSP	 	& D3:6--10		& JSR 		& --- \\
				DC	 	& --- 			& HC 		& 1:326 \\
				HDDC 	& ---			& TCHC 		& 1:238--239 \\
				\end{tabular}
				% HDDC = woodford's DC diss; JSR = Marquardt's JS Revelations; TCHC = Vogel HC
			
			\begin{description}
				\item[Online Access] \href{https://www.josephsmithpapers.org/paper-summary/letter-to-noah-c-saxton-12-february-1833/1}{Here}
				% \item[Analysis] \hyperref[XX]{Here}
			\end{description}
			
			\begin{description}
				\item[Background]
			\end{description}
			
				% It might also be useful to give a two sentence context of the period: e.g., "This speech was given in Nauvoo to a small congregation one month prior to the destruction of the Nauvoo Expositor. These and other tensions may have been on the prophet's mind when he gave the speech."
			
			\begin{description}
				\item[Original Source]
			\end{description}
			
				% brief description of original source material, with reference to JSP or somewhere else for more info
				% Background material: it might be helpful to have a note that says whether a given document was presented as a speech, was written in a journal, etc.
				% Also a comment here about how reliable the source might be, or what shape it is in, if there are large omissions or parts that are hard to understand, and so on.
			
			\begin{description}
				\item[Text]
			\end{description}
			
					\hypertarget{EMS-1833-JS-12-FEB-1833-a}%
				To
					\marginnote{a}%
				N C Sexton [Noah C. Saxton] Rochester
				
				Dear sir
				
				I was somewhat disappointed on receiving my paper with only a part of my letter inserted in \sout{that} it. The letter which I wrote you for publication I wrote by the commandment of God, and I am quite anxious to have it all laid before the public for it is of \uline{importance} to them, But I have no clame upon you, neither do I wish to urge you beyond that which is reasonable to do it. I have only to appeal to your extended geneorsity to all religious societies that claim that Christ has come in the flesh and also tell you what will be the consequen[c]es of a neglect to publish it--- some parts of the letter were very severe upon the wickedness of sectarianism---
					\hypertarget{EMS-1833-JS-12-FEB-1833-b}%
				I
					\marginnote{b}%
				acknowledge and the truth, remember is hard and severe against all iniquity and wickedness, but this is no reason why it should not be published but the very reason why it should, It lays the axe at the root of the tree and I long to see many of the sturdy oaks which \sout{I} have long cumbred the ground fall prostrate. I now say unto you that if you wish to clear your garments from the blood of you[r] readers I exhort you to publish that letter \uline{entire} but if not the sin be upon your head---
				
				\begin{tightright}
				Accept sir the good wishes and tender regard of your unworthy servant---
				
				Joseph Smith Jr
				\end{tightright}
				
				Kirtland 12\supers{th} Feby. 1833
				
		% -----
		\subsubsection{Revelation, 27 February 1833 [D\&C 89]}
		% -----
			\label{EMS-1833-JS-27-FEB-1833}
			% \label{EMS-1830-JS-DAY-3LETTERMONTH-YEAR}
		
			\begin{description}
				\item[Print Sources]
			\end{description}

				\begin{tabular}{p{0.5in}p{1in}p{0.5in}p{1in}}
				JSP	 	& D3:11--24			& JSR 		& 232--233 \\
				DC	 	& Section 89 		& HC 		& 1:327--329 \\
				HDDC 	& 1169--1181		& TCHC 		& 1:240 \\
				\end{tabular}
				% HDDC = woodford's DC diss; JSR = Marquardt's JS Revelations; TCHC = Vogel HC
			
			\begin{description}
				\item[Online Access] \href{https://www.josephsmithpapers.org/paper-summary/revelation-27-february-1833-dc-89/1}{Here}
				% \item[Analysis] \hyperref[DC89]{Here}
			\end{description}
			
			\begin{description}
				\item[Background]
			\end{description}
			
				% It might also be useful to give a two sentence context of the period: e.g., "This speech was given in Nauvoo to a small congregation one month prior to the destruction of the Nauvoo Expositor. These and other tensions may have been on the prophet's mind when he gave the speech."
			
			\begin{description}
				\item[Original Source]
			\end{description}
			
				% brief description of original source material, with reference to JSP or somewhere else for more info
				% Background material: it might be helpful to have a note that says whether a given document was presented as a speech, was written in a journal, etc.
				% Also a comment here about how reliable the source might be, or what shape it is in, if there are large omissions or parts that are hard to understand, and so on.
			
			\begin{description}
				\item[Text]
			\end{description}
			
				\begin{tightcenter}
					A \uline{Word of Wisdom}
				\end{tightcenter}
				
				A word of wisdom for the benefit of the Saints in these last days and also the Saints in Zion to be sent greeting, not by commandment or Constraint, but by Revelation \& the word of wisdom shewing forth the order \& will of God in the temporal salvation of all Saints, given for a principle with promise, adapted to the Capacity of the weak \& the weakest of all Saints who are or can be called Saints---
				
				Behold verily thus Saith the Lord unto you in consequence of evils \& designs which will exist in the hearts of conspiring men in these last days, I have warned you \& forewarned you by giving unto you this word of wisdom by Revelation, that inasmuch as any man drinketh wine or Strong drink among you behold it is not good, neither mete in the sight of your Father, only in assembling your[s]elves in your Sacraments before him, \& behold this should be wine of your own make \& again Strong drinks are not for the belly, but for the washing of your bodies, \& Tobacco is not for \uline{man} but is for bruises \& all sick cattle to be used with judgement \& skill.
				
				And again hot drinks are not for the body or belly, all wholesome herbs God hath ordained for the constitution \& nature \& use of man, every herb in the season thereof \& every fruit in the season thereof, all these to be used with prudence \& thanksgiving, yea flesh also of beasts \& of fowls, I the Lord hath ordained for the use of man with thanksgiving--- Nevertheless they are to be used sparingly \& it is pleasing unto me that they should not be used only in times of winter or of famine--- All grain is for the use of man \& of beasts to be the staff of life not only for man, but for the beasts \& for the fowls, and all wild animals that run or creep on the earth \& these hath God [made] for the use of man only in times of famine or excess of hunger all grain is good for the use of man \sout{\& of beasts} as also the fruit of the vine that which beareth fruit whether in the ground or above ground. Nevertheless \uline{wheat} for \uline{man} \& corn for the Ox \& Oats for the horse. Rye for the fowls \& the swine \& for all beasts of the field and Barley for all useful animals \& for mild drinks as also other grains--- and all saints who remember to keep \& do these sayings walking in obedience to the commands shall receieve health in their navel \& marrow to their bones \& shall find wisdom \& great treasures of wisdom \& knowledge even hidden treasures \& shall run \& not be weary \& walk \& not faint. And I the Lord give unto them a promise that the destroying Angel shall pass them by as they did by the Children of Israel \& not slay them
				
		% -----
		\subsubsection{Revelation, 8 March 1833 [D\&C 90]}
		% -----
			\label{EMS-1833-JS-8-MAR-1833}
			% \label{EMS-1830-JS-DAY-3LETTERMONTH-YEAR}
		
			\begin{description}
				\item[Print Sources]
			\end{description}

				\begin{tabular}{p{0.5in}p{1in}p{0.5in}p{1in}}
				JSP	 	& D3:24--32			& JSR 		& 234--236 \\
				DC	 	& Section 90 		& HC 		& 1:329--331 \\
				HDDC 	& 1182--1194		& TCHC 		& 1:241--242 \\
				\end{tabular}
				% HDDC = woodford's DC diss; JSR = Marquardt's JS Revelations; TCHC = Vogel HC
			
			\begin{description}
				\item[Online Access] \href{https://www.josephsmithpapers.org/paper-summary/revelation-8-march-1833-dc-90/1}{Here}
				% \item[Analysis] \hyperref[DC90]{Here}
			\end{description}
			
			\begin{description}
				\item[Background]
			\end{description}
			
				% It might also be useful to give a two sentence context of the period: e.g., "This speech was given in Nauvoo to a small congregation one month prior to the destruction of the Nauvoo Expositor. These and other tensions may have been on the prophet's mind when he gave the speech."
			
			\begin{description}
				\item[Original Source]
			\end{description}
			
				% brief description of original source material, with reference to JSP or somewhere else for more info
				% Background material: it might be helpful to have a note that says whether a given document was presented as a speech, was written in a journal, etc.
				% Also a comment here about how reliable the source might be, or what shape it is in, if there are large omissions or parts that are hard to understand, and so on.
			
			\begin{description}
				\item[Text]
			\end{description}
			
				\begin{tightcenter}
				Kirtland 8\supers{th} of March 1833
				\end{tightcenter}
				
				A Commandment given unto Joseph saying thus saith the Lord verily verily I say unto you my son thy sins are forgiven thee according to thy petition for thy prayers and the prayers of thy brethren have come up into my ears therefore thou art blessed from henceforth that bear the keys of the kingdom given unto you which kingdom is coming forth for the Last time verily I say unto you the keys of this kingdom shall never be taken from you whilst thou art in the world neither in the world to come never[the]less through you shall the oricles be given unto another yea even unto the church and all they who receive the oricles of God let them be aware how they hold them lest they are accounted as a light thing and are brought under condemnation thereby and stumble and fall when the storms descend <\&> the winds blow and the rains descend and beat upon their house and again verily I say unto thy brethren Sidney [Rigdon] and Fredrick [Frederick G. Williams] there sins are forgiven them also and they are accounted as equal with thee in holding the keys of this Last Kingdom as also through your administration the <keys of the> School of the prophets which I have commanded to be organized that thereby they may be perfected in their minstry for the salvation of Zion and of the Nations of Israel and of the Gentiles as many as will believe that through your administration they may receive the word and through their administration the word may go forth unto the ends of the earth unto the Gentiles first and then behold and Lo they shall turn unto the Jews and then cometh the day when the arm of the Lord shall be reveiled in power in convincing the nations the heathen nations the house of Joseph of the Gospel of their salvation for it shall come to pass in that day that evry man shall hear the fulness of the Gospel in his own Tongue and in his own Language through thou who are ordained unto this power by the administration of the comforter shed forth upon them for the revelation of Jesus Christ and now verely I say unto you I give unto you a commandment that you continue in this ministry and presidency and when you have finished the translation of the prophets you shall from \sout{them} <thence> forth preside over the affairs of the Church and the School and from time to time as shall be manifest by the comfo[r]ter receive revelations to unfold the mystres [mysteries] of the kingdom and set in order the churches and study and Learn and become acquainted with all good books and with Languages tongues and people \&c \&c and this shall be your business and mission in all your Lives to preside in council and set in order all the affairs of this Church and kingdom be not ashamed neither confounded but be admonished in all your high mindedness and pride for it bringeth a snare upon your souls set in order your houseses keep slothfulness and uncleanliness far from you now verily I say unto you let there be a place provided as soon as it is possable \sout{for} \sout{you} for the family of thy councellor \& scribe even Frederick and <let> mine Aged servant Joseph [Smith Sr.] continue with his family upon the place <where he now lives> and let it not be sold untill the mouth of the Lord shall name and let thy councellor even Sidney remain where he now resides untill the mouth of the Lord shall name <and> let the Bishop search dilligently to obtain an agent and let it be a man who has got riches in store a man of God and of strong faith that thereby he may be enabled to discharge evry debt that the store house of the Lord may not be brought in to disrepute before the eyes of the people search diligently pray always and be believing and all things shall work together for your good if ye walk uprightly and remember the covenant where with ye have covenanted one with another let your families be small especially \sout{my} mine aged Servant Joseph as pertaining to thou who do not belong to your families that those things that are provided for you to bring to pass my work are not taken from you and given to those that are not worthy and thereby you are hindred in accomplishing <those things which> I have commanded you and again verely I say unto you it is my will that my hand maid\sout{en} Viana [Vienna Jaques] should receive money to bear her expences and go up unto the Land of Zion and the residue of her money I will consecrate unto myself and reward her in mine own due time verely I say unto you <that> it is meet in mine eyes that she should go up unto the Land of Zion and receive an inheritance from the hand of the Bishop that she may settle down in peace <in as much as she is faithful> and not be Idle in her days from thenceforth and behold verely I say unto you that ye shall write this commandment and say unto your brethren in Zion in Love greeting that I have called you also to preside over Zion in mine own due time therefore let them cease wear[y]ing me concerning this matter behold I say unto you that your brethren in Zion begin to repent and the Angels rejoice over them nevertheless I am not well pleased with many things and I am not well please[d] with my servant William E McLel[l]in neithe[r] with my servant sidney Gilbert and the Bishop also and others have many things to repent of but verely I say unto you that I the Lord will contend with Zion and plead with her strong ones and chasten her untill she overcome and are clean before me for she shall not be moved out of her place I the Lord have spoken it--- Am[en]---
				
		% -----
		\subsubsection{Revelation, 9 March 1833 [D\&C 91]}
		% -----
			\label{EMS-1833-JS-9-MAR-1833}
			% \label{EMS-1830-JS-DAY-3LETTERMONTH-YEAR}
		
			\begin{description}
				\item[Print Sources]
			\end{description}

				\begin{tabular}{p{0.5in}p{1in}p{0.5in}p{1in}}
				JSP	 	& D3:32--35			& JSR 		& 236 \\
				DC	 	& Section 91 		& HC 		& 1:331--332 \\
				HDDC 	& 1195--1200		& TCHC 		& 1:242 \\
				\end{tabular}
				% HDDC = woodford's DC diss; JSR = Marquardt's JS Revelations; TCHC = Vogel HC
			
			\begin{description}
				\item[Online Access] \href{https://www.josephsmithpapers.org/paper-summary/revelation-9-march-1833-dc-91/1}{Here}
				% \item[Analysis] \hyperref[DC91]{Here}
			\end{description}
			
			\begin{description}
				\item[Background]
			\end{description}
			
				% It might also be useful to give a two sentence context of the period: e.g., "This speech was given in Nauvoo to a small congregation one month prior to the destruction of the Nauvoo Expositor. These and other tensions may have been on the prophet's mind when he gave the speech."
			
			\begin{description}
				\item[Original Source]
			\end{description}
			
				% brief description of original source material, with reference to JSP or somewhere else for more info
				% Background material: it might be helpful to have a note that says whether a given document was presented as a speech, was written in a journal, etc.
				% Also a comment here about how reliable the source might be, or what shape it is in, if there are large omissions or parts that are hard to understand, and so on.
			
			\begin{description}
				\item[Text]
			\end{description}
			
				\begin{tightcenter}
				Kirtland 9th of March 1833
				\end{tightcenter}
				
				A Revelation given concerning Apocrypha
				
				Verily thus saith the Lord unto you concerning the Apocrypha there are many things contained therein that are true and it is mostly translated correct--- there are many things contained therein that are not true which are interpelations by the hands of men varely I say unto you that it is not needful that the Apocrypha should be translated therefore whoso readeth it let him understand for the spirit manifesteth truth and and whoso is enlightened by the spirit shall obtain benifit therefrom and whoso receiveth not the spirit cannot be benefited; Therefore it is not needful that it should be translated. Amen
				
		% -----
		\subsubsection{Revelation, 15 March 1833 [D\&C 92]}
		% -----
			\label{EMS-1833-JS-15-MAR-1833}
			% \label{EMS-1830-JS-DAY-3LETTERMONTH-YEAR}
		
			\begin{description}
				\item[Print Sources]
			\end{description}

				\begin{tabular}{p{0.5in}p{1in}p{0.5in}p{1in}}
				JSP	 	& D3:35--37			& JSR 		& 236--237 \\
				DC	 	& Section 92 		& HC 		& 1:333 \\
				HDDC 	& 1201--1206		& TCHC 		& 1:243 \\
				\end{tabular}
				% HDDC = woodford's DC diss; JSR = Marquardt's JS Revelations; TCHC = Vogel HC
			
			\begin{description}
				\item[Online Access] \href{https://www.josephsmithpapers.org/paper-summary/revelation-15-march-1833-dc-92/1}{Here}
				% \item[Analysis] \hyperref[DC92]{Here}
			\end{description}
			
			\begin{description}
				\item[Background]
			\end{description}
			
				% It might also be useful to give a two sentence context of the period: e.g., "This speech was given in Nauvoo to a small congregation one month prior to the destruction of the Nauvoo Expositor. These and other tensions may have been on the prophet's mind when he gave the speech."
			
			\begin{description}
				\item[Original Source]
			\end{description}
			
				% brief description of original source material, with reference to JSP or somewhere else for more info
				% Background material: it might be helpful to have a note that says whether a given document was presented as a speech, was written in a journal, etc.
				% Also a comment here about how reliable the source might be, or what shape it is in, if there are large omissions or parts that are hard to understand, and so on.
			
			\begin{description}
				\item[Text]
			\end{description}
			
				\begin{tightcenter}
				Kirtland 15th March 1833---
				\end{tightcenter}
				
				Verely thus saith the Lord I give unto the united firm organized agreeable to the commandment previously given a revelation \& commandment concerning my servant Frederick [G. Williams] that ye shall r[e]ceive him into the firm what I say unto one I say unto all. and again I say unto you my servent Frederick thou shalt be a lively member in this firm and inas much as thou art faithful in keeping all former commandments thou shalt be blessed for ever Amen------
				
		% -----
		\subsubsection{Minutes, 18 March 1833}
		% -----
			\label{EMS-1833-JS-18-MAR-1833}
			% \label{EMS-1830-JS-DAY-3LETTERMONTH-YEAR}
		
			\begin{description}
				\item[Print Sources]
			\end{description}

				\begin{tabular}{p{0.5in}p{1in}p{0.5in}p{1in}}
				JSP	 	& D3:39--43		& JSR 		& --- \\
				DC	 	& --- 			& HC 		& 1:334--335 \\
				HDDC 	& ---			& TCHC 		& 1:245--246 \\
				\end{tabular}
				% HDDC = woodford's DC diss; JSR = Marquardt's JS Revelations; TCHC = Vogel HC
			
			\begin{description}
				\item[Online Access] \href{https://www.josephsmithpapers.org/paper-summary/minutes-18-march-1833/1}{Here}
				% \item[Analysis] \hyperref[XX]{Here}
			\end{description}
			
			\begin{description}
				\item[Background]
			\end{description}
			
				% It might also be useful to give a two sentence context of the period: e.g., "This speech was given in Nauvoo to a small congregation one month prior to the destruction of the Nauvoo Expositor. These and other tensions may have been on the prophet's mind when he gave the speech."
			
			\begin{description}
				\item[Original Source]
			\end{description}
			
				% brief description of original source material, with reference to JSP or somewhere else for more info
				% Background material: it might be helpful to have a note that says whether a given document was presented as a speech, was written in a journal, etc.
				% Also a comment here about how reliable the source might be, or what shape it is in, if there are large omissions or parts that are hard to understand, and so on.
			
			\begin{description}
				\item[Text]
			\end{description}
			
					\hypertarget{EMS-1833-JS-18-MAR-1833-a}%
				Kirtland
					\marginnote{a}
				\phantom{ }\sout{1833} 18 March 1833---
				
				This day an assembly of the high Priests meet at the school room of the prophets and were organized in due form by solemn prayer to the most high by sidney Rigdon then proceded to ordain Doctor [Philastus] Hurlbut to be an elder under the hand of Sidney Rigdon after which Bro Sidney arose and desired that he and Bro Frederick [G. Williams] Should be ordained to the office that they had been called Viz to the of President of the high Priesthood and to be equal in holding the keys of the Kingdom with Broth[e]r Joseph\sout{s} Smith J---
					\hypertarget{EMS-1833-JS-18-MAR-1833-b}%
				according
					\marginnote{b}%
				to a revelation given on the [8]th day of March 1833 in Kirtland sayi[ng] thus, and again verely I say unto \sout{thy} <my> brethren Sidney and Frederick their sins are forgiven them also and they are accounted equal in holding the keys of this <last> Kingdom. and again I give unto you a commandment that you continue in this ministry and presidency and when you have finished the translating of the prophets you shall from thenceforth preside over the affairs of the Church and the school from time to time as shall be manifest by the comfo[r]ter receive revelation to unfold the mysteries of the Kingdom and set in order the church.
					\hypertarget{EMS-1833-JS-18-MAR-1833-c}%
				Accordingly
					\marginnote{c}%
				Bro Joseph proceded to and ordained them by the Laying on of the hands to be equal with him in holding the Keys of the Kingdom and also to the Presidency of the high Priest hood. After which several exertations were given to faithfulness and obedience to the commandments of God and much useful instruction given for the benefit of the saints with a promise that the pure in heart that were present should see a heavenly vision,
					\hypertarget{EMS-1833-JS-18-MAR-1833-d}%
				and
					\marginnote{d}%
				after remaining for a short time in secret prayer the promise was verified to many present having the eyes of their understanding<s> opened so as to behold many things afte[r] which the bread and wine was distributed by Bro Joseph after which many of the brethren saw a heavenly vision of the saviour and concourses of angels and many othe[r] thing[s] of which each one has a reccord of what they saw \&c------
				
				\begin{tightright}
				<F G Williams Ck PT. [pro tempore]>
				\end{tightright}
				
		% -----
		\subsubsection{License for Frederick G. Williams, 20 March 1833}
		% -----
			\label{EMS-1833-JS-20-MAR-1833}
			% \label{EMS-1830-JS-DAY-3LETTERMONTH-YEAR}
		
			\begin{description}
				\item[Print Sources]
			\end{description}

				\begin{tabular}{p{0.5in}p{1in}p{0.5in}p{1in}}
				JSP	 	& D3:43--46		& JSR 		& --- \\
				DC	 	& --- 			& HC 		& --- \\
				HDDC 	& ---			& TCHC 		& --- \\
				\end{tabular}
				% HDDC = woodford's DC diss; JSR = Marquardt's JS Revelations; TCHC = Vogel HC
			
			\begin{description}
				\item[Online Access] \href{https://www.josephsmithpapers.org/paper-summary/license-for-frederick-g-williams-20-march-1833/1}{Here}
				% \item[Analysis] \hyperref[XX]{Here}
			\end{description}
			
			\begin{description}
				\item[Background]
			\end{description}
			
				% It might also be useful to give a two sentence context of the period: e.g., "This speech was given in Nauvoo to a small congregation one month prior to the destruction of the Nauvoo Expositor. These and other tensions may have been on the prophet's mind when he gave the speech."
			
			\begin{description}
				\item[Original Source]
			\end{description}
			
				% brief description of original source material, with reference to JSP or somewhere else for more info
				% Background material: it might be helpful to have a note that says whether a given document was presented as a speech, was written in a journal, etc.
				% Also a comment here about how reliable the source might be, or what shape it is in, if there are large omissions or parts that are hard to understand, and so on.
			
			\begin{description}
				\item[Text]
			\end{description}
			
				A Conference of High Priests of the Church of Jesus Christ certify that Frederick G. Williams the bearer of this after due examination of his moral character and Christian attainments was found worthy to receive their testamonials from under our hands. We therefore Certify that he has been regularly ordained to the Presidency of the High priesthood under the hand of Joseph Smith Jr after the holy order of God according to a commandment given the 8th day of March AD 1833 to preside over the affairs of the Church and Kingdom of Jesus Christ as established in these last days.
				
				Given under our hands at Kirtland Geauga Co Ohio North America the 20\supers{th} day of March 1833
				
				\begin{tightright}
				Joseph Smith Jr Prd
				
				\uline{Sidney} \uline{Rigdon}
				\end{tightright}
				
		% -----
		\subsubsection{Letter to John S. Carter, 13 April 1833}
		% -----
			\label{EMS-1833-JS-13-APR-1833}
			% \label{EMS-1830-JS-DAY-3LETTERMONTH-YEAR}
		
			\begin{description}
				\item[Print Sources]
			\end{description}

				\begin{tabular}{p{0.5in}p{1in}p{0.5in}p{1in}}
				JSP	 	& D3:56--64		& JSR 		& --- \\
				DC	 	& --- 			& HC 		& 1:338--340 \\
				HDDC 	& ---			& TCHC 		& 1:248--251 \\
				TPJS 	& 20--22		& 	 		&  \\
				\end{tabular}
				% HDDC = woodford's DC diss; JSR = Marquardt's JS Revelations; TCHC = Vogel HC
			
			\begin{description}
				\item[Online Access] \href{https://www.josephsmithpapers.org/paper-summary/letter-to-john-s-carter-13-april-1833/1}{Here}
				% \item[Analysis] \hyperref[XX]{Here}
			\end{description}
			
			\begin{description}
				\item[Background]
			\end{description}
			
				% It might also be useful to give a two sentence context of the period: e.g., "This speech was given in Nauvoo to a small congregation one month prior to the destruction of the Nauvoo Expositor. These and other tensions may have been on the prophet's mind when he gave the speech."
			
			\begin{description}
				\item[Original Source]
			\end{description}
			
				% brief description of original source material, with reference to JSP or somewhere else for more info
				% Background material: it might be helpful to have a note that says whether a given document was presented as a speech, was written in a journal, etc.
				% Also a comment here about how reliable the source might be, or what shape it is in, if there are large omissions or parts that are hard to understand, and so on.
			
			\begin{description}
				\item[Text]
			\end{description}
			
				\begin{tightcenter}
						\hypertarget{EMS-1833-JS-13-APR-1833-a}%
					Kirtlan[d]
						\marginnote{a}
					13th April 1833
				\end{tightcenter}
				
				Dear Broth[er] [John S.] Carter your Letter to Broth Jared [Carter] is just put into my hand and I have carefully purrused its contents and imbrace this oppertunity to answer it, we proceed to answer your questions first concerning your labour in the region where you live we acquiesce in your feelings on this subject until the the mouth of the Lord shall name and as it respects the vision you speak of we do not consider ourselves bound to receive any revelation forom [from] any one man or woman without being legally constituted and ordained to that authority and given sufficien[t] proof of it,
					\hypertarget{EMS-1833-JS-13-APR-1833-b}%
				I
					\marginnote{b}%
				will inform you that it is contrary to the economy of God for any member of the Church or \sout{any} \sout{<one>} any one to receive instruction for those in authority hig[h]er than themselves, therefore you will see the impropriety of giving <heed> to them, but if any <have a vision> \sout{heavenly} or a visitation from an hevenaly [heavenly] messenger it must be for their own benefit and instruction,
					\hypertarget{EMS-1833-JS-13-APR-1833-c}%
				for
					\marginnote{c}%
				the fundimental principals, government and doctrine of the church is invested in the keys of the kingdom as it respects an apostate or one who has been cut off from the Church and wishes to come in again the law of our church expresly says that such shall repent, and be babtised and be admited the same as at the first, the duty of a high priest is to administer spriitual and holy things and to hold Communeion <with> God but not to exorcise [exercise] \sout{monarchy} monarchal government or to appoint meetings for the Elders without their concent and again it is the high priests duty to be bet[t]er qualifide to teach principles and doctrines than the Elder
					\hypertarget{EMS-1833-JS-13-APR-1833-d}%
				for
					\marginnote{d}%
				the office of Elders is an appendege to the high priesthood and it <centers \&> concentrates in one, and again the \sout{proper} <process> \sout{way} of Labouring with a member we are to deal with them percisely as the scripturs direct if thy brother trespass against the[e] take him betwen him and thee alone and if he maketh the satisfaction thou hast saved thy brother and if not proce[e]d to take another with the[e] \& \textit{[illegible]} when there is no bishop they are to be tried by the voice of the Church
					\hypertarget{EMS-1833-JS-13-APR-1833-e}%
				and
					\marginnote{e}%
				if an Elder or an high priest be presant they are to take the lead in managing the business if not by such that have the highest authority, with respect to preparing to go to Zion first it would be pleasing to the lord that \sout{that} the Church or Churches going to Zion should be organised, and appoint\sout{ed} a suitable person who is well acquainted with the Conditions of the Church \& and be sent to Kirtland to inform the Bishop and procure licence from him agreeable to \sout{the} revelation
					\hypertarget{EMS-1833-JS-13-APR-1833-f}%
				so
					\marginnote{f}%
				doing you will prevent confusion and disorder and escape many difficulties that attend an unorganised band in Journey[ing] \sout{to} in the last days and again those in debt should in all cases pay their debts and the rich \sout{in} are in no wise to cast out the poor or leave them behind for it is said that the poor shall inherit the earth you quoted a pass[a]ge in Jeremiah with regard to Journey to Zion the word of God stands sure so let it be done
				
					\hypertarget{EMS-1833-JS-13-APR-1833-g}%
				There
					\marginnote{g}%
				are two paragraphs in your letter which I do not commend as they are writen blind speaking of the Elder being sent like lightning from the bow of Judah the second, no secret in the councils of Zion you mention this as if fear rested upon your mind otherwise we cannot understand it and again we never enquire \sout{of} at the hand of \sout{the Lord} God for special revelation only in case of ther being no previous revelation to suit the case and that in a court of high Priests for further information on the subject you have writen
					\hypertarget{EMS-1833-JS-13-APR-1833-h}%
				I
					\marginnote{h}%
				will refer you to the Elders who have recently left here for the east by commandment some of whom you will probably see soon you may depend on any information you receive from them that are faithful you may expect to see Bro Orson [Pratt] \& Lyman [Johnson] for whom we have great fellowship, it <is> a great thing to enquire at the hand of God or to come into his presence and we feel fearful to appro[a]ch him upon subject[s] that are of little or no consequen[ce] to satisfy the enqueries of individuals especially <about> things the knowledge of which men aught to obtain in all cencerity before God for themselves in humility by the prayer of faith, and more especially a teacher or a high Priest in the Church of Christ
					\hypertarget{EMS-1833-JS-13-APR-1833-i}%
				I
					\marginnote{i}%
				speak these things not by way of reproach but by way of instruction and I speak as being acquainted whereas we are strangers to each other in the flesh I love your soul and the souls of the Children of men and pray and do all I can for the salvation of all I now close by sending you a salutation of peace in the name of the Lord Jesus Christ Amen the blessings of our Lord Jesus Christ be and abide with you all Amen------
				
				\begin{tightright}
					\uline{Joseph} \uline{Smith} \uline{Jr}
				
					F[rederick] G Williams
				\end{tightright}
				
					\hypertarget{EMS-1833-JS-13-APR-1833-j}%
				NB
					\marginnote{j}%
				if it is inconvenient to send a deleget to Kirtland to procure licen[s]e for the brethren to go to Zion it can be done by two or more Elders we have received two letters from bro [Henry G.] Sherwood stating the order and condition of the Church and respecting the vision of his wife but on account of a multitude of business they have not been answered by us you will please to read this letter to Bro Sherwood and for further information appeal to the Elders you will please to inform them at the east who write to us to pay the postage as we are receiving letters from all parts and have to pay a great sum of money otherwise we shall be under the necessity of letting them remain in the office as we are not able to pay so much
				
				\begin{tightright}
					JS
					
					\uline{FGW}
				\end{tightright}
				
		% -----
		\subsubsection{Letter to Church Leaders in Jackson County, Missouri, 21 April 1833}
		% -----
			\label{EMS-1833-JS-21-APR-1833}
			% \label{EMS-1830-JS-DAY-3LETTERMONTH-YEAR}
		
			\begin{description}
				\item[Print Sources]
			\end{description}

				\begin{tabular}{p{0.5in}p{1in}p{0.5in}p{1in}}
				JSP	 	& D3:64--70		& JSR 		& --- \\
				DC	 	& --- 			& HC 		& 1:340--342 \\
				HDDC 	& ---			& TCHC 		& 1:251--253 \\
				\end{tabular}
				% HDDC = woodford's DC diss; JSR = Marquardt's JS Revelations; TCHC = Vogel HC
			
			\begin{description}
				\item[Online Access] \href{https://www.josephsmithpapers.org/paper-summary/letter-to-church-leaders-in-jackson-county-missouri-21-april-1833/1}{Here}
				% \item[Analysis] \hyperref[XX]{Here}
			\end{description}
			
			\begin{description}
				\item[Background]
			\end{description}
			
				% It might also be useful to give a two sentence context of the period: e.g., "This speech was given in Nauvoo to a small congregation one month prior to the destruction of the Nauvoo Expositor. These and other tensions may have been on the prophet's mind when he gave the speech."
			
			\begin{description}
				\item[Original Source]
			\end{description}
			
				% brief description of original source material, with reference to JSP or somewhere else for more info
				% Background material: it might be helpful to have a note that says whether a given document was presented as a speech, was written in a journal, etc.
				% Also a comment here about how reliable the source might be, or what shape it is in, if there are large omissions or parts that are hard to understand, and so on.
			
			\begin{description}
				\item[Text]
			\end{description}
			
				 ...it is not the will of the Lord to print any of the new translation in the Star but when it is published it will all go to the world together in a volume by itself, and the new Testament and the book of Mormon will be printed together...
				 
		% -----
		\subsubsection{Letter to Edward Partridge, 2 May 1833}
		% -----
			\label{EMS-1833-JS-2-MAY-1833}
			% \label{EMS-1830-JS-DAY-3LETTERMONTH-YEAR}
		
			\begin{description}
				\item[Print Sources]
			\end{description}

				\begin{tabular}{p{0.5in}p{1in}p{0.5in}p{1in}}
				JSP	 	& D3:71--77		& JSR 		& --- \\
				DC	 	& --- 		& HC 		& --- \\
				HDDC 	& ---		& TCHC 		& --- \\
				\end{tabular}
				% HDDC = woodford's DC diss; JSR = Marquardt's JS Revelations; TCHC = Vogel HC
			
			\begin{description}
				\item[Online Access] \href{https://www.josephsmithpapers.org/paper-summary/letter-to-church-leaders-in-jackson-county-missouri-21-april-1833/1}{Here}
				% \item[Analysis] \hyperref[XX]{Here}
			\end{description}
			
			\begin{description}
				\item[Background]
			\end{description}
			
				% It might also be useful to give a two sentence context of the period: e.g., "This speech was given in Nauvoo to a small congregation one month prior to the destruction of the Nauvoo Expositor. These and other tensions may have been on the prophet's mind when he gave the speech."
			
			\begin{description}
				\item[Original Source]
			\end{description}
			
				% brief description of original source material, with reference to JSP or somewhere else for more info
				% Background material: it might be helpful to have a note that says whether a given document was presented as a speech, was written in a journal, etc.
				% Also a comment here about how reliable the source might be, or what shape it is in, if there are large omissions or parts that are hard to understand, and so on.
			
			\begin{description}
				\item[Text]
			\end{description}
			
				\begin{tightcenter}
						\hypertarget{EMS-1833-JS-2-MAY-1833-a}%
					Kirtland
						\marginnote{a}
					Ohio May 2\supers{d} 1833
				\end{tightcenter}
				
				Beloved Brother Edward [Partridge],
				
				I commence answering your letter \& sincere request to me, by begging your pardon for not having addressed you, more particularly in letters which I have written to Zion, for I have always felt, as though a letter written to any one in authority in Zion, would be the property of all, \& it mattered but little to whom it was directed. But I am satisfied that this is an error, for instruction that is given pointedly, and expressly to us, designating our names as individuals, seems to have double power and influence over our minds,
					\hypertarget{EMS-1833-JS-2-MAY-1833-b}%
				I
					\marginnote{b}%
				am thankful to the Lord for the testimony of his spirit, which he has given me, concerning your honesty, and sincerity before him, and the Lord loveth you, and also Zion, for he chasteneth whom he loveth, and scourgeth every son \& daughter whom he receiveth, and he, will not suffer you to be confounded, and of this thing you may rest assured,
					\hypertarget{EMS-1833-JS-2-MAY-1833-c}%
				notwithstanding,
					\marginnote{c}%
				all the threatning of the enemy, and your perils among false brethren, For verily I say unto you, that this is my prayer, and I verily believe the prayer of all the saints in Kirtland, recorded in heaven, in these words, Heavenly Father in the name of Jesus Christ thy son, preserve brother Edward, the bishop of thy church, and give him wisdom, knowledge \& power, \& the holy ghost, that he may impart to thy saints in Zion. their inheritance, \& to every man his portion of meat in due season, and now, this is our confidence \& record on high, therefore fear not little flock, for it has been your fathers good will to giv[e] you the king[dom],.
					\hypertarget{EMS-1833-JS-2-MAY-1833-d}%
				and
					\marginnote{d}%
				now, I will proceed to tell you my views, concerning consecration, property, and giving inheritances \&c. The law of the Lord, binds you to receive, whatsoever property is consecrated, by deed, The consecrated property, is considered the residue kept for the Lords store house, and it is given for this consideration, for to purchase inheritaces for the poor, this, any man has a right to do, agreeable to all laws of our country, to donate, give or consecrate all that he feels disposed to give, and it is your duty, to see that whatsoever is given, is given legally, therefore, it must be given for the consideration of the poor saints, and in this way no man can take any advantage of you in law,
					\hypertarget{EMS-1833-JS-2-MAY-1833-e}%
				Again,
					\marginnote{e}%
				concerning inheritances, you are bound by the law of the Lord, to give a deed, secureing to him who receives inheritances, his inheritance, for an everlasting inheritance, or in other words, to be his individual prope[r]ty, his privat ste[wa]rdship, and if he is found a transgressor \& should be cut off, out of the church, his inheritance is his still and he is dilivere[d] over to the buffetings of satan, till the day of redemption, But the property which he consecrated to the poor, for their benefit, \& inheritance, \& stewardship, he cannot obtain again by the law of the Lord,
					\hypertarget{EMS-1833-JS-2-MAY-1833-f}%
				Thus
					\marginnote{f}%
				you see the propriety of this law, that rich men cannot have power to disinherit the poor by obtaining again that which they have consecrated, which is the residue, signified in the law, that you will find in the second paragraph of the extract from the law, in the second number, And now b[r]other Edward, be assured that we all feel thankful, that the brethren in Zion are beginning to humble themselves, \& trying to keep the commandments of the Lord, which is our prayer to God, you may all be able to do, and now, may the grace of God be with all, amen.
				
				\begin{tightright}
					Joseph Smith Jun
				\end{tightright}
				
				The above is a true copy of a letter, directed \& sent, \& subscribed agreeable thereto
				
		% -----
		\subsubsection{Revelation, 6 May 1833 [D\&C 93]}
		% -----
			\label{EMS-1833-JS-6-MAY-1833}
			% \label{EMS-1830-JS-DAY-3LETTERMONTH-YEAR}
		
			\begin{description}
				\item[Print Sources]
			\end{description}

				\begin{tabular}{p{0.5in}p{1in}p{0.5in}p{1in}}
				JSP	 	& D3:83--91			& JSR 		& 237--239 \\
				DC	 	& Section 93		& HC 		& 1:343--346 \\
				HDDC 	& 1207--1221		& TCHC 		& 1:255--256 \\
				\end{tabular}
				% HDDC = woodford's DC diss; JSR = Marquardt's JS Revelations; TCHC = Vogel HC
			
			\begin{description}
				\item[Online Access] \href{https://www.josephsmithpapers.org/paper-summary/revelation-6-may-1833-dc-93/1}{Here}
				% \item[Analysis] \hyperref[DC93]{Here}
			\end{description}
			
			\begin{description}
				\item[Background]
			\end{description}
			
				% It might also be useful to give a two sentence context of the period: e.g., "This speech was given in Nauvoo to a small congregation one month prior to the destruction of the Nauvoo Expositor. These and other tensions may have been on the prophet's mind when he gave the speech."
			
			\begin{description}
				\item[Original Source]
			\end{description}
			
				% brief description of original source material, with reference to JSP or somewhere else for more info
				% Background material: it might be helpful to have a note that says whether a given document was presented as a speech, was written in a journal, etc.
				% Also a comment here about how reliable the source might be, or what shape it is in, if there are large omissions or parts that are hard to understand, and so on.
			
			\begin{description}
				\item[Text]
			\end{description}
			
				\begin{tightcenter}
				------Kirtland May 6--- 1833------
				\end{tightcenter}
				
				Verely thus saith the Lord, it shall come to pass, that evry soul who forsaketh their sins and cometh unto me and calleth on my name and obeyeth my voice and keepeth all my commandments shall see my face and know that I am and that I am the true light that lighteth evry man who cometh into the world; and that I am in the fathe[r] and the father in me and the fathe[r] and I are one the father because he gave me of his fulness and the son becaus I was in the world and made flesh my tabernacle and dwelt among the sons of men I was in the world and received of my father, and the works of him were plainly manifest and John saw and bear record of the fulness of \sout{his} <my> glory and the fulness of Johns reccord is hereafter to be reveiled and he bear reccord saying I saw his glory that he was in the begining before the world was therefore in the beg[inn]ing the word was for he was the word even the messenger of salvation the light and the redeemer of the world the spirit of truth who came into the world becaus the world was made by him and in him was the \sout{<fife>} \sout{light} <Life> of men and the light of men the worlds were made by him men were made by him all things were made by him and through him and of him and I, John bear reccord that I beheld his glory as the glory of the only begotten of the fathe[r] full of grace and truth even the spirit of truth which came and dwelt in flesh and dwelt among us and I John saw that he received not of the fulness at the first but received grace for grace and he received not of the fulness but continued from grace to grace until he received a fulness and thus he was called the son of God because he received not of the fulness at the first and I John bear reccord and lo the heavens were opened and the holy ghost decended upon him in the form of a dove and set upon him and there came a voice out of heaven saying this is my beloved son, and I John bear reccord that he received a fulness of the glory of the \sout{eternal} \sout{God} father and he receivd all power both in heaven and on earth and the glory of the father was with him for he dwelt in him and it shall come to pass that if you are faithful you shall receive the fulness of the reccord of John I give unto you these sayings that you may understand and know how to worship and know what you worship that you may come unto the fathe[r] in my name and in due time receive of his fulness for if you keep my commandments you shall receive of his fulness and be glor[i]fied in me as I am glor[i]fied in the father, therefore I say <unto> you you shall receive grace for grace and now verely I say unto you I was in the begining with the fathe[r] and am the first born and all those who are begotten through me are partakers of the glory of the same and are the church of the first born, ye were also in the begining with the fathe[r] that which is Spirit even the spirit of truth and truth is knowledge of things as they are and as they were and as they are to come and whatsoever is more or less than these is the spirit of that wicked one who was a liar from the begining the spirit of truth is of God, I am the spirit of truth, and John bear reccord of me say<ing> he received a fullness of truth yea even all truth and no man receiveth a fulness unless he keepeth his commandments he that keepeth his commandments receiveth truth and light untill he is glor[i]fied in truth and knoweth all things, man was also in the begining with God, inteligence or the Light of truth was not created or made neith[er] indeed can be all truth is independent in that sphere <in which God has placed it---> to act for itself as all inteligenc[e] also otherwise there is no existance behold here is the agency of man and here is the condemnation of man because that which was from the begining is p[l]ainly manifest unto them and they receive not the light and evry man whose spirit rec[e]iveth not the light for man is spirit the Elements are eternal and spirit and element inseperably connected receiveth a fulness of Joy and when seperated man cannot receiv <a> fulness of Joy the elements are the tabernacle of God, yea man is the tabernacle of God even temples and whatsoeve[r] temple is defiled God shall distroy that temple, the glory of God is inteligence or in other words light \& truth light and truth forsaketh that evil one evry spirit of man was innocent in the begining, and God having redeemed man from, the fall <man> became again in their infant state <innocent> before God and that wicked one cometh and taketh away light and truth through disobeidienc [disobedience] from the children of men and becam[e] of the tradition of their fathers but I have commanded you to bring up your Children in light and truth, but verily I say unto you my servant Frederick [G. Williams] you have continued under this condemnation you have not taught your Children light and truth according to the Commandments and that wicked one hath power as yet over you and this is the caus of your affliction and now a commandment I give unto you and if ye will be delivered you shall set in order your own house for there are many things that are not right in your house verely I say unto my servant Sidney [Rigdon] that in some things he hath not kept the commandments concerning his children therefore firstly set in order thy house, and verely I say unto my servant Joseph (or in othe[r] words I will call you friends) for ye are my friends) and ye shall have an inheritance with me I called you servants for the worlds sake and ye are their servants for my sake and now verely I say unto you Joseph you have not kept the commandments and must needs stand rebuked before the lord your family must needs repent and forsake some things and give more earnest heed unto your sayings or be removed out of their place what I say unto one I say unto all pray always lest that wicked one have power in you and remove you out of your place my servant Newel [K. Whitney] also the Bishop of my Church hath need to be chastened and set in order his family and see that they are more diligent and concerned at home and pray always or they shall be removed out of their place now I say unto you my friends let my servant Sidny go his Journey and make haste and also proclaim the acceptable year of the Lord and the gospel of salvation as I shall give him utterence and by your prayr of faith with one consent I will uphold him and let my servants Joseph \& Frederick make haste also and it shall be given them even according to the prayer of faith
				
				and inasmuch as you keep my sayings you shall not be confounded in this world nor in the world to come \sout{Am}
				
				and verely I say unto you that it is my will that ye should hasten to translate my Scriptures and to obtain a knowledg of history and of Countries and of Kingdoms and of laws, of \sout{man} \sout{\& of God} <God \& man> and all these for the salvation of Zion Amen
				
		% -----
		\subsubsection{Revelation, 1 June 1833 [D\&C 95]}
		% -----
			\label{EMS-1833-JS-1-JUN-1833}
			% \label{EMS-1830-JS-DAY-3LETTERMONTH-YEAR}
		
			\begin{description}
				\item[Print Sources]
			\end{description}

				\begin{tabular}{p{0.5in}p{1in}p{0.5in}p{1in}}
				JSP	 	& D3:104--108		& JSR 		& 239--240 \\
				DC	 	& Section 95 		& HC 		& 1:350--352 \\
				HDDC 	& 1232--1240		& TCHC 		& 1:261 \\
				\end{tabular}
				% HDDC = woodford's DC diss; JSR = Marquardt's JS Revelations; TCHC = Vogel HC
			
			\begin{description}
				\item[Online Access] \href{https://www.josephsmithpapers.org/paper-summary/revelation-1-june-1833-dc-95/1}{Here}
				% \item[Analysis] \hyperref[DC95]{Here}
			\end{description}
			
			\begin{description}
				\item[Background]
			\end{description}
			
				% It might also be useful to give a two sentence context of the period: e.g., "This speech was given in Nauvoo to a small congregation one month prior to the destruction of the Nauvoo Expositor. These and other tensions may have been on the prophet's mind when he gave the speech."
			
			\begin{description}
				\item[Original Source]
			\end{description}
			
				% brief description of original source material, with reference to JSP or somewhere else for more info
				% Background material: it might be helpful to have a note that says whether a given document was presented as a speech, was written in a journal, etc.
				% Also a comment here about how reliable the source might be, or what shape it is in, if there are large omissions or parts that are hard to understand, and so on.
			
			\begin{description}
				\item[Text]
			\end{description}
				
				\begin{tightcenter}
				Kirtland June 1\supers{st.} 1833---
				\end{tightcenter}
				
				Verily thus saith the Lord unto you whom I love, and whom I love I also chasten that their sins may be forgiven, for with the chastisement I prepare a way for their deliverance in all things out of temptation and I have loved you therefore ye must needs be chast[e]ned and stand rebuked before my face. for ye have sinned against me a verry grievous sin in that ye have not considered the great commandment in all things that I have given unto you concerning the building of mine house for the preparation wherewith I deign to prepare mine Apostles to prune my vineyard for the last time that I may bring to pass my strange act that I may pour out my spirit upon all flesh. But behold verily I say unto you there are many who have been ordained among you whom I called but few of them are chosen. they who are not chosen have sinned a verry grievous sin in that they are walking in darkness at noon day, and for this cause I gave unto you a commandment that you should call your solem assembly that your fastings and your mourning might come up into the ears of the Lord of sabaoth which is by interpretation the creator of \sout{all things} the first day the beginning and the end. Yea verily I say unto you I gave unto you a commandment that you should build an house in the which house I design to endow those whom I have chosen with power from on high, for this is the promise of the Father unto you. Therefore, I commanded you to tarry even as mine Apostles at Jerusalem. nevertheless my servants sinned a verry grievous sin and contentions arose in the school of the prophets, which was verry grievous unto me saith your Lord. therefore I sent them forth to be chastened. Verily I say unto you, it is my will that you should build an house. If ye keep my commandments ye shall have power to build it. If ye keep not my commandments the love of the father shall not continue with you therefore ye shall walk in darkness. now here is wisdom and the mind of the Lord, Let the house be built not after the manner of \sout{this} <the> world. for I give not unto you that ye shall live after the manner of the world. Therefore let it be built after the manner which I shall show unto three of you whom ye shall appoint and ordain unto this power and the size thereof shall be fifty and five feet in width and let it be sixty and five feet in length in the inner court thereof, and let the lower part of the inner court \sout{thereof} be dedicated unto me for your sacrament offering and for your preaching and your fasting and your praying and the offering up your most holy desires unto me saith your lord, and let the higher part of the inner court be dedicated unto me for the school of mine Apostles saith Son Ah Man, or in otherwords Alphas, or in other words Omegas even Jesus Christ your lord Amen.------
				
		% -----
		\subsubsection{Revelation, 4 June 1833 [D\&C 96]}
		% -----
			\label{EMS-1833-JS-4-JUN-1833}
			% \label{EMS-1830-JS-DAY-3LETTERMONTH-YEAR}
		
			\begin{description}
				\item[Print Sources]
			\end{description}

				\begin{tabular}{p{0.5in}p{1in}p{0.5in}p{1in}}
				JSP	 	& D3:110--112		& JSR 		& 241 \\
				DC	 	& Section 96 		& HC 		& 1:352--353 \\
				HDDC 	& 1241--1248		& TCHC 		& 1:262--263 \\
				\end{tabular}
				% HDDC = woodford's DC diss; JSR = Marquardt's JS Revelations; TCHC = Vogel HC
			
			\begin{description}
				\item[Online Access] \href{https://www.josephsmithpapers.org/paper-summary/revelation-4-june-1833-dc-96/1}{Here}
				% \item[Analysis] \hyperref[DC96]{Here}
			\end{description}
			
			\begin{description}
				\item[Background]
			\end{description}
			
				% It might also be useful to give a two sentence context of the period: e.g., "This speech was given in Nauvoo to a small congregation one month prior to the destruction of the Nauvoo Expositor. These and other tensions may have been on the prophet's mind when he gave the speech."
			
			\begin{description}
				\item[Original Source]
			\end{description}
			
				% brief description of original source material, with reference to JSP or somewhere else for more info
				% Background material: it might be helpful to have a note that says whether a given document was presented as a speech, was written in a journal, etc.
				% Also a comment here about how reliable the source might be, or what shape it is in, if there are large omissions or parts that are hard to understand, and so on.
			
			\begin{description}
				\item[Text]
			\end{description}
			
				\begin{tightcenter}
				Kirtland June 4\supers{th} 1833---
				\end{tightcenter}
				
				Behold I say unto you here is wisdom whereby ye may know how to act concerning this matter. for it is expedient in me that this stake that I have set for the strength of Zion should be made strong. Therefore let my servant Newel [K. Whitney] take charge of the place which is named among you upon which I design to build mine holy house, and again let it be divided into lots according to wisdom for the benefit of those who seek inheritances as it shall be determined in council among you. Therefore take heed that ye see to this matter, and that portion that is necessary to benefit the firm for the purpose of bringing forth my word to the children of men, for Behold verily I say unto you, this is the most expedient in me that my word should go forth unto the children of men for the purpose of subdueing the hearts of the children of men for your good even so Amen--- and again verily I say unto you it is wisdom and expedient in me that my servant John Johnson whose offering I have accepted and whose prayers I have heared, unto whom I give a promise of Eternal life inasmuch as he keepeth my commandments from hence forth, for he is a descendant of Joseph and a partaker of the blessings of the promise made unto his fathers. Verily I say unto you it is expedient in me that he should become a member of the firm that he may assist in bringing forth my word unto the children of men. Therefore ye shall ordain him unto this blessing, and he shall seek dilligently to take away incumberances that are upon the house named among you that he may dwell there<in> even so Amen------
				
		% -----
		\subsubsection{Letter to Church Leaders in Jackson County, Missouri, 25 June 1833}
		% -----
			\label{EMS-1833-JS-25-JUN-1833}
			% \label{EMS-1830-JS-DAY-3LETTERMONTH-YEAR}
		
			\begin{description}
				\item[Print Sources]
			\end{description}

				\begin{tabular}{p{0.5in}p{1in}p{0.5in}p{1in}}
				JSP	 	& D3:147--158		& JSR 		& --- \\
				DC	 	& --- 				& HC 		& 1:362--368 \\
				HDDC 	& ---				& TCHC 		& 1:272--277 \\
				TPJS 	& 22--25			& 	 		&  \\
				\end{tabular}
				% HDDC = woodford's DC diss; JSR = Marquardt's JS Revelations; TCHC = Vogel HC
			
			\begin{description}
				\item[Online Access] \href{https://www.josephsmithpapers.org/paper-summary/letter-to-church-leaders-in-jackson-county-missouri-25-june-1833/1}{Here}
				% \item[Analysis] \hyperref[XX]{Here}
			\end{description}
			
			\begin{description}
				\item[Background]
			\end{description}
			
				% It might also be useful to give a two sentence context of the period: e.g., "This speech was given in Nauvoo to a small congregation one month prior to the destruction of the Nauvoo Expositor. These and other tensions may have been on the prophet's mind when he gave the speech."
			
			\begin{description}
				\item[Original Source]
			\end{description}
			
				% brief description of original source material, with reference to JSP or somewhere else for more info
				% Background material: it might be helpful to have a note that says whether a given document was presented as a speech, was written in a journal, etc.
				% Also a comment here about how reliable the source might be, or what shape it is in, if there are large omissions or parts that are hard to understand, and so on.
			
			\begin{description}
				\item[Text]
			\end{description}
			
				\begin{tightright}
						\hypertarget{EMS-1833-JS-25-JUN-1833-a}%
					Kirtland
						\marginnote{a}%
					June 25--- 1833
				\end{tightright}
			
					
				Brethren we have received your last containing a number of questions which you desire us to answer, this we do the more readily as we desire with all our hearts, the prosperity of Zion and the peace of her inhabitents for we have as great an intrest in the welfare of Zion as you can have--- First as respects getting the book of Commandments bound we think it is not necessary[.] they will be sold as well without binding and there is no book binder to be had as we know off, nor is there materials to be had for binding without keeping the books too long from circulation------
				
					\hypertarget{EMS-1833-JS-25-JUN-1833-b}%
				With
					\marginnote{b}%
				regard to the books of Mormon which are in the hand of brothe[r] Burket we say to you get them from brothe[r] Burkeet give him receipt for them in the name of the \sout{littery} firm Let Bro. [Sidney] Gilbert pay bro Chapin his mony------
				
					\hypertarget{EMS-1833-JS-25-JUN-1833-c}%
				We
					\marginnote{c}%
				have not found the book of Jasher nor any of the othe[r] lost books mentioned in the bible as yet nor wille we obtain \sout{obtain} them at present--- respecting the Apochraphy [Apocrypha] the Lord Said to us that there were many things in it which were true and ther were many things in it which were not true and to those who desired, it should be given by the spirit to know true from the false, we have received some revelations within a short time back which you will obtain in due time. As soon as we can get time, we will review the manuscripts of the Book of Mormon, after which they will be forwarded to you---
					\hypertarget{EMS-1833-JS-25-JUN-1833-c}%
				We
					\marginnote{c}%
				commend the plan highly of your choossing a teacher to instruct the High Priests that they may be able to silence gainsayers--- Concerning Bishops we recommend the following, to let Bro I[saac] Morley be ordained second Bishop in Zion. And let Bro John Carrel [Corrill] be ordained third--- Let Bro Edward [Partridge] choose as counsellors in their place, Bro Parley P Pratt and Bro Titus Billings, ordaining Bro Billings to the High Priesthood. Let Bro Morley choose for his counsellors Bro Christian Whitmer whom ordain to the High Priesthood and Bro Newel Knights [Knight]--- Let Bro Carrol choose bro Daniel Staunton [Stanton] and Bro Hezekiah Peck for his counsellors, let Bro Hezekiah also be ordained to the high Priesthood---
					\hypertarget{EMS-1833-JS-25-JUN-1833-d}%
				Bro
					\marginnote{d}%
				John Johnson of Hiram has been received as a member of the firm by commandment and has just come to Kirtland to live. As soon as we get a power of Agency signed agreeable to law for Bro Partridge, we will forward it to him, and will immediately expect one from that part of the firm to Bro Newel K Whitney signed in the same manner. We would again say to Bro Edward to be sure to get a form according to law for secureing a gift, we have found by examineing the Law that a gift cannot be retained without this--- The truth triumphs gloriously in the East, Multitudes are embraceing it.
					\hypertarget{EMS-1833-JS-25-JUN-1833-e}%
				I
					\marginnote{e}%
				Sidney [Rigdon] who write this letter in behalf of the Presidency have had the privilege of seeing my aged Mother Baptized into the faith of the Gospel a few weeks since at the advanced age of seventy five, she now resides with me. We send by this mail a draft of the City of Zion with explanations, and a draft of the house to be built immediately in Zion for the presidency as well as all purposes of Religion and instruction. Kirtland, the Stake of Zion, is strengthening continually, when they [the?] enemies look at her, they wag their heads and march along. We anticipate the day when the enemies will have fled away and be far from us.
					\hypertarget{EMS-1833-JS-25-JUN-1833-f}%
				You
					\marginnote{f}%
				will remember that the power of agency must be signed by the wives as well as their husbands, and the wives must be examined separate \& apart from the husband the same as signing a deed, and a specification to that effect inserted at the bottom by the Justice before whom said acknowledgement is made, otherwise the power will be of non[e] effect. Clarrissa Bachellor of Boston wants her paper discontinued because she has gone from the place, and she has turned from the faith. Send a paper to Joshua Baley Andover Vermont
				
					\hypertarget{EMS-1833-JS-25-JUN-1833-g}%
				Should
					\marginnote{g}%
				you not understand the explanations Sent with the drafts you will inform us, so as you may have a propper understanding, for it is meet that all things should be done according to the pattern--- The following errors we have found in the commandments as printed 40\supers{th.} Chap 10\supers{th.} verse third line, instead of corruptable put corrupted 14 verse of the same chapter 5\supers{th.} line instead of respecter to persons, put respecter of persons. 21\supers{st.} verse 2\supers{nd.} line of the same chapter, instead of respecter to, put respecter of 44 Chapter 12 verse last line, instead of hands, put heads---
			
					\hypertarget{EMS-1833-JS-25-JUN-1833-h}%
				Bro
					\marginnote{h}%
				Edward Sir, I proceed to answer your questions concerning the consecration of Property. First, it is not right to condescend to verry great paticulars in takeing inventories. the fact is this, that a man is bound by the law of the church to consecrate to the Bishop before he can be considered a legal heir to the Kingdom of Zion, and this too, without constraint, and unless he does this, he cannot be acknowledged before the Lord on the Church Book. Therefore, to condescend to particulars, I will tell you that every man must be his own judge how much he should receive, and how much he should suffer to remain in the hands of the Bishop.
					\hypertarget{EMS-1833-JS-25-JUN-1833-i}%
				I
					\marginnote{i}%
				speak of those who consecrate more than they need for the support of themselves and family The matter of consecration must be done by the mutual consent of both parties--- For, to give the Bishop power to say how much every man shall have and he be obliged to comply with the Bishops judgment, is giveing to the Bishop more power than a King has and upon the other hand, to let every man say how much he needs and the Bishop obliged to comply with his judgment, is to throw Zion into confusion and make a Slave of the Bishop.
					\hypertarget{EMS-1833-JS-25-JUN-1833-j}%
				The
					\marginnote{j}%
				fact is, there must be a balance or equalibrium of power between the bishop and the people, and thus harmony and good will may be preserved among you. Therefore, those persons consecrating property to the Bishop in Zion, and then receiveing an inheritance back, must show reasonably to the Bishop that he wants as much as he claims. but in case the two parties can<not> come to a mutual agreement, the Bishop is to have nothing to do about receiveing their consecrations and the case must be laid before a council of twelve high Priests, the Bishop not being one of the council, but he is to lay the case before them.
					\hypertarget{EMS-1833-JS-25-JUN-1833-k}%
				Say
					\marginnote{k}%
				to Bro Gilbert that we have no means in our power to assist him in a pecuniary point, as we know not the hour when we shall be Sued for debts which we have contracted ourselves in N[ew] York. Say to him that he must exert himself to the utmost to obtain means himself to replenish his Store for it must be replenished and it is his duty to attend to it. We were not a little surprised to hear that some of our letters of a public nature which we sent for the good of Zion <have been> kept back from the Bishop, this is conduct which we \uline{highly} \uline{disapprobate},
					\hypertarget{EMS-1833-JS-25-JUN-1833-l}%
				Answers
					\marginnote{l}%
				to queries in Bro Phelps [William W. Phelps's] letter of June 4\supers{th.} First in relation to the poor, when the Bishops are appointed according to our reccommendation, it will involve upon them to see to the poor according to the laws of the church. In regard to the printing of the New translation it cannot be done until we can attend to it ourselves, and this we will do as soon as the Lord permit--- As to Bro Frederick [G. Williams], all members of the \uline{United} \uline{Firm} are considered one,
					\hypertarget{EMS-1833-JS-25-JUN-1833-m}%
				The
					\marginnote{m}%
				order of the literary firm is a matter of stewardship which is of the greatest importance and the mercantile establishment God commanded \sout{for} to be devoted to the support thereof, \uline{and} \uline{God} \uline{will} \uline{bring} \uline{every} \uline{transgressor} \uline{into} \uline{judgment}. Say to the Brethren Hulits and to all others that the Lord never authorized them to say that the Devil nor his angels nor the Sons of perdition should ever be restored, for their state of destiny was not revealed to man, is not revealed, nor ever shall be revealed save to those who are made partakers thereof, consequently, those who teach this doctrine have not received it of the Spirit of the Lord,
					\hypertarget{EMS-1833-JS-25-JUN-1833-n}%
				Truly,
					\marginnote{n}%
				Bro Oliver [Cowdery] declared it to be the doctrine of devils. We therefore, command that this doctrine be taught no more in Zion--- We sanction the decission of the Bishop and his council in relation to this doctrine being a bar of communion The number of disciples in K[irtland] is, about 150 We have commenced building the House of the Lord in this place, and it goes on rapidly---
					\hypertarget{EMS-1833-JS-25-JUN-1833-o}%
				Good
					\marginnote{o}%
				news from the East and South, of the success of the Labourers is often saluteing our ears, a general time of health among us, families all well and day and night <we pray> for the salvation of Zion. We deliver Bro Ziba [Peterson] over to the buffetings of Satan in the name of the Lord, that he may learn not to transgress the commandments of God. We conclude our letter by the usual salutation in token of the new and everlasting covenant
				
				We hasten to a close because the mail is just going
				
				\begin{tightright}
					Joseph Smith Jr
					
					Sidney Rigdon
					
					F, G, Williams
					
					Martin Harris
				\end{tightright}
				
					\hypertarget{EMS-1833-JS-25-JUN-1833-p}%
				Orson
					\marginnote{p}%
				Hyde Clk \sout{of} for the presidency
				
				P.S. We feel gratified with the way in which Bro William is conducting the Star at present we hope he will seek to render it more and more interesting. In [re]lation to the size of Bishopricks when Zion is once properly regulated there will be a Bishop to each to each square of the size of the one we send you with this; but at present it must be done according to wisdom. It is needful Brethren that you should be all of one heart and of one mind in doing the will of the Lord. there should exist the greatest freedom and familiarity among the Rulers in Zion. We were exceeding sorry to hear the complaint which was made in Bro Edwards letter that the letters attending the olive leaf had been kept from him as it is meet that he should know all thing[s] in relation to Zion as the Lord has appointed him to be a judge in Zion. We hope dear Brethren that the like circumstance will not take place again---
					\hypertarget{EMS-1833-JS-25-JUN-1833-q}%
				When
					\marginnote{q}%
				we direct letters to Zion to any of the High Priests which pertains to the regulation thereof, we always design that they Should be laid before the Bishop so as to enable him to perform his duty, we say say so much hopeing that it will be received in Kindness and our Brethren be careful about each others feelings and walk in love honoring one another more than themselves as is required of the Lord...
				
		% -----
		\subsubsection{Letter to Church Leaders in Jackson County, Missouri, 2 July 1833}
		% -----
			\label{EMS-1833-JS-2-JUL-1833}
			% \label{EMS-1830-JS-DAY-3LETTERMONTH-YEAR}
		
			\begin{description}
				\item[Print Sources]
			\end{description}

				\begin{tabular}{p{0.5in}p{1in}p{0.5in}p{1in}}
				JSP	 	& D3:165--168		& JSR 		& --- \\
				DC	 	& --- 				& HC 		& 1:368--370 \\
				HDDC 	& ---				& TCHC 		& 1:278--280 \\
				TPJS 	& 25--26			& 	 		&  \\
				\end{tabular}
				% HDDC = woodford's DC diss; JSR = Marquardt's JS Revelations; TCHC = Vogel HC
			
			\begin{description}
				\item[Online Access] \href{https://www.josephsmithpapers.org/paper-summary/letter-to-church-leaders-in-jackson-county-missouri-2-july-1833/1}{Here}
				% \item[Analysis] \hyperref[XX]{Here}
			\end{description}
			
			\begin{description}
				\item[Background]
			\end{description}
			
				% It might also be useful to give a two sentence context of the period: e.g., "This speech was given in Nauvoo to a small congregation one month prior to the destruction of the Nauvoo Expositor. These and other tensions may have been on the prophet's mind when he gave the speech."
			
			\begin{description}
				\item[Original Source]
			\end{description}
			
				% brief description of original source material, with reference to JSP or somewhere else for more info
				% Background material: it might be helpful to have a note that says whether a given document was presented as a speech, was written in a journal, etc.
				% Also a comment here about how reliable the source might be, or what shape it is in, if there are large omissions or parts that are hard to understand, and so on.
			
			\begin{description}
				\item[Text]
			\end{description}
			
					\hypertarget{EMS-1833-JS-2-JUL-1833-a}%
				...we
					\marginnote{a}%
				are excedingly fateagued owing to a great press of business we this day finished the translating of the Scriptures for which we returned gratitude to our heavenly father...
				
					\hypertarget{EMS-1833-JS-2-JUL-1833-b}%
				...As
					\marginnote{b}%
				to the gift of tongues, all we can say is that in this place we have received it as the ancients did we wish you however to be careful lest in this you be deceived guard against evils which may arise from any accounts given of women or otherwise be careful in all things lest any root of bitterness spring up among you and thereby many be defiled. Satan will no doubt trouble you about the Gift of tongues unless you are careful you cannot watch him too closly nor pray to[o] much may the Lord give you wisdom in all things, in a letter mailed last week you will doubtless \sout{see} before you receive this <have> obtained information about the new translation.
					\hypertarget{EMS-1833-JS-2-JUL-1833-c}%
				Consign
					\marginnote{c}%
				the Box of the book of \sout{the} Commandments to NK Whitney \& Co Kirtland Geauga Co Ohio care of Killy and Walworth Cleaveland Cuyahoga County Ohio, I Sidney [Rigdon] write this in great haste in answer to yours to Bro Joseph as I am going off immediately in company with Bro Frederick [G. Williams], to proclaim the gospel
					\hypertarget{EMS-1833-JS-2-JUL-1833-d}%
				we
					\marginnote{d}%
				think of starting to morrow having finished the translation of the bible a few hours since and needing some recreation we know of no way we can spend our time more to divine exceptence then endevoring to build up <his> Zion in these last days as we are not willing to Idle any time which can be spent to useful purpose doors are opening continually for proclaiming the spirit of bitterness among the people is fast subsiding and a spirit of enquiry is taking its place
					\hypertarget{EMS-1833-JS-2-JUL-1833-e}%
				I
					\marginnote{e}%
				proclaimed last Sunday at Chardon our county seat [I] had the court house, there was a general turn out, good attention and a pressing invitation for more meetings which will be granted if the Lord will when we return from this tower [tour], Bro Joseph is going to take a tower [tour] with Bro George James of Brownhelm as soon as he (George) comes to this place, we hope our brethren that the greatest freedom and frankness will exist between you and the Bishop not withholding from each other any information from us but communicate with the greatest freedom lest you should produce evils of a serious \sout{nature} character and the Lord becomes offended
					\hypertarget{EMS-1833-JS-2-JUL-1833-f}%
				for
					\marginnote{f}%
				know assuredly if we by our wickedness bring evil on our own heads the Lord will let us bear it till we get weary and hate eniquity Bro Frederick wants you to say to Bro [John] Burk that the man from whom he expected to get the mill stones has run off so he will not be able to get them but Brother Burk can get them at St Louis of the same mans make
				
					\hypertarget{EMS-1833-JS-2-JUL-1833-g}%
				We
					\marginnote{g}%
				conclude by giving our heartiest approbation to evry measure calculated for the spread of <the> truth in these last days and our strongest desires and sincerest prayers for the prosperity of Zion--- Say to all the brethren and sisters in Zion that they have our hearts our best wishes and the strongest desires of our spirits for their welfare Temporal Spiritual, and Eternal
				
				As ever we salute you in the name of the Lord Jesus Amen
				
				\begin{tightright}
					Sidney Rigdon
					
					Joseph Smith Jr
					
					F. G. Williams---
				\end{tightright}
				
		% -----
		\subsubsection{Letter to John Smith, 2 July 1833}
		% -----
			\label{EMS-1833-JS-2-JUL-1833-JS}
			% \label{EMS-1830-JS-DAY-3LETTERMONTH-YEAR}
		
			\begin{description}
				\item[Print Sources]
			\end{description}

				\begin{tabular}{p{0.5in}p{1in}p{0.5in}p{1in}}
				JSP	 	& D3:168--172		& JSR 		& --- \\
				DC	 	& --- 				& HC 		& 1:370 \\
				HDDC 	& ---				& TCHC 		& 1:277--278 \\
				\end{tabular}
				% HDDC = woodford's DC diss; JSR = Marquardt's JS Revelations; TCHC = Vogel HC
			
			\begin{description}
				\item[Online Access] \href{https://www.josephsmithpapers.org/paper-summary/letter-to-john-smith-2-july-1833/1}{Here}
				% \item[Analysis] \hyperref[XX]{Here}
			\end{description}
			
			\begin{description}
				\item[Background]
			\end{description}
			
				% It might also be useful to give a two sentence context of the period: e.g., "This speech was given in Nauvoo to a small congregation one month prior to the destruction of the Nauvoo Expositor. These and other tensions may have been on the prophet's mind when he gave the speech."
			
			\begin{description}
				\item[Original Source]
			\end{description}
			
				% brief description of original source material, with reference to JSP or somewhere else for more info
				% Background material: it might be helpful to have a note that says whether a given document was presented as a speech, was written in a journal, etc.
				% Also a comment here about how reliable the source might be, or what shape it is in, if there are large omissions or parts that are hard to understand, and so on.
			
			\begin{description}
				\item[Text]
			\end{description}
			
				\begin{tightright}
						\hypertarget{EMS-1833-JS-2-JUL-1833-JS-a}%
					Kirtland
						\marginnote{a}%
					2\supers{d} July 1833---
				\end{tightright}
				
				Brother John Smith--- (2\supers{d})
				
				We <have> just received your letter of the 8\supersu{th} of June, which seems to have been wrttin [written] in a Spirit of Justifycation on your part; You well recollect that previous to your leaveing this place you was tryd before the Bishops Court which found you guilty of misdemeanor--- and desided that you should no longer retain your authaurety in the Church. all of which we \sout{approove} as presidents of of the High Preasthood Sanction. \sout{you name} you nam[e] something in your letter that took place at Bro--- Alneys [Oliver Olney's] in Shalersville on the 27\supers{th}, and 28\supers{th}, of August which we perfectly \sout{Agree} recollect and had you have made such Confession as you was required to at Chipiway [Chippewa], all things would have worked together for <your> good;
					\hypertarget{EMS-1833-JS-2-JUL-1833-JS-b}%
				and
					\marginnote{b}%
				as I told you, but you did not manifest that degree of humility to the bretheren that was required, but remained obstenate, for that <reason> God withdrew his Spirit from you and left <you> in \sout{the} in darkness. \sout{darkness} in your letter you say many hard things against the Bretheren Especially Father Smith Br Reyolds [Reynolds Cahoon] the Bishop \&c--- all which we highly disaprove,--- it seems also that your son Eden [Smith] is confederate with you and needs to repent together with yourself in all humility; before the Lord; or you must expect to be dealt with agreeable to the Laws of the Church.
				
				we say you are no more than a private member in the Church,
				
				\begin{tightright}
				Joseph Smith Jr. Pr
				
				F[rederick] G. Williams, Pd
				\end{tightright}
				
		% -----
		\subsubsection{Letter to Church Leaders in Eugene, Indiana, 2 July 1833}
		% -----
			\label{EMS-1833-JS-2-JUL-1833-CL}
			% \label{EMS-1830-JS-DAY-3LETTERMONTH-YEAR}
		
			\begin{description}
				\item[Print Sources]
			\end{description}

				\begin{tabular}{p{0.5in}p{1in}p{0.5in}p{1in}}
				JSP	 	& D3:172--175		& JSR 		& --- \\
				DC	 	& --- 				& HC 		& 1:371 \\
				HDDC 	& ---				& TCHC 		& 1:280--281 \\
				\end{tabular}
				% HDDC = woodford's DC diss; JSR = Marquardt's JS Revelations; TCHC = Vogel HC
			
			\begin{description}
				\item[Online Access] \href{https://www.josephsmithpapers.org/paper-summary/letter-to-church-leaders-in-eugene-indiana-2-july-1833/1}{Here}
				% \item[Analysis] \hyperref[XX]{Here}
			\end{description}
			
			\begin{description}
				\item[Background]
			\end{description}
			
				% It might also be useful to give a two sentence context of the period: e.g., "This speech was given in Nauvoo to a small congregation one month prior to the destruction of the Nauvoo Expositor. These and other tensions may have been on the prophet's mind when he gave the speech."
			
			\begin{description}
				\item[Original Source]
			\end{description}
			
				% brief description of original source material, with reference to JSP or somewhere else for more info
				% Background material: it might be helpful to have a note that says whether a given document was presented as a speech, was written in a journal, etc.
				% Also a comment here about how reliable the source might be, or what shape it is in, if there are large omissions or parts that are hard to understand, and so on.
			
			\begin{description}
				\item[Text]
			\end{description}
			
				\begin{tightright}
						\hypertarget{EMS-1833-JS-2-JUL-1833-CL-a}%
					Kirtland
						\marginnote{a}
					2 July 1833
				\end{tightright}
				
				Dear Brethren it is truly painful to be under the \sout{painful} necessity of writing on a subject which engages our attention at this time viz the case of John Smith and Eden Smith his son. we have Just received a letter from you concerning their standing in the Church, we do not hold them in fellowship we would inform you that John Smith has been dealt with and his authority taken from him and you are required not to receive his teachings but to treat him as a transgressor until he repents and humbles himself before the Lord to the entire satisfaction of the Church and also you have authority to call a conference and sit in Judgment on Edens case and deal with him as the law directs
					\hypertarget{EMS-1833-JS-2-JUL-1833-CL-b}%
				we
					\marginnote{b}%
				feel to rebuke the Elders of that <branch of the> church of Christ for not magnifying their office and letting the transgressor go unpunished we therefore enjoin upon you to be watchful on your part and search out enequity and put it down wherever it may be found you will see by this brethren that you have authority to sit in council on the Smiths and if found Guilty to deal with them accordingly we have this day directed a letter [to] John Smith thereby making known to him our disaprobation of the course he has persued.---
				
				We commend you [to] God and his grace ever praying he will keep and preserve you blameless till he come.
				
				\begin{tabular}{p{3in}p{2in}}
				Joseph Smith Jr &  \\
				Sidney Rigdon & Presdts of HPHd \\
				F G. William [Frederick G. Williams] &  \\
				\end{tabular}
				
				Postscrip by Br Newel [K. Whitney]---
				
				\begin{tightcenter}
						\hypertarget{EMS-1833-JS-2-JUL-1833-CL-c}%
					Kirtland
						\marginnote{c}
					2d July 1833
				\end{tightcenter}
				
				Dear Brethren---
				
				Yours of the 3d of June came safe to hand the last mail and John Smith which was directed to Br Joseph now my brethren on this sheet you have Bro Joseph sanction to my proceedings and the letter I last wrote you and you will govern yourselves accordingly for you have full power and authority to call the two Bro Smiths to an account for their conduct and unless they repent and make satisfaction not only to your branch of the church but also to this branch they must be cut off from the body for under existing circumstances we have no fellowship with them
					\hypertarget{EMS-1833-JS-2-JUL-1833-CL-d}%
				Bro
					\marginnote{d}%
				John Smith authority as an officer in the church was taken from him before he left and he ought to have given up his License but he went away without doing so and it seams he has made use of it to impose upon you--- as to the two sisters you spoke of in your last. if there is no testamony on either side all you can do is to forbid them to partake of the Sacrament unworthyly and pray much and God will bring all thing to light
				
				\begin{tightright}
					N.K. Whitney Bishop
				\end{tightright}
				
		% -----
		\subsubsection{Classification of Scriptures, not before 17 July 1833}
		% -----
			\label{EMS-1833-JS-17-JUL-1833}
			% \label{EMS-1830-JS-DAY-3LETTERMONTH-YEAR}
		
			\begin{description}
				\item[Print Sources]
			\end{description}

				\begin{tabular}{p{0.5in}p{1in}p{0.5in}p{1in}}
				JSP	 	& D3:176--186		& JSR 		& --- \\
				DC	 	& --- 		& HC 		& --- \\
				HDDC 	& ---		& TCHC 		& --- \\
				\end{tabular}
				% HDDC = woodford's DC diss; JSR = Marquardt's JS Revelations; TCHC = Vogel HC
			
			\begin{description}
				\item[Online Access] \href{https://www.josephsmithpapers.org/paper-summary/classification-of-scriptures-not-before-17-july-1833/1}{Here}
				% \item[Analysis] \hyperref[XX]{Here}
			\end{description}
			
			\begin{description}
				\item[Background]
			\end{description}
			
				% It might also be useful to give a two sentence context of the period: e.g., "This speech was given in Nauvoo to a small congregation one month prior to the destruction of the Nauvoo Expositor. These and other tensions may have been on the prophet's mind when he gave the speech."
			
			\begin{description}
				\item[Original Source]
			\end{description}
			
				% brief description of original source material, with reference to JSP or somewhere else for more info
				% Background material: it might be helpful to have a note that says whether a given document was presented as a speech, was written in a journal, etc.
				% Also a comment here about how reliable the source might be, or what shape it is in, if there are large omissions or parts that are hard to understand, and so on.
			
			\begin{description}
				\item[Text]
			\end{description}
			
					\hypertarget{EMS-1833-JS-17-JUL-1833-a}%
				...Scriptures
					\marginnote{a}
				on High Priesthood

				Genesis.------

				Scriptures relating to the order of the High Priesthood

				Genesis--- Chapt 1 ver 1--- 3 verse--- 4th verse--- 5th verse--- 6 verse--- 7 verse--- 8, 9, 10, 11, 12, 13, 14, 16, 17, 18 19, 20, 21, 22, 23, 24--- 25 verses------

				Genesis II Chapter 1 \sout{verse} <Paragraph> \sout{verse} \sout{7} 8 <Paragraph--->

				III Section 8 Paragraph------

				IV, Section 1 Paragraph--- 2\supers{d} Pr--- 13 p\supers{r.}--- 20 Par

				VII--- Sect 1 Paragrap[h]--- 2 Par 3\supers{rd.} Par------

					\hypertarget{EMS-1833-JS-17-JUL-1833-b}%
				VIII
					\marginnote{b}%
				Section all relative to the High Priesthood to the 10 Par--- and also the 11\supersu{th}\supers{.} 13\supersu{th}\supers{.} 14\supersu{th}\supers{.} 15\supersu{th}\supers{.} 16\supersu{th}\supers{.} 19\supers{th.} Par--- 20\supers{th.} \uline{par}.--- 22\supers{d.} Par--- 23\supers{d.} Par--- 24\supers{th} Pr 25\supers{th.}--- Par--- 26\supers{th.} Par \& 27\supers{th.} Paragraphs

				IX Section 1\supers{st.} Paragraph, 2\supers{d.} Par--- 3\supers{d.} Par--- 4\supers{th.} Par 5\supers{th.} Par--- 6\supers{th.} Par------ 7\supers{th.} Par--- 8--- Par------ 9 Par 10 Par 11 Par,--- 12 Par, 13\supers{th.} Par, \sout{1$\diamond$} 21 Par 22\supers{d.} Par---

				X Section--- 4 Paragraph--- \sout{6} 7\supers{th.} Par--- \sout{9} Par---...

					\hypertarget{EMS-1833-JS-17-JUL-1833-c}%
				...Priesthood
					\marginnote{c}
				after the order of Aaron

				I Section 9 Par

				III Sec--- 9 Par 10 Pa 12th Par 13\supers{th} P---

				IV Sec--- 6--- 7--- 8--- and 12 Par

				V Sec--- 1--- 2--- 3--- \& 4 Paragraph also 5th \& 6 Par

				6--- Section--- 1--- 2d--- 5--- \& 7--- Paragraph

				VIII Section--- 1 Par...
				
		% -----
		\subsubsection{Revelation, 2 August 1833--A [D\&C 97]}
		% -----
			\label{EMS-1833-JS-2-AUG-1833-A}
			% \label{EMS-1830-JS-DAY-3LETTERMONTH-YEAR}
		
			\begin{description}
				\item[Print Sources]
			\end{description}

				\begin{tabular}{p{0.5in}p{1in}p{0.5in}p{1in}}
				JSP	 	& D3:198--203		& JSR 		& 241--243 \\
				DC	 	& Section 97 		& HC 		& 1:400--402 \\
				HDDC 	& 1249--1259		& TCHC 		& 1:308--309 \\
				\end{tabular}
				% HDDC = woodford's DC diss; JSR = Marquardt's JS Revelations; TCHC = Vogel HC
			
			\begin{description}
				\item[Online Access] \href{https://www.josephsmithpapers.org/paper-summary/revelation-2-august-1833-a-dc-97/1}{Here}
				% \item[Analysis] \hyperref[DC97]{Here}
			\end{description}
			
			\begin{description}
				\item[Background]
			\end{description}
			
				% It might also be useful to give a two sentence context of the period: e.g., "This speech was given in Nauvoo to a small congregation one month prior to the destruction of the Nauvoo Expositor. These and other tensions may have been on the prophet's mind when he gave the speech."
			
			\begin{description}
				\item[Original Source]
			\end{description}
			
				% brief description of original source material, with reference to JSP or somewhere else for more info
				% Background material: it might be helpful to have a note that says whether a given document was presented as a speech, was written in a journal, etc.
				% Also a comment here about how reliable the source might be, or what shape it is in, if there are large omissions or parts that are hard to understand, and so on.
			
			\begin{description}
				\item[Text]
			\end{description}
			
				\begin{tightcenter}
				Kirtland 2\supers{d.} of August 1833---
				\end{tightcenter}
				
				The word of the Lord unto Joseph Sidney [Rigdon] and Frederick [G. Williams], verely I say unto you my friends I speak unto you with my voice even the voice of my spirit that I may shew unto you my will concerning your brethren in the Land of Zion many of whom are truly humble and are seeking dilligently to Learn wisdom and to find truth, verely verely I say unto you blessed are all such for they shall obtain for I the Lord sheweth mercy unto all the meek and upon all whomsoever I will that I may be justified when I shall bring them into judgment behold I say unto you concerning [the school] in Zion I the Lord am well pleased that there should be a school in Zion and also with my servant Parley [P. Pratt] for he abideth in me and inasmuch as he continue to abide in me he shall continue to preside over the school in the Land of Zion and I will bless him with a multiplicity of blessings in expounding all scriptures and mysteries to the Edefication of the school and of the church in Zion and to the residue of the school I the Lord am willing to shew mercy nevertheless there are those that must needs be chastened and their work shall be mad[e] known, the ax is laid at the root of the trees and evry tree that bringeth not forth good fruit shall be hewn down and cast into the fire I the Lord hath spoken it, verily I say unto you all among them who know their hearts are honest and are broken and their spirits contrite and are willing to observe their covenants by sacrifice yea every sacrifice which I the Lord shall command they are all accepted of me for I the Lord will cause them to bring forth as a very fruitful tree which is planted in a goodly land by a pure stream that yealdeth much precious fruit, Verely I say unto you that it is my will that an house should be built unto me in the Land of Zion like unto the pattern which I have given you yea let it be built speedely by the tithing of my people behold this is the tithing and the sacrifice which I the Lord require at \sout{the} \sout{hand} \sout{of} there hands that there may be an house built unto me for the salvation of Zion for a place of thanksgiving for all saints and for a place of instruction for all those who are called to the work of the ministry in all their several callings and offices that they may be perfected in the understanding of their ministry in theory in principle and in doctrine in all thing[s] pertaining to the Kingdom of God on the earth the keys of which kingdom have been confered upon you and inas much as my people build an house unto me in the name of the Lord and do not suffer any unclean thing to come into it that it be not defiled my glory shall rest upon it yea and my presence shall be therefor I will come into it and all the pure in heart that shall come into it shall see God but if it be defiled I will not come into it and my glory shall not be there for I will not come into unholy temples and now behold if Zion do these things she shall prosper and spread herself and become very glorious very great and very tereble and the Nations of the earth shall honor her and shall say surely Zion is the city of our God and surely Zion cannot fall neithe[r] be mooved out of her place for God is there and the hand of the Lord is there and he hath sworn by the power of his might to be her salvation and her high tower therefore verily thus saith the Lord let Zion rejoice for this Zion the pure in heart therefore let Zion rejoice while all the wicked shall mourn for behold and lo vengence cometh speedily upon the ungodly as the whirlwind and who shall escape it the Lords scurge shall shall pass over by night and by day and the report thereof shall vex all people yet it shall not be staid until the Lord come for the indignation of the Lord is kindled against their abominations and all their wicked works nevertheless Zion shall escape if she observe to do all things whatsoever I have commanded her but if she observe not to do whatsoeve[r] I have commanded her I will visit her according to all her works with sore afflictions with pestilence with plague with sword with vengence with devouring fire nevertheless let it be read this once in their ears that I the Lord have accepted of their offering but if she sin no more none of these things shall come upon her and I will bless her with blessings and multiply a multiplicity of blessing upon her and upon her generation for ever and ever saith the Lord your God Amen
				
		% -----
		\subsubsection{Revelation, 2 August 1833--B [D\&C 94]}
		% -----
			\label{EMS-1833-JS-2-AUG-1833-B}
			% \label{EMS-1830-JS-DAY-3LETTERMONTH-YEAR}
		
			\begin{description}
				\item[Print Sources]
			\end{description}

				\begin{tabular}{p{0.5in}p{1in}p{0.5in}p{1in}}
				JSP	 	& D3:203--207		& JSR 		& 243--244 \\
				DC	 	& Section 94 		& HC 		& 1:346--347 \\
				HDDC 	& 1222--1231		& TCHC 		& 1:256--257 \\
				\end{tabular}
				% HDDC = woodford's DC diss; JSR = Marquardt's JS Revelations; TCHC = Vogel HC
			
			\begin{description}
				\item[Online Access] \href{https://www.josephsmithpapers.org/paper-summary/revelation-2-august-1833-b-dc-94/1}{Here}
				% \item[Analysis] \hyperref[DC94]{Here}
			\end{description}
			
			\begin{description}
				\item[Background]
			\end{description}
			
				% It might also be useful to give a two sentence context of the period: e.g., "This speech was given in Nauvoo to a small congregation one month prior to the destruction of the Nauvoo Expositor. These and other tensions may have been on the prophet's mind when he gave the speech."
			
			\begin{description}
				\item[Original Source]
			\end{description}
			
				% brief description of original source material, with reference to JSP or somewhere else for more info
				% Background material: it might be helpful to have a note that says whether a given document was presented as a speech, was written in a journal, etc.
				% Also a comment here about how reliable the source might be, or what shape it is in, if there are large omissions or parts that are hard to understand, and so on.
			
			\begin{description}
				\item[Text]
			\end{description}
			
				\begin{tightcenter}
				Kirtland 2\supers{d} August 1833
				\end{tightcenter}
				
				And again verely I say unto you my friends a commandment I give unto you that ye shall commence a work of laying out and preparing a begining and foundation of the city of the stake of Zion here in the land of Kirtland begining at my house and behold it must be done according to the pattern which I have given unto you and let the first lot on the south <be> consecrated unto me for the building of an house for the presidency for the work of the presidency in obtaining revelations and for the work of the ministry of the presidency in all things pertaining to the Church and Kingdom verely I say unto you that it shall be built fifty five by sixty five in the width thereof and in the length thereof in the inner court and there shall be a lower court and an higher court according to the pattern which shall be given unto you hereafter and it shall be dedecated unto the Lord from the foundation thereof according to the order of the Priesthood according to the pattern which shall be given unto you hereafter and it shall be wholly dedecated unto the Lord for the work of the presedency and ye shall not suffer any unclean thing to come into it and my glory shall be there and my presence shall be there but if ther shall come into it any unclean thing my glory shall not be there and my presence shall not come into it and again verily I say unto you the seccond lot on the south shall be dedecated unto me for the building of an house unto me for the work of the printing of the translation of my scripturs and all things whatsoever I shall command you and it shall be fifty five by sixty five in the width therof and in the length there of--- in the inne[r] court and there shall be a lower and an higher court and this house shall be wholly dedecated unto the Lord from the foundation there of for the work of the printing in all thing[s] whatsoever I shall command you to be holy and undefiled according to the pattern in all things as it shall be given unto you and on the third lot shall my servant Hyrum [Smith] receive his inheritance and on the first and seccond lots on the North shall my servants Reynolds [Cahoon] and Jared [Carter] receive their inheritence that they may do the work which I have appointed unto them to be a committe[e] to build mine houses according to the commandment which I the Lord God have given unto you and now I give unto you no more at this time Amen
				
		% -----
		\subsubsection{Revelation, 6 August 1833 [D\&C 98]}
		% -----
			\label{EMS-1833-JS-6-AUG-1833}
			% \label{EMS-1830-JS-DAY-3LETTERMONTH-YEAR}
		
			\begin{description}
				\item[Print Sources]
			\end{description}

				\begin{tabular}{p{0.5in}p{1in}p{0.5in}p{1in}}
				JSP	 	& D3:221--228		& JSR 		& 244--247 \\
				DC	 	& Section 98 		& HC 		& 1:403--406 \\
				HDDC 	& 1260--1275		& TCHC 		& 1:309--311 \\
				\end{tabular}
				% HDDC = woodford's DC diss; JSR = Marquardt's JS Revelations; TCHC = Vogel HC
			
			\begin{description}
				\item[Online Access] \href{https://www.josephsmithpapers.org/paper-summary/revelation-6-august-1833-dc-98/1}{Here}
				% \item[Analysis] \hyperref[DC98]{Here}
			\end{description}
			
			\begin{description}
				\item[Background]
			\end{description}
			
				% It might also be useful to give a two sentence context of the period: e.g., "This speech was given in Nauvoo to a small congregation one month prior to the destruction of the Nauvoo Expositor. These and other tensions may have been on the prophet's mind when he gave the speech."
			
			\begin{description}
				\item[Original Source]
			\end{description}
			
				% brief description of original source material, with reference to JSP or somewhere else for more info
				% Background material: it might be helpful to have a note that says whether a given document was presented as a speech, was written in a journal, etc.
				% Also a comment here about how reliable the source might be, or what shape it is in, if there are large omissions or parts that are hard to understand, and so on.
			
			\begin{description}
				\item[Text]
			\end{description}
			
				\begin{tightcenter}
				Kirtland 6th of August 1833---
				\end{tightcenter}
				
				Verily I say unto you my friends fear not let your hearts be comforted yea rejoice ever more and in evry thing give thanks waiting patiently on the Lord for your prayers have entered into the ears of the of the Lord of sabbaoth and are recorded with this seal and testament the Lord hath sworn and decreed that they shall be granted therefore he giveth this promise unto you with an immutable covenant that they shall be fulfilled and all things wherewith you have been afflicted shall work together for your good and to my names glory saith the Lord and now verely I say unto you concerning the Laws of the land it is my will that my people should observe to do all things whatsoever I command them and that Law of the land which is constitutonal suporting the principles of freedom in maintaning rights and privealiges belonging to all mankind is justifyable before me therefore I the [Lord] justify you and your Brotheren of my Chirch in befriending that Law which [is] the constitutoonal Law of the land and as pertaining to law of man whatsoever is more or less then this cometh of evil I the Lord your God maketh you free therefore ye are free indeed and the law also maketh you free--- nevertheless when the wicked rule the people mourn wherefore honest men and wise men should be sought for dilligently and good men and wise men ye should observe to uphold otherwise whatsoever is less then these cometh of evil and I give unto you a commandment that ye shall forsake all evil and cleave unto all good that ye shall live by evry word that procedeth forth out of the mouth of God for he will give unto the faithful \sout{line} <line> upon line precept upon precept and I will try you and prove you herewith and whoso layeth down his life in my cause for my name sake shall find it again even life eternal therefore be not afraid of your enemies for I have decreed in my heart saith the Lord that I will prove you in all things whether you will abide in my covenant even unto death that ye may be found worthy for if ye will not abide in my covenant ye are not worthy of me therefore renounce war and proclaim peace and seek dilligently to turn the hearts of the Children to their fathers and the hearts of the father to the Children \sout{and the hearts of the father to the Children}, and again the hearts of the Jews unto the prophets and the prophets unto the Jews lest I come and smite the whole earth with a curse <and> \sout{all} all flesh be consumed before me let not your hearts be troubled for in my father['s] house \sout{is} are many mansions and I have prepared a place for <you> and where my father and I am there ye shall be also behold I the Lord am not well pleased with many who are in the church at Kirtland for they do not forsake their sins and their wicked ways the pride of their hearts and their covetiousness and all their detestable things and \sout{the} observe the words of wisdom and eternal Life which I have given unto them verily I say unto you that I the Lord will chasten them \sout{verily} and will do whatsoever I List if they do not repent and observe all things whatsoever I have said unto \sout{you} them and again I say unto you if ye observe to do whatsoevr I command you I the Lord will turn all wrath and indignation from you and the gates of hell shall not prevail against you, Now I speak unto you concerning your families, if men will smite you or your families once and ye bear it patiently and revile not against them neither \sout{revile} seek revenge ye shall be rewarded but if ye bear it not patiently it shall be accounted unto them as being meeted out a just measure unto you and again if your enemies shall smite you a seccond time and you revile not against your \sout{enemies} enemy \& bear it patiently your reward shall be an hundred fold and again if he shall smite you a third time and ye bear it patiently your reward shall be doubled unto you four fold and these three testamonies shall stand against your enemy if he repent not and shall not be blotted out, and now verily I say unto you if that enemy shall escape my vengence that he be not brought into judgment before me then ye shall see to it that ye warn him in my name that he come no more upon you neither upon your families even your children nor your childrens Children unto the third and fourth generation and then if he shall come upon you or your Children or your childrens Children unto the third and forth generation nevertheless thine enemy is in thine hands and and if thou reward him according to his works thou art Justified if he has saugh [sought] thy life and thy life is endangered by him, thine enemy is in thine hands and thou art justified behold this is the Law I gave unto my servant Nephi and thy father Joseph and Jacob Isaac and Abraham and all mine ancient\sout{s} prophets and Apostles and again this is the Law that I gave unto mine ancients that they should not go out unto battle against any Nation Kindred Tongue or people save I the Lord commanded them and if any Nation Tongue or people should proclaim war against them they should first lift a standard of peace unto that people Nation or tongue and if that people did not accept the offering of peace neithr the seccond nor third time they should bring their testamonies before the Lord then the Lord would give unto them a commandment and justify them in going out to battle against that Nation Tongue or people and I the Lord would fight their battles and their childrens battles and their children['s] children['s] until they had avenged themselves \sout{of} upon all their enemies unto the third and forth generation, behold this is an ensample unto all people saith the Lord your God for justification befor me and again verely I say unto you if after thine enemy has come upon thee the first time and he repent and come unto thee praying thy forgivness thou shalt forgive him and shall hold it no more as a testamony against thine enemy and so on unto the seccond and third time and as oft as thine enemy repenteth of the trespass wherewith he has trespassed--- against thee thou shalt forgive him until seventy times seven and if he trespass against thee and repent not the first time nevertheless thou shalt forgive him and if he trespass against thee seccond time and repent not never theless thou shalt forgive him and if he tresspass against thee the third time and repent not thou shalt also forgive him but if he trespass against \sout{him} the forth time thou shalt not forgive him but shall bring these testamones [testimonies] before the Lord and they shall not be blotted out untill he repent and rewards \sout{them} thee four fold in all things wherewith he has trespased against them [thee] and if and if he do this thou shalt forgive him with all thine heart and if he do not this I the Lord will avenge thee of thine \sout{enemies} enemy and [an] hundred fold and upon his children and upon his childrens children of all them that hate me unto the third and forth generation but if the children shall repent or the childrens children and turn unto the Lord their God with all their heart with all their might mind and strenght and restore four fould for all their trespasses wherewith they have trespassed or wherewith their fathers have trespassed or their fathers fathers then thine indignation shall be turned away and vengence shall no more come upon them saith the Lord your God and their trespasses shall neve[r] be brought any more as a testamony before the Lord against them Amen------
				
		% -----
		\subsubsection{Letter to Church Leaders in Jackson County, Missouri, 10 August 1833}
		% -----
			\label{EMS-1833-JS-10-AUG-1833}
			% \label{EMS-1830-JS-DAY-3LETTERMONTH-YEAR}
		
			\begin{description}
				\item[Print Sources]
			\end{description}

				\begin{tabular}{p{0.5in}p{1in}p{0.5in}p{1in}}
				JSP	 	& D3:238--243		& JSR 		& --- \\
				DC	 	& --- 		& HC 		& --- \\
				HDDC 	& ---		& TCHC 		& --- \\
				\end{tabular}
				% HDDC = woodford's DC diss; JSR = Marquardt's JS Revelations; TCHC = Vogel HC
			
			\begin{description}
				\item[Online Access] \href{https://www.josephsmithpapers.org/paper-summary/letter-to-church-leaders-in-jackson-county-missouri-10-august-1833/1}{Here}
				% \item[Analysis] \hyperref[XX]{Here}
			\end{description}
			
			\begin{description}
				\item[Background]
			\end{description}
			
				% It might also be useful to give a two sentence context of the period: e.g., "This speech was given in Nauvoo to a small congregation one month prior to the destruction of the Nauvoo Expositor. These and other tensions may have been on the prophet's mind when he gave the speech."
			
			\begin{description}
				\item[Original Source]
			\end{description}
			
				% brief description of original source material, with reference to JSP or somewhere else for more info
				% Background material: it might be helpful to have a note that says whether a given document was presented as a speech, was written in a journal, etc.
				% Also a comment here about how reliable the source might be, or what shape it is in, if there are large omissions or parts that are hard to understand, and so on.
			
			\begin{description}
				\item[Text]
			\end{description}
			
					\hypertarget{EMS-1833-JS-10-AUG-1833-a}%
				...I
					\marginnote{a}%
				will now proceed to give you some advice concerning your business \& the advice may be relied upon. It is wisdom that you look out another place to locate on; be wise in your selection, \& commence in the best situation you can find, is not the land before you? \& an other place of beginning will be no injury to Zion in the end, \& though you may be wearied, yet count it joy, for the Lord will reward you more than a hundred fold for all your sufferings in righteousness,--- Make out your bill of damages immediately, if you have not, \& get the pay; do not remove any faster to your new home than you bound yourselves to, but pray for the Lord to deliver, for this is his will that you should, \& fear not for his arm will be revealed, \& it will fall upon the wicked \& they cannot escape.
				
					\hypertarget{EMS-1833-JS-10-AUG-1833-b}%
				For
					\marginnote{b}%
				the comfort of those who offered their lives \& made the compromise to remove, I just say that the Lord was well pleased with that act, that is, the agreement to remove, \& there was no other way to save the lives of all the church in Zion, or the most: and any who are dissatisfied with that move, are not right \& have cause to repent, \& call upon the Lord for grace to support them in the moment of tribulation. This great tribulation would not have come upon Zion had it not been for rebelion: Firstly there were rebelions against the one to whom were intrusted the keys, \& from thence it has spread down to the lowest \& least member! not this alone, but those who were void of understanding were continually telling that which was not true, \& putting false coloring to the things of God!
					\hypertarget{EMS-1833-JS-10-AUG-1833-c}%
				I
					\marginnote{c}%
				mean those whose mouths are continually open, \& whose tongues cannot be stayed from tatling! \& the church will never have peace while such remain in her, therefore, brethren purge them out, \& have no confidence in any except such as will lay down their lives for this sacred cause for none others are worthy of it. It was necessary that these things should come upon us: not only justice demands it, but there was no other way to cleanse the church.
					\hypertarget{EMS-1833-JS-10-AUG-1833-d}%
				Fear
					\marginnote{d}%
				not, brethren, the Lord is yet for you \& though the heavens \& the earth pass away, yet the elect will be saved, \& Zion will be the joy of all saints \& they will possess her forever \& ever; \& though the atmosphere looks dark yet the Son of righteousness will soon appear with healing in his wings, \& he will spare his people as a man spareth his own Son who serveth him...
				
		% -----
		\subsubsection{Letter to Church Leaders in Jackson County, Missouri, 18 August 1833}
		% -----
			\label{EMS-1833-JS-18-AUG-1833}
			% \label{EMS-1830-JS-DAY-3LETTERMONTH-YEAR}
		
			\begin{description}
				\item[Print Sources]
			\end{description}

				\begin{tabular}{p{0.5in}p{1in}p{0.5in}p{1in}}
				JSP	 	& D3:258--269		& JSR 		& --- \\
				DC	 	& --- 		& HC 		& --- \\
				HDDC 	& ---		& TCHC 		& --- \\
				\end{tabular}
				% HDDC = woodford's DC diss; JSR = Marquardt's JS Revelations; TCHC = Vogel HC
			
			\begin{description}
				\item[Online Access] \href{https://www.josephsmithpapers.org/paper-summary/letter-to-church-leaders-in-jackson-county-missouri-18-august-1833/1}{Here}
				% \item[Analysis] \hyperref[XX]{Here}
			\end{description}
			
			\begin{description}
				\item[Background]
			\end{description}
			
				% It might also be useful to give a two sentence context of the period: e.g., "This speech was given in Nauvoo to a small congregation one month prior to the destruction of the Nauvoo Expositor. These and other tensions may have been on the prophet's mind when he gave the speech."
			
			\begin{description}
				\item[Original Source]
			\end{description}
			
				% brief description of original source material, with reference to JSP or somewhere else for more info
				% Background material: it might be helpful to have a note that says whether a given document was presented as a speech, was written in a journal, etc.
				% Also a comment here about how reliable the source might be, or what shape it is in, if there are large omissions or parts that are hard to understand, and so on.
			
			\begin{description}
				\item[Text]
			\end{description}
			
				\begin{tightright}
						\hypertarget{EMS-1833-JS-18-AUG-1833-a}%
					Kirtland,
						\marginnote{a}%
					August 18, 1833.
				\end{tightright}
				
				Brother William [W. Phelps], John [Whitmer], Edward [Partridge], Isaac [Morley], John [Corrill] and Sidney [Gilbert].
				
				O thou disposer of all Events, thou dispencer of all good! in the name of Jesus Christ I ask thee to inspire my heart indiht [indite] my thaughts guide my peen [pen] to note some kind word to these my Brotheren in Zion that like the rays of the sun upon the Earth wormeth [warmeth] the face thereof so let this word I write worm [warm] the hearts of my Brotheren or as the gentle rain deceneth [descendeth] upon the earth or the dews upon the mountanes refresheth the face of nature and Causeth her to smile
					\hypertarget{EMS-1833-JS-18-AUG-1833-b}%
				so
					\marginnote{b}%
				give unto thy servent Joseph \sout{have} a word that shall refresh the hearts and revi[v]e the spir[i]ts yea souls <of> those afflicted ones who have been called to leave their homes and go to a strange land not knowing what should befall them behold this is like Abraham a strikeing <evidence> of their acceptance before the <Lord> in this thing but this is not all <they are> \sout{but} called to contend with the beast of the wilderness for a long time whos[e] Jaws <are> \sout{were} open to devour them thus did Abraham and also Paul at Ephesus b[e]hold thou art like \sout{him} <them> and again the affliction of my Brotheren reminds me of Abraham offering up Isaac his only son but my Brotheen [brethren] have have been called to give up even more than this their wives and their children yea and their own life also
					\hypertarget{EMS-1833-JS-18-AUG-1833-c}%
				O
					\marginnote{c}%
				Lord what more dost thou require at their hands before thou wilt come and save them may I not say thou wilt yea I will <say> Lord thou wilt save them out of the hands of their enemies thou hast tried them in the fu[r]nace of affliction a furnace of thine own choseing [choosing]and couldst thou have tried them more then thou hast O Lord then let this suffice and from henceforth <let> this <be> reco[r]ded <be> in heaven for thine angels to look upon and for a testimony against all those ungodly men who have commited those ungodly deeds forever and ever
					\hypertarget{EMS-1833-JS-18-AUG-1833-d}%
				and
					\marginnote{d}%
				<yea> \sout{let} thine anger <is> \sout{be} enkindled against them \sout{and <let> them} <and they shall> be consumed before thy face and be far removed from Zion O <they will \sout{go}> \sout{let} \sout{them} go down to <the> pit and give pl[a]ce for thy saints for thy spirit will not always strive with man therefore I fear for all these things yet O Lord glorify thyself thy will be done and not mine but I must conclude my pray[er] my heart being full of real desire for all such are not reprobate that they cannot be saved------
				
					\hypertarget{EMS-1833-JS-18-AUG-1833-e}%
				Dear
					\marginnote{e}%
				Brotheren in fellowship and <love> towards you and with a broken heart and a contrite Spirit I take the pen to address you but I know not what to--- say to you and the thaught <that> this \sout{<of>} letter will be so long coming to you my heart faints within me I feel to exclaim O Lord let the desire of my heart be felt and realizied this moment <upon you hearts> and teach you all things thy servent would communicate to \sout{would} you my Brotheren since the inteligence of the Calamity of Zion has reached the ears of the wicked there is no saifty for us here but evevery man has to wa[t]ch their houses every night to keep off the Mob[b]ers
					\hypertarget{EMS-1833-JS-18-AUG-1833-f}%
				Satan
					\marginnote{f}%
				has Come down in Great wrath upon all the Chirch of God and the[re] is no saifty only in the arm of Jehovah none else can deliver and he will not deliver unless we \sout{do} prove ourselves faithful to him in the severeest trouble for he that will have his robes washed in the blood of the Lamb must come up throught great tribulation even the greatest of all affliction but know this when men thus deal with you and speak all maner of evil of you falsly for the sake of Christ that he is your friend and I verily know that he will spedily deliver Zion for I have his immutible covenant that this shall be the case 
					\hypertarget{EMS-1833-JS-18-AUG-1833-g}%
				but
					\marginnote{g}%
				god is pleased to keep it hid from mine eyes the means how exactly the thing will be done the chirch in Kirtland concluded with one accord to die with you or redeem you and never at any time have I felt as I now feel that pure love and for you my Brotheren the wormth [warmth] and Zeal for you saf[e]ty that we can scarcely hold our spirits but wisdom I trust will keep us from madness and desperation and the power of the Go[s]pel will enable us to stand and and bear with patience the Great affliction that is falling upon us on all side[s] for we <are> no safer here in Kirtland then you are in Zion
					\hypertarget{EMS-1833-JS-18-AUG-1833-h}%
				the
					\marginnote{h}%
				cloud is gethering arou[nd] us with great fury and all pharohs host or in other words all hell and the com[bined] pow[e]rs of Earth are Marsheling their forces to overthrow us and we like the chilldrn [children] of Issarel [Israel] with the red Sea before \sout{us} them and the Egyptions ready to fall upon them to distroy them and no arm could deliver but the arm of God and this is the case with us we must wait on God to be gratious and call on him with out ceaseing to make bare his arm for our defence for naught but the arm of the almighty can Save us we are all well here as can be expected yea altogether so with the exception of some little ailments feavers \&c.------
				
					\hypertarget{EMS-1833-JS-18-AUG-1833-i}%
				Brother
					\marginnote{i}%
				Oliver [Cowdery] is now Sitting before me and is faithful and true and his heart bleeds as it were for Zion yea never did the hart pant for the cooling streem as doth the heart of thy Brothe[r] Oliver for thy salvation yea and I \sout{may} may say this is the Case with the whole Chirch and all the faithful Oliver will or aught rather to stay with me or in this land until I am permitted to Come with him for I know that if God shall spare my life that he will permit me to settle on an inhe[r]itance on the land of Zion <in due time> but when I do not know but this I do know that I have been keept from going <up> as yet for your sa[k]es and the day will come that Zion will be keept for our sakes therefore be of good cheer and the cloud shall pass over and the sun shall shine as clear and as fair as heaven itself and the Event shall be Glorious
					\hypertarget{EMS-1833-JS-18-AUG-1833-j}%
				Oliver
					\marginnote{j}%
				can stay here to good advantage and have his wife come to him and he can be instrumental of doing great good in this pla[ce] and god will <give> Brother william [W. Phelps] more help and Grace to stand as an\sout{d} ensign to the people for it must be lifted up--- and cursed sha[ll] every man be that lifts his arm to <hinder> this great work and god is my witness of this truth it shall be done and let all the saints say amen------
				
					\hypertarget{EMS-1833-JS-18-AUG-1833-k}%
				Dear
					\marginnote{k}%
				Brotheren we must wait patiently until the Lord come[s] and resto[res] unto us all things and build the waist places again for he will do it in his time and now what shall I say to cumfort your hearts well I will tell you that you have my whole confidence yea there is not one doubt in <my heart> not one place in me but what is filld with perfect confidence and love for you
					\hypertarget{EMS-1833-JS-18-AUG-1833-l}%
				and
					\marginnote{l}%
				this affliction is sent upon us not for your sins but for the sins of the chirch and that all the ends of the Earth may know that you are not speculiting [speculating] with the$\diamond\diamond$ for Lucre but you are willing to die for the cause you have espoused you know that the chirch have tre[a]ted lightly the commandments of the Lord and for this cause they are not worthy to receive them yet god has suffered it not for your sins but that he might preprare you for a grateer [greater] work that you might be prepared for the endowment from on high
					\hypertarget{EMS-1833-JS-18-AUG-1833-m}%
				we
					\marginnote{m}%
				cast no reflections upon you we are of one heart and one mind on this subject which I speak in the name of the chirch all seem to wax strong as th[e]y see the day <of> tribulation approcing [approaching] and if our kingdom were of this world then we would fight but our weapons are not carnal yet mighty and <will> bind satan ere long under our feet
					\hypertarget{EMS-1833-JS-18-AUG-1833-n}%
				we
					\marginnote{n}%
				shall get a press immediately in this place and print th[e] Star until you can obtain deliverence and git up again if god permit and we believe he will we think it would be wise in yo[u] to try to git influence by offering to print a paper in favor of the goverment as you know we are all friends to the Constitution yea true friends to that Country \sout{we hea} for which our fathers bled in the mean time god will send Embasadors to the authorities of the government and sue for protection and redress that they may be left with out excuse that a ritious [righteous] Judgement might be upon them and thus the testimony of the Kingdom must go unto all and there are many ways that God designs to bring about his ritious purposees and in the day of Judgement he designs to make us the Judges of the \sout{whole} \sout{world} generation in which we live
					\hypertarget{EMS-1833-JS-18-AUG-1833-o}%
				O
					\marginnote{o}%
				how unsearchable are the depths of his mysteries and his ways past finding out Brotheren the testamony which you have given of your honesty and the truth of this work will be felt Eteaaly [eternally] by this generation for it will be proclaimed to Ends of the Earth that there are men now liveing who have offered up their lives \sout{for} \sout{this} as a testimony of their religion our Brotheren in the East will handle this testimony to good advantage it seems to inspire every heart to a lively sence of faith and to arm them <with> double fortitude and power and the harder the persicution the greater the gifts of God upon his chirch yea all things shall work together for good to them who are willing to lay down their lives for Christ sake
					\hypertarget{EMS-1833-JS-18-AUG-1833-p}%
				we
					\marginnote{p}%
				are suffering great persicution on account of one man by the name of Docter Hurlburt [Doctor Philastus Hurlbut] who has been expeled from the chirch for lude and adulterous conduct and to spite us he is lieing in a wonderful manner and the peapl [people] are running after him and giveing him mony to b[r]ake  down mormanism which much endangers <our lives> at preasnt [present] but god will put a stop to his carear soon and all will be well my heart this moment <is made> glad for Zion we have just receivd you[r] letter containing the bond with which our enemies bound themselves \sout{and} to distroy Zion and also the blessing <of> god in poreing out \sout{upon} his spirit upon you and we have had the word of the Lord that you shall [be] deliverd from you[r] dainger and <shall> again flurish in spite of hell [that i]s god has communicated to m[e] by the gift of the holy ghost that this should be <the case> after much p[rayer] and suplication
					\hypertarget{EMS-1833-JS-18-AUG-1833-q}%
				and
					\marginnote{q}%
				also that an other printing office must be built the Lord knows how and also it is the will of the Lord that the Store shud [should] be kept and that <not> one foot of <land> \sout{the} perchased should <be> given to the enimies of God or sold to them but if any is sold let it be sold to the chirch we cannot git the consent of the Lord that we shall give the ground to the enemies yet let those who are bound to leave the land \sout{to} make a show as if to do untill the Lord delivr[.]
					\hypertarget{EMS-1833-JS-18-AUG-1833-r}%
				a
					\marginnote{r}%
				word to the wise is sufficient therefore Jud[g]e what I say for know assuredly that every foot of ground that falls into the hands of the enimies with consent is not easy to be obtained again O be wise and not let the knowledge I give unto <you> be known abroad for your sak[e]s hold fast that which you have received trust in god considder Elijah when he prayed for rain go often to your holy plases and <look> for a cloud of light to apper to your help
					\hypertarget{EMS-1833-JS-18-AUG-1833-s}%
				O
					\marginnote{s}%
				God I ask thee in the name of Jesus of nazereth to Save all things concerning Zion and build up her wait [waste] places and restore all things O god send forth Judgement unto victory O come down and cause the moutans [mountains] to flow down at thy presance and now I conclude by telling you that we w[a]it the Comand of God to do whatever \sout{we} he ple[a]se and if <he> shall say go up to Zion and defend thy Brotheren by <the sword> we fly and we count not \sout{dear} our live dear to us I am your Brother in Christ
				
				\begin{tightright}
					Joseph Smith Jr
				\end{tightright}
				
				Edward Partri[d]ge
				
				Independence Jackson County
				
				Missouri
				
		% -----
		\subsubsection{Letter to Vienna Jaques, 4 September 1833}
		% -----
			\label{EMS-1833-JS-4-SEP-1833}
			% \label{EMS-1830-JS-DAY-3LETTERMONTH-YEAR}
		
			\begin{description}
				\item[Print Sources]
			\end{description}

				\begin{tabular}{p{0.5in}p{1in}p{0.5in}p{1in}}
				JSP	 	& D3:288--296		& JSR 		& --- \\
				DC	 	& --- 				& HC 		& 1:407--409 \\
				HDDC 	& ---				& TCHC 		& 1:313--315 \\
				TPJS 	& 26--27			& 	 		&  \\
				\end{tabular}
				% HDDC = woodford's DC diss; JSR = Marquardt's JS Revelations; TCHC = Vogel HC
			
			\begin{description}
				\item[Online Access] \href{https://www.josephsmithpapers.org/paper-summary/letter-to-vienna-jaques-4-september-1833/1}{Here}
				% \item[Analysis] \hyperref[XX]{Here}
			\end{description}
			
			\begin{description}
				\item[Background]
			\end{description}
			
				% It might also be useful to give a two sentence context of the period: e.g., "This speech was given in Nauvoo to a small congregation one month prior to the destruction of the Nauvoo Expositor. These and other tensions may have been on the prophet's mind when he gave the speech."
			
			\begin{description}
				\item[Original Source]
			\end{description}
			
				% brief description of original source material, with reference to JSP or somewhere else for more info
				% Background material: it might be helpful to have a note that says whether a given document was presented as a speech, was written in a journal, etc.
				% Also a comment here about how reliable the source might be, or what shape it is in, if there are large omissions or parts that are hard to understand, and so on.
			
			\begin{description}
				\item[Text]
			\end{description}
			
				\begin{tightcenter}
						\hypertarget{EMS-1833-JS-4-SEP-1833-a}%
					Kirtland
						\marginnote{a}
					Sept 4th 1833------
				\end{tightcenter}
				
				Dear Siste[r]
				
				Having a few Leisur moments I sit down to communicate to you a few wordes which I know I am under obligation to improve \sout{for} to your Satisfaction if it should be a satisfaction for you to receive a few words from your unworthy brother in Christ, I received your Letter some time since containing a history of your Journey and your safe arival for which I bless the Lord
					\hypertarget{EMS-1833-JS-4-SEP-1833-b}%
				I
					\marginnote{b}%
				have often felt a whispering since I received your letter like this Joseph thou art indebted to thy God for the offering of thy Sister Viana [Vienna Jaques] which proved a Savior of life as pertaining to thy pecunary concern therefor she should not be forgotten of thee for the Lord hath done this and thou shouldst remember her in all thy prayers and also by letter for she oftentimes calleth on the Lord saying O Lord inspire thy Servant Joseph to communicate by letter some word to thine unworthy handmaid canst thou not speak peaciably unto thine handmaid
					\hypertarget{EMS-1833-JS-4-SEP-1833-c}%
				and
					\marginnote{c}%
				say all my sins are forgiven and art thou not content with the chastisement wherewith thou hast chastised thy handmaid yea siste[r] this seams to be the whisperings of a spirit and Judge ye what spirit it is I was sensable, when you left Kirtland that the Lord would chasten you but I pray<ed> fervantly in the name of Jesus that you might live to receive your inheritance agreeable to the commandm[e]nt which was given concerning you
					\hypertarget{EMS-1833-JS-4-SEP-1833-d}%
				I
					\marginnote{d}%
				am not at all astonished at what has happened to you neithe[r] to what has happened to Zion and I could tell all the why's \& wherefores of all there calamities but alas it is in vain to warn and give precepts for all men are naturally disposed to walk in their own paths as they are pointed out by their own fingers and are not willing to considder and walk in the path which is pointed out by another saying this is the way walk ye in it altho he should be an uner[r]ing director and the Lord his God sent him nevertheless I do not feel disposed to cast any reflections but I feel to cry mightily unto the Lord that all things might work together for good which has happened
					\hypertarget{EMS-1833-JS-4-SEP-1833-e}%
				yea
					\marginnote{e}%
				I feel to say O Lord let Zion be comforted let her waste places be built up and established an hundred fold let thy saints come unto Zion out of every nation let her be exalted to the third heavens and let thy Judgments be sent forth unto victory and after this great tribulation let thy blessings fall upon thy people and let thy handmaid live till her soul shall be satisfied in beholding the glory of Zion notwithstanding her present affliction she shall yet arise and put on her beautiful garments and be the Joy and \sout{praise} <glory> of the whole earth
					\hypertarget{EMS-1833-JS-4-SEP-1833-f}%
				therefore
					\marginnote{f}%
				let your heart be comferted live in strict obedience to the commandments of God and walk humble before him and he will exalt thee in his own due time the brethren in this place are gaining ground in spiritual things and are trying to overcome all things that is not well pleasing to their heavenly father \sout{we have common} there has many brethren mooved to this place from different parts of the country so much so that one house is not sufficient to contain them for public worship and we have divided and hold meetings in two sepperete places namely at the school house on the flats and Uncle John Smiths who lives on brother [Joseph] Coes place
					\hypertarget{EMS-1833-JS-4-SEP-1833-g}%
				we
					\marginnote{g}%
				have commenced building the house of the Lord in this place and are making great progress in it so much so that I feel great hopes that by spring it will be finished so that we can have a place to worship where we shall not be molested a few days since Brothe[r] Ball and Siste[r] Elizabeth [Eaton] Chase arived here <from> boston broth[er] ball has gone about three miles from this place to work at his trade and Siste[r] Elizabeth lives with me at present Agnes [Coolbrith] \& Mary [Bailey] Lives with father Smith------
				
					\hypertarget{EMS-1833-JS-4-SEP-1833-h}%
				I
					\marginnote{h}%
				will assure you that the Lord has respect unto the offering you made he is a God that changes not and and his word cannot fail remember what he has said in the book of mormon respecting those who should assist in bringing this work forth we frequently have inteligenc from our elder abroad that are proclaming the word that God is working with them for they have attained to great faith insomuch that signs do follow them that believe Brothe[r] David Pettin [David W. Patten] has Just returned from his tour \sout{from} <to> the east and gives us great satisfaction as to his ministry he has raised up a church of about Eighty <three> members in that part of the country where his friends live in the state of New York
					\hypertarget{EMS-1833-JS-4-SEP-1833-i}%
				many
					\marginnote{i}%
				were healed \sout{by} through his instrumantality several criples were restored as many as twelve that were afflicted came at \sout{at} a time from a distanc[e] to be healed he <and others> administered in the name of Jesus and they were made whole thus you see that the Laborers in the Lords vineyard are Labouring with their mights while the day lasts knowing the night soon cometh wherein no man can work
					\hypertarget{EMS-1833-JS-4-SEP-1833-j}%
				I
					\marginnote{j}%
				wish you to say to brothe[r] [Edward] Partridge that we received his lette[r] of the 13 August directed to Bro Frederick [G. Williams] requesting an explination on the Plan of the house which is to be built in Zion and also of the City Platt that <the> brothern whom we have recently sent to Zion will give them all the information they need about it I have but little time to write at present for I am Labouring on the house of the Lord with my own hands therefor I must bid you farewell and subscribe myself your unworthy brother in christ amen
				
				\begin{tightright}
					[Joseph Smith Jr.]
				\end{tightright}
				
				Viana Jaquish
				
				\begin{tabular}{p{3in}p{1in}}
				Kirtland Mills O & \\
				 & 2$\diamond$37 \\
				Sept 11 & \\
				\end{tabular}
				
				Viana Jaquish
				
				Independence Jackson [County]
				
				Missou[ri]
				
		% -----
		\subsubsection{Letter to Silas Smith, 26 September 1833}
		% -----
			\label{EMS-1833-JS-26-SEP-1833}
			% \label{EMS-1830-JS-DAY-3LETTERMONTH-YEAR}
		
			\begin{description}
				\item[Print Sources]
			\end{description}

				\begin{tabular}{p{0.5in}p{1in}p{0.5in}p{1in}}
				JSP	 	& D3:301--308		& JSR 		& --- \\
				DC	 	& --- 		& HC 		& --- \\
				HDDC 	& ---		& TCHC 		& --- \\
				\end{tabular}
				% HDDC = woodford's DC diss; JSR = Marquardt's JS Revelations; TCHC = Vogel HC
			
			\begin{description}
				\item[Online Access] \href{https://www.josephsmithpapers.org/paper-summary/letter-to-silas-smith-26-september-1833/1}{Here}
				% \item[Analysis] \hyperref[XX]{Here}
			\end{description}
			
			\begin{description}
				\item[Background]
			\end{description}
			
				% It might also be useful to give a two sentence context of the period: e.g., "This speech was given in Nauvoo to a small congregation one month prior to the destruction of the Nauvoo Expositor. These and other tensions may have been on the prophet's mind when he gave the speech."
			
			\begin{description}
				\item[Original Source]
			\end{description}
			
				% brief description of original source material, with reference to JSP or somewhere else for more info
				% Background material: it might be helpful to have a note that says whether a given document was presented as a speech, was written in a journal, etc.
				% Also a comment here about how reliable the source might be, or what shape it is in, if there are large omissions or parts that are hard to understand, and so on.
			
			\begin{description}
				\item[Text]
			\end{description}
			
				\begin{tightright}
						\hypertarget{EMS-1833-JS-26-SEP-1833-a}%
					Kirtland
						\marginnote{a}%
					Mills, Ohio, Sept \sout{20} 26th, 1833
				\end{tightright}

				Respected Uncle Silas:--- It is with feelings of deep interest for the welfare of mankind which fill my mind on the reflection that all were formed by the hand of Him who will call the same to give \sout{and} an impartial account of all their works in that great day to which you and myself in common with them are bound, that I take up my pen and seat myself in an attitude to address a few though imperfect lines to you for your perusal.

					\hypertarget{EMS-1833-JS-26-SEP-1833-b}%
				I
					\marginnote{b}%
				have no doubt but you will agree with me that men will be held accountable for the things they have, and not for the things they have not, or, that all the light and intelligence communicated to them from their Beneficent Creator, whether it is much or little, by the same they in justice will be judged; and that they are required to yield obedience to, and improve upon that, and that only, which is given; for man is not to live by bread alone, but by every word that proceeds out of the mouth of the Lord.

					\hypertarget{EMS-1833-JS-26-SEP-1833-c}%
				Seeing
					\marginnote{c}%
				that the Lord has never given them to understand by anything heretofore revealed that He had ceased to speak, forever, to his creatures, when sought unto in a proper manner, why should it be thought a thing incredible that He should be pleased to speak again, in these last days for their salvation?

				Perhaps you may be surprised at this assertion. That I should say for the salvation of his creatures in these last days, since we have already in our possession a vast volume of his word, which he has previously given

					\hypertarget{EMS-1833-JS-26-SEP-1833-d}%
				But
					\marginnote{d}%
				you will admit that the word spoken to Noah was not sufficient for Abraham, or it was not required of <him> to leave the land of his nativity, and seek an inheritance in a strange country upon the word spoken to Noah, but, for himself he obtained promises from the hand of the Lord, and walked in that perfection that he was called the friend of God.

				Isaac, the promised seed, was not required to rest his hope alone upon the promises made to his father Abraham, but was privileged with the assurance of his approbation in the sight of Heaven, by the direct voice of the Lord to him.

					\hypertarget{EMS-1833-JS-26-SEP-1833-e}%
				If
					\marginnote{e}%
				one man can live upon the revelations to another, might I not with propriety ask, why the necessity then, of the Lord's speaking to Isaac as he did, as is recorded in the twenty sixth chapter of Genesis? For the Lord there repeats, or rather, promises again to perform the oath which he had previously sworn to Abraham, and why this repetition to Isaac? Why was not the first promise as sure for Isaac as it was for Abraham? Was not Isaac Abraham's son, and could he not place implicit confidence in the veracity of his father as <being> a man of God?

					\hypertarget{EMS-1833-JS-26-SEP-1833-f}%
				Perhaps
					\marginnote{f}%
				you may say that he was a very peculiar man, and different from men in these last days, consequently the Lord favored him with blessings, peculiar and different, as he was different from men in this age.

				I admit that he was a peculiar man, and was not only peculiarly blessed, but greatly blessed.

				But all the peculiarity that I can discover in the man, or all the difference between him and men in this age, is, that he was more holy and more perfect before God, and came to Him with a purer heart, and more faith than men in this day.

					\hypertarget{EMS-1833-JS-26-SEP-1833-g}%
				The
					\marginnote{g}%
				same might be said on the subject of Jacob's history. Why was it that the Lord spake to him concerning the same promise, after He had made it once to Abraham, and renewed it to Isaac? Why could not Jacob rest contented upon the word spoken to his fathers? When the time of the promise drew nigh for the deliverance of the children of Israel from the land of Egypt, why was it necessary that the Lord should begin to speak to them?

					\hypertarget{EMS-1833-JS-26-SEP-1833-h}%
				The
					\marginnote{h}%
				promise or word to Abraham, was, that his seed should serve in bondage, and be afflicted, four hundred years, and after that they should come out with great substance. Why did they not rely upon this promise, and when they had remained in Egypt, in bondage, four hundred years, come out, without waiting for further revelations, but act entirely upon the promise given to Abraham that they should come out?

				Paul said to his Hebrew brethren, that God might more abundantly show unto the heirs of promise the immutability of His counsel, He confirmed it by an oath. He also exhorts them, who, through faith and patience inherit the promises.

					\hypertarget{EMS-1833-JS-26-SEP-1833-i}%
				Notwithstanding,
					\marginnote{i}%
				we (said Paul) have fled for refuge to lay hold upon the hope set before us, which hope we have as an anchor to the soul, both sure and stedfast, and which entereth into that within the vail, yet he was careful to press upon them the necessity of continuing on until they, as well as those who then inherited the promises, might have the assurance of their salvation confirmed to them, by an oath from the mouth of Him who could not lie; for that seemed to be the example anciently, and Paul holds it out to his Hebrew brethren as an object attainable in his day.

					\hypertarget{EMS-1833-JS-26-SEP-1833-j}%
				And
					\marginnote{j}%
				why not? I admit that by reading the Scriptures of truth the Saints, in the days of Paul, could learn, beyond the power of contradiction, that Abraham, Isaac, and Jacob, had the promise of eternal life confirmed to them by an oath of the Lord, but that promise or oath was <no> assurance to them of their salvation; but they could by walking in the footsteps and continuing in the faith of their fathers, obtain, for themselves an oath for confirmation that they were meet to be partakers of the inheritance, with the Saints in light.

					\hypertarget{EMS-1833-JS-26-SEP-1833-k}%
				If
					\marginnote{k}%
				the Saints in the days of the Apostles were priviledged to take the Ancients for examples, and lay hold of the same promises, and attain to the same exalted privilege of knowing that their names were written in the Lamb's Book of Life and that they were sealed there as a perpetual memorial before the face of the Most High, will not the same faithfulness, the same purity of heart and the same Faith, bring the same assurance of eternal life, and that in the same manner, to the children of men now in this age of the world?
					\hypertarget{EMS-1833-JS-26-SEP-1833-l}%
				I
					\marginnote{l}%
				have no doubt but that the holy Prophets and Apostles and Saints in ancient days, were saved in the Kingdom of God; neither do I doubt but that they held converse and communion with Him while they were in the flesh, as Paul said to his Corinthian brethren that the Lord Jesus showed Himself to above five hundred Saints at one time after His resurrection. Job said that he knew that his Redeemer lived and that he should see Him in the flesh in the latter days. I may believe that Enoch walked with God and by faith was translated. I may believe that Noah was a perfect man in his generation and also walked with God.
					\hypertarget{EMS-1833-JS-26-SEP-1833-m}%
				I
					\marginnote{m}%
				may believe that Abraham communed with God and conversed with angels. I may believe that Isaac obtained a renewal of the covenant made to Abraham by the direct voice of the Lord. I may believe that Jacob conversed with holy angels, and heard the voice of his Maker, that he wrestled with the angel until he prevailed and obtained the blessing. I may believe that Elijah was taken to Heaven in a chariot of fire with fiery horses. I may believe that the saints saw the Lord and conversed with Him face to face after His resurrection.
					\hypertarget{EMS-1833-JS-26-SEP-1833-n}%
				I
					\marginnote{n}%
				may believe that the Hebrew Church came to Mount Zion, and unto the city of the Living God the Heavenly Jerusalem, and to an innumerable company of angels. I may believe that they looked into eternity, and saw the Judge of all, and Jesus the Mediator of the New Covenant. But will all this purchase an assurance for me, and waft me to the regions of eternal day and seat me down in the presence of the King of Kings with my garments spotless pure and white?

					\hypertarget{EMS-1833-JS-26-SEP-1833-o}%
				Or
					\marginnote{o}%
				must I not rather obtain for myself by my own faith and diligence in keeping the commandments of the Lord, an assurance of salvation for myself? And have I not an equal privilege with the ancient Saints? And will not the Lord hear my prayers and listen to my cries as soon as he ever did to theirs, if I come to him in the manner they did? Or, is he a respecter of persons?

					\hypertarget{EMS-1833-JS-26-SEP-1833-p}%
				So
					\marginnote{p}%
				I must close this subject for want of time, and I may with propriety say at the beginning

				We would be glad to see you in Kirtland, we would be glad to see you embrace the New Covenant and be one with us, we sometimes think you are now one with us in heart.
				
				\begin{tightcenter}
					I remain yours affectionately
				\end{tightcenter}
				
				\begin{tightright}
					Joseph Smith Jun
				\end{tightright}

				To Silas Smith
				
		% -----
		\subsubsection{Revelation, 12 October 1833 [D\&C 100]}
		% -----
			\label{EMS-1833-JS-12-OCT-1833}
			% \label{EMS-1830-JS-DAY-3LETTERMONTH-YEAR}
		
			\begin{description}
				\item[Print Sources]
			\end{description}

				\begin{tabular}{p{0.5in}p{1in}p{0.5in}p{1in}}
				JSP	 	& D3:320--325		& JSR 		& 247--248 \\
				DC	 	& Section 100 		& HC 		& 1:420--421 \\
				HDDC 	& 1282--1289		& TCHC 		& 1:330 \\
				\end{tabular}
				% HDDC = woodford's DC diss; JSR = Marquardt's JS Revelations; TCHC = Vogel HC
			
			\begin{description}
				\item[Online Access] \href{https://www.josephsmithpapers.org/paper-summary/revelation-12-october-1833-dc-100/1}{Here}
				% \item[Analysis] \hyperref[DC100]{Here}
			\end{description}
			
			\begin{description}
				\item[Background]
			\end{description}
			
				% It might also be useful to give a two sentence context of the period: e.g., "This speech was given in Nauvoo to a small congregation one month prior to the destruction of the Nauvoo Expositor. These and other tensions may have been on the prophet's mind when he gave the speech."
			
			\begin{description}
				\item[Original Source]
			\end{description}
			
				% brief description of original source material, with reference to JSP or somewhere else for more info
				% Background material: it might be helpful to have a note that says whether a given document was presented as a speech, was written in a journal, etc.
				% Also a comment here about how reliable the source might be, or what shape it is in, if there are large omissions or parts that are hard to understand, and so on.
			
			\begin{description}
				\item[Text]
			\end{description}
			
				\begin{tightcenter}
				Prereysburgh [Perrysburg] Chatocqua [Chautauque] Co NY
				\end{tightcenter}
				
				Saturday October 12\supersu{th} 1833
				
				Verily thus saith the Lord unto you my friends Sidney [Rigdon] \& Joseph your families are well they <are> in mine hands and I will do with them as seemeth me good for in me there is all power therefore follow me and listen to the council which I shall give unto you behold and lo I have much people in this place in the regeons round about and an effectual door shall be opened in the regeons round about in this eastern land therefore I the Lord have suffired you to come unto this place for thus it was expedient in me for the salvation of souls therefore verely I say unto you lift up your voices unto this people speak the thoughts that I shall put into your hearts and ye shall not be confounded before men for it shall be given you in the very hour yea in the very moment what ye shall say but a commandment I give unto you that ye shall declare whatsoever things ye declare in my name in solemnity of heart in the spirit of meekness in all things and I give unto you this promise that inasmuch as ye do this the holy Ghost shall be shed forth in bearing record unto all things whatsoever ye shall say and it is expedient in me that you Sidney should be spokesman unto this people yea verily I will ordain you unto this calling even to be a spokesman unto my servant Joseph and I will give unto him power to be mighty in testimony and I will give unto the<e> power to be mighty in expounding all schriptures that thou mayest be a spokesman unto him and he shall be a revelator unto thee that thou mayest know the certanty of all things pertaining to the things of my kingdom on the earth Therefore continue your journey and let your hearts rejoice for behold and lo I am with you even unto the end \sout{and}
				
				And now I give unto you a word concerning Zion Zion shall be redeemed altho she is chasened for a little season thy breatheren my servents Orson [Hyde] and John [Gould] are in my hands and inasmuch as they keep my commandments they shall be saved therefore let your hearts be comforted for all things shall work together for good to them that walk uprightly and to the sactifycation [sanctification] of the church for I will raise up unto myself a pure people that will serve me in righteousness and all that call on the name of the Lord and keep his commandments shall be saved even so Amen
				
		% -----
		\subsubsection{Record, 27--28 October 1833}
		% -----
			\label{EMS-1833-JS-27-28-OCT-1833}
			% \label{EMS-1830-JS-DAY-3LETTERMONTH-YEAR}
		
			\begin{description}
				\item[Print Sources]
			\end{description}

				\begin{tabular}{p{0.5in}p{1in}p{0.5in}p{1in}}
				JSP	 	& J1:15--16	& JSR 		& --- \\
				DC	 	& ---	& HC 		& --- \\
				HDDC 	& ---	& TCHC 		& --- \\
				\end{tabular}
				% HDDC = woodford's DC diss; JSR = Marquardt's JS Revelations; TCHC = Vogel HC
			
			\begin{description}
				\item[Online Access] \href{https://www.josephsmithpapers.org/paper-summary/journal-1832-1834/16}{Here}
				% \item[Analysis] \hyperref[XX]{Here}
			\end{description}
			
			\begin{description}
				\item[Background]
			\end{description}
			
				% It might also be useful to give a two sentence context of the period: e.g., "This speech was given in Nauvoo to a small congregation one month prior to the destruction of the Nauvoo Expositor. These and other tensions may have been on the prophet's mind when he gave the speech."
			
			\begin{description}
				\item[Original Source]
			\end{description}
			
				% brief description of original source material, with reference to JSP or somewhere else for more info
				% Background material: it might be helpful to have a note that says whether a given document was presented as a speech, was written in a journal, etc.
				% Also a comment here about how reliable the source might be, or what shape it is in, if there are large omissions or parts that are hard to understand, and so on.
			
			\begin{description}
				\item[Text]
			\end{description}
			
					\hypertarget{EMS-1833-JS-27-28-OCT-1833-a}%
				Sunday
					\marginnote{a}%
				26 [27] held a meeting in Mount plesent [Pleasant]. to a large congregation twelve came forward and was baptized and many more were deeply impressed appointed a meeting for this day monday the 27 [28] at the request of some who desires to be baptized at candle lighting held a meeting for confirmation we broke bread laid on hands for the gift of the holy spirit had a good meeting the spirit was given in great \sout{to} power to some and the rest had great \sout{pease} peace may God carry on his work in this place till all shall know him Amen.
					\hypertarget{EMS-1833-JS-27-28-OCT-1833-b}%
				Held
					\marginnote{b}%
				meeting yesterday at 10 o clock after meeting two came forward and were baptized confirmed them at the watters edge held meeting last evening Ordained br E[leazer] F[reeman] Nickerson to the office of Elder had a good meeting one of the sisters got the gift of toungues which made the saints rejoice may God increse the gifts among them for his sons sake this morning we bend our course for home may the Lord prosper our journey Amen
				
		% -----
		\subsubsection{Record, 5--13 November 1833}
		% -----
			\label{EMS-1833-JS-5-13-NOV-1833}
			% \label{EMS-1830-JS-DAY-3LETTERMONTH-YEAR}
		
			\begin{description}
				\item[Print Sources]
			\end{description}

				\begin{tabular}{p{0.5in}p{1in}p{0.5in}p{1in}}
				JSP	 	& J1:16--18	& JSR 		& --- \\
				DC	 	& ---	& HC 		& --- \\
				HDDC 	& ---	& TCHC 		& --- \\
				\end{tabular}
				% HDDC = woodford's DC diss; JSR = Marquardt's JS Revelations; TCHC = Vogel HC
			
			\begin{description}
				\item[Online Access] \href{https://www.josephsmithpapers.org/paper-summary/journal-1832-1834/19}{Here}
				% \item[Analysis] \hyperref[XX]{Here}
			\end{description}
			
			\begin{description}
				\item[Background]
			\end{description}
			
				% It might also be useful to give a two sentence context of the period: e.g., "This speech was given in Nauvoo to a small congregation one month prior to the destruction of the Nauvoo Expositor. These and other tensions may have been on the prophet's mind when he gave the speech."
			
			\begin{description}
				\item[Original Source]
			\end{description}
			
				% brief description of original source material, with reference to JSP or somewhere else for more info
				% Background material: it might be helpful to have a note that says whether a given document was presented as a speech, was written in a journal, etc.
				% Also a comment here about how reliable the source might be, or what shape it is in, if there are large omissions or parts that are hard to understand, and so on.
			
			\begin{description}
				\item[Text]
			\end{description}
			
				\textbf{November 13\supers{th} nothing \sout{of} of note transpired from the 4\supers{th} of Noveber u[n]til this day in the morning at 4 O\supers{h} clock I was awoke by Brother Davis knocking at <my> door saying Brother Joseph come git <up> and see the signs in the heavens and I arrose and beheld to my great Joy the stars fall from heaven yea they fell like hail stones a litteral fullfillment of the word of God as recorded in the holy scriptures and a sure sign that the coming of Christ is clost at hand Oh how marvellous are thy works Oh Lord and I thank thee for thy me[r]cy unto me thy servent Oh Lord save me in thy kingdom for Christ sake Amen}
				
		% -----
		\subsubsection{Record, 14--19 November 1833}
		% -----
			\label{EMS-1833-JS-14-19-NOV-1833}
			% \label{EMS-1830-JS-DAY-3LETTERMONTH-YEAR}
		
			\begin{description}
				\item[Print Sources]
			\end{description}

				\begin{tabular}{p{0.5in}p{1in}p{0.5in}p{1in}}
				JSP	 	& J1:18--19	& JSR 		& --- \\
				DC	 	& ---	& HC 		& --- \\
				HDDC 	& ---	& TCHC 		& --- \\
				TPJS 	& 30--31 &  &  \\
				\end{tabular}
				% HDDC = woodford's DC diss; JSR = Marquardt's JS Revelations; TCHC = Vogel HC
			
			\begin{description}
				\item[Online Access] \href{https://www.josephsmithpapers.org/paper-summary/journal-1832-1834/21}{Here}
				% \item[Analysis] \hyperref[XX]{Here}
			\end{description}
			
			\begin{description}
				\item[Background]
			\end{description}
			
				% It might also be useful to give a two sentence context of the period: e.g., "This speech was given in Nauvoo to a small congregation one month prior to the destruction of the Nauvoo Expositor. These and other tensions may have been on the prophet's mind when he gave the speech."
			
			\begin{description}
				\item[Original Source]
			\end{description}
			
				% brief description of original source material, with reference to JSP or somewhere else for more info
				% Background material: it might be helpful to have a note that says whether a given document was presented as a speech, was written in a journal, etc.
				% Also a comment here about how reliable the source might be, or what shape it is in, if there are large omissions or parts that are hard to understand, and so on.
			
			\begin{description}
				\item[Text]
			\end{description}
			
					\hypertarget{EMS-1833-JS-14-19-NOV-1833-a}%
				\phantom{ }\textbf{November}
					\marginnote{a}%
				\phantom{ }\textbf{19\supers{th} from the 13\supers{th} u[n]till this date \sout{of} nothing of note has transpired since the great sign in the heavins this day my <hart> is somewhat sorrowfull but feel to trust in the Lord the god of Jacob \sout{I} I have learned in my travels that man is trecheous [treacherous] and selfish but few excepted Brother <Sidney [Rigdon]> is a man whom I love but is not capab[le] of that pure and stedfast love for those who are his benefactors as should \sout{posess} possess the breast of \sout{an man} a Presedent of the chu[r]ch of Christ}
					\hypertarget{EMS-1833-JS-14-19-NOV-1833-b}%
				\phantom{ }\textbf{this}
					\marginnote{b}%
				\phantom{ }\textbf{with some other little things such as a selfish and indipendance of mind which to[o] often manifest distroys the confidence of those who would lay down their lives for him but notwithstanding these things he is <a> very great and good man a man of great power of words and can <gain> the friendship of his hearrers very quick he is a man whom god will uphold if he will continue faithful to his calling}
					\hypertarget{EMS-1833-JS-14-19-NOV-1833-c}%
				\phantom{ }\textbf{O}
					\marginnote{c}%
				\phantom{ }\textbf{God grant that he may for the Lords sake Amen the man who willeth to do well we should extall his virtues and speak not of his faults behind his back a man who willfuly turneth away from his friend without a cause is not \sout{lightly} <easily> \sout{to be} \sout{fogiven} <forgiven.> the kindness of a man <should> \sout{is} never \sout{to} be forgotten that person who never forsaketh his trust should ever have the highest place for regard in our hearts and our love should never fail but increase more and more and this my disposition and sentiment \&c}
					\hypertarget{EMS-1833-JS-14-19-NOV-1833-d}%
				\phantom{ }\textbf{Amen}
					\marginnote{d}%
				\phantom{ }\textbf{Brother Frederick [G. Williams] \sout{is a}\sout{man} \sout{who} who is one of those men in whom I place the greatest confidence and trust for I have found him ever full of love and Brotherly kindness he is not a man of many words but is ever wining because of his constant mind he shall ever have place in my heart and is ever intitled to my confidence} He is perfectly honest and upright, and seeks with all his heart to magnify his presidency in the church of ch[r]ist, but fails in many instances, in consequence of a \sout{lack} <want> of confidence in himself:
					\hypertarget{EMS-1833-JS-14-19-NOV-1833-e}%
				God
					\marginnote{e}%
				grant that he may overcome all evil: Blessed be brother Frederick, for he shall never want a friend; and his generation after him shall flourish. The Lord hath appointed him an inheritance upon the land of Zion. Yea, and his head shall blossom<. And he shall be> as an olive branch that is bowed down with fruit: even so; Amen.
				
					\hypertarget{EMS-1833-JS-14-19-NOV-1833-f}%
				And
					\marginnote{f}%
				again, blessed be brother Sidney, also, notwithstanding he shall be high and lifted up, yet he shall bow down under the yoke like unto an ass that coucheth beneath his burthen; that learneth his master's <will> by the stroke of the rod: thus saith the Lord. Yet the Lord will have mercy on him, and he shall bring forth much fruit; even as the \sout{vin} vine of the choice grape when her clusters are <is> ripe, before the time of the gleaning of the vintage:
					\hypertarget{EMS-1833-JS-14-19-NOV-1833-g}%
				and
					\marginnote{g}%
				the Lord shall make his heart merry as with sweet wine because of him who putteth forth his hand and lifteth him up \sout{from} <out of> deep mire, and pointeth him out the way, and guideth his feet when he stumbles; and humbleth him in his pride. Blessed are his generations. Nevertheless, one shall hunt after them as a man hunteth after an ass that hath strayed in the wilderness, \& straitway findeth him and bringeth him into the fold. Thus shall the Lord watch over his generation that they may be saved: even so; Amen.
				
		% -----
		\subsubsection{Letter to Moses Nickerson, 19 November 1833}
		% -----
			\label{EMS-1833-JS-19-NOV-1833}
			% \label{EMS-1830-JS-DAY-3LETTERMONTH-YEAR}
		
			\begin{description}
				\item[Print Sources]
			\end{description}

				\begin{tabular}{p{0.5in}p{1in}p{0.5in}p{1in}}
				JSP	 	& D3:355--360		& JSR 		& --- \\
				DC	 	& --- 				& HC 		& 1:441--443 \\
				HDDC 	& ---				& TCHC 		& 1:351--353 \\
				TPJS 	& 27--30			& 	 		&  \\
				\end{tabular}
				% HDDC = woodford's DC diss; JSR = Marquardt's JS Revelations; TCHC = Vogel HC
			
			\begin{description}
				\item[Online Access] \href{https://www.josephsmithpapers.org/paper-summary/letter-to-moses-nickerson-19-november-1833/1}{Here}
				% \item[Analysis] \hyperref[XX]{Here}
			\end{description}
			
			\begin{description}
				\item[Background]
			\end{description}
			
				% It might also be useful to give a two sentence context of the period: e.g., "This speech was given in Nauvoo to a small congregation one month prior to the destruction of the Nauvoo Expositor. These and other tensions may have been on the prophet's mind when he gave the speech."
			
			\begin{description}
				\item[Original Source]
			\end{description}
			
				% brief description of original source material, with reference to JSP or somewhere else for more info
				% Background material: it might be helpful to have a note that says whether a given document was presented as a speech, was written in a journal, etc.
				% Also a comment here about how reliable the source might be, or what shape it is in, if there are large omissions or parts that are hard to understand, and so on.
			
			\begin{description}
				\item[Text]
			\end{description}
			
					\hypertarget{EMS-1833-JS-19-NOV-1833-a}%
				...When
					\marginnote{a}%
				I contemplate the rapidity with which the great and glorious day of the coming of the Son of Man advances, when he shall come to receive his saints unto himself where they shall dwell in his presence and be crowned with glory \& immortality; when I consider that soon the heavens are to be shaken, and the earth tremble and reel to and fro; and <that> the heavens are to be unfolded as a scroll when it is \sout{folded} <rolled> up, that every mountain and island are to flee \sout{away} <away> I cry out in my heart, What manner of person ought I to be in all holy conversasion and godliness!
				
					\hypertarget{EMS-1833-JS-19-NOV-1833-b}%
				You
					\marginnote{b}%
				remember the testimony which I bore in the name of the Lord Jesus, concerning the great work which he has brought forth in the last days. You know my manner of communication, how that in weakness and simpleness I declared to you what the Lord had brought forth by the ministering of his holy angels to me, for this generation. I pray that the Lord may enable you to treasure these things up in your mind; for I know that his Spirit will bear testimony to all who seek diligently after knowledge from him. I hope you will search the scriptures, to see whether these things are not also consistant with those things that the ancient prophets and apostles have written.
				
					\hypertarget{EMS-1833-JS-19-NOV-1833-c}%
				I
					\marginnote{c}%
				remember brother Freeman [Eleazer Freeman Nickerson] and Wife, Ranson also, and sister Lydia [Goldthwaite Bailey], and little Charles, with all the brethren and sisters. I intreat for an interest in all your prayers before the throne of mercy in the name of Jesus. I hope that the Lord will grant that I may see you all again, and above all that we may overcome and set down together in the Kingdom of our Father...
				
		% -----
		\subsubsection{Letter to Church Leaders in Geneseo, New York, 23 November 1833}
		% -----
			\label{EMS-1833-JS-23-NOV-1833}
			% \label{EMS-1830-JS-DAY-3LETTERMONTH-YEAR}
		
			\begin{description}
				\item[Print Sources]
			\end{description}

				\begin{tabular}{p{0.5in}p{1in}p{0.5in}p{1in}}
				JSP	 	& D3:360--366		& JSR 		& --- \\
				DC	 	& --- 		& HC 		& --- \\
				HDDC 	& ---		& TCHC 		& --- \\
				\end{tabular}
				% HDDC = woodford's DC diss; JSR = Marquardt's JS Revelations; TCHC = Vogel HC
			
			\begin{description}
				\item[Online Access] \href{https://www.josephsmithpapers.org/paper-summary/letter-to-church-leaders-in-geneseo-new-york-23-november-1833/1}{Here}
				% \item[Analysis] \hyperref[XX]{Here}
			\end{description}
			
			\begin{description}
				\item[Background]
			\end{description}
			
				% It might also be useful to give a two sentence context of the period: e.g., "This speech was given in Nauvoo to a small congregation one month prior to the destruction of the Nauvoo Expositor. These and other tensions may have been on the prophet's mind when he gave the speech."
			
			\begin{description}
				\item[Original Source]
			\end{description}
			
				% brief description of original source material, with reference to JSP or somewhere else for more info
				% Background material: it might be helpful to have a note that says whether a given document was presented as a speech, was written in a journal, etc.
				% Also a comment here about how reliable the source might be, or what shape it is in, if there are large omissions or parts that are hard to understand, and so on.
			
			\begin{description}
				\item[Text]
			\end{description}
			
				\begin{tightcenter}
						\hypertarget{EMS-1833-JS-23-NOV-1833-a}%
					Kirtland, November 23, 1833.
						\marginnote{a}
				\end{tightcenter}
				
				An epistle from a counsel of high priests of the church of christ, organized on the 6th of April, A.D. 1830, to their brethren of the same church, residing at Geneseo, Livingston County, New York:
				
				Dearly beloved brethren,
				
					\hypertarget{EMS-1833-JS-23-NOV-1833-b}%
				It
					\marginnote{b}%
				is with \sout{deep} feelings of <deep> interest for your welfare, that we address ourselves to you by this Epistle, which we send by the hands of our worthy brethren, Orson Pratt, and Lyman Johnson, both personly known to us, whom we recommend to your fellowship, as men of good morals and of firm and unshaken integrity in the gospel of our Lord Jesus Christ, to which ministry they have been called and regularly ordained by the hands of this church, and set apart to this office after having been received into the same by baptism according to the Articles and Covenants thereof.
					\hypertarget{EMS-1833-JS-23-NOV-1833-c}%
				It
					\marginnote{c}%
				is just for us, for your sakes, to say, that our brother Orson Pratt, was one of those who first embraced this gospel, and was soon set apart to the work of the ministry, and during an excessive labor of three <years> has conducted himself with that p[ro]priety, and has made such advances in the knowledge of the doctrine of the Kingdom of Christ, that we recommend him in full confidence as a man cabable of setting in order the ordinances and requisitions of the same.
					\hypertarget{EMS-1833-JS-23-NOV-1833-d}%
				Our
					\marginnote{d}%
				brother Lyman Johnson has labored in the ministry more than two years, during which he has showed himself worthy of the high responsibility, and is justly entitled to the confidence of all the saints with whom he has labored, and is fully qualified to assist our brother Orson Pratt in setting in order all matters of difficulty that may be among you.
				
					\hypertarget{EMS-1833-JS-23-NOV-1833-e}%
				Dear
					\marginnote{e}%
				brethren, we have learned with painful feelings, that divisions and strifes in a degree have made their appearance among you, which evidently is the work of the adversary of our souls, to disaffect your minds toward the truth, \sout{that} <and grieve> the Holy Spirit <that it> withdraws, and leaves you in darkness, to be led captive down to destruction: and with great anxiety of heart we have called upon our heavenly Father in the name of Jesus for you.* We need not prove to you by argument, brethren, that where there <are> contentions, and unbelief in the sacred things communic[ated] to the saints by revelation, that discord, hardness, jealousies, and numberless evils will inevitably insue.
					\hypertarget{EMS-1833-JS-23-NOV-1833-f}%
				When
					\marginnote{f}%
				we reflect upon the holiness and perfections of our great Master, who has opened a way whereby we may come unto him, even by the sacrifice of himself, our hearts melt within us for <his> condescension. And when we reflect also, that he has called us to be perfect in all things, that we may be prepared to meet him in peace when he comes in his glory with all the holy angels, we feel to exhort our brethren with boldness, to be humble and prayerful, to walk indeed as children of the light and of the day, that they may have grace to withstand every temptation, and to overcome every evil in the worthy name of our Lord Jesus Christ. For be assured brethren, that the day is truly near when the Master of the house will rise up and shut the door, and none but such as have on a wedding garment will be permitted to enjoy a se[a]t at the marriage supper!
				
					\hypertarget{EMS-1833-JS-23-NOV-1833-g}%
				Therefore,
					\marginnote{g}%
				dear brethren, we have sent our brethren aforementioned to you, hoping you will receive them in our name, and in the name of the church in Kirtland, and receive their teachings and instructions as from us, for they have been set apart to this work, to act in this author[i]ty, and have received the prayers of this counsel, and this church.
				
				We conclude this short letter by earnestly desiring an interest in your prayers, and commending you to the mercy and favor of our Lord Jesus Christ.
				
				We subscribe ourselves your brethren in the bonds of the new and everlasting Covenant,
				
				\begin{tightright}
					Joseph Smith Jr Moderator of co[u]ncil
				\end{tightright}
				
				Oliver Cowdery, \uline{Clerk}
				
					\hypertarget{EMS-1833-JS-23-NOV-1833-h}%
				* <We
					\marginnote{h}%
				have been informed that when> These two brethren visited you previously, having authority from us to teach you the doctrine of this church, and [to] expound the revelations to your understandings, <that they> learned that brother Ezra Landin [Landon] did not believe <all of> the \sout{Vision given to us of} revelations which had been delivered to this chuch by inspiration by the appointment of heaven. Our brother Orson Pratt while reading the Vision to a certain brother, while in brother Landin's house, was threatened of being turned out of at the door except he should desist.
					\hypertarget{EMS-1833-JS-23-NOV-1833-i}%
				They
					\marginnote{i}%
				then called a counsel of \sout{the} High priests to labor with brother L. who, when on the point of being cut off from the church said that he believed the Vision and would teach it to the church. On the return of said brethren from the east this fall, they learned that brother L did not teach, neither believe the Vision. We have also learned from other brethren that he does not walk worthy of his high calling before the Lord, and without speedy repentance and deep humility will have his office and also membership in this church taken from him.
				
					\hypertarget{EMS-1833-JS-23-NOV-1833-j}%
				We
					\marginnote{j}%
				want you to understand, dear brethren, that the conduct of our brother L. has greatly grieved us, and this church. We want you to understand, that we hold no communion, nor have no fellowship for those who do not believe the book of Mormon, and the revelations which God has given to us in these last days. We are informed that our brother L. endeavors to excuse himself for not believing the Vision, saying, that it is not a revelation, but a Vision. We want you to understand from us, that we pronounce such t$\diamond\diamond\diamond$ings the works of the devil, and are calculated to ensnare the souls of the saints.
					\hypertarget{EMS-1833-JS-23-NOV-1833-k}%
				We
					\marginnote{k}%
				plainly declare, and as men that expect to, and must be judged by the searcher of all hearts, that those who do not believe \sout{and} <all> the revelations and vision given to this church, \sout{by} that they do not believe the book of Mormon, and consequently have no fellowship with us. We write plain, for we are bound so to do, \sout{and we write in the name of our Lord Jesus Christ,} and we hope that what we write may be heeded, for we write in the name of our Lord Jesus Christ, by virtue of our calling in his church.
				
		% -----
		\subsubsection{Letter to Edward Partridge, 5 December 1833}
		% -----
			\label{EMS-1833-JS-5-DEC-1833}
			% \label{EMS-1830-JS-DAY-3LETTERMONTH-YEAR}
		
			\begin{description}
				\item[Print Sources]
			\end{description}

				\begin{tabular}{p{0.5in}p{1in}p{0.5in}p{1in}}
				JSP	 	& D3:366--374		& JSR 		& --- \\
				DC	 	& --- 				& HC 		& 1:448--451 \\
				HDDC 	& ---				& TCHC 		& 1:358--362 \\
				TPJS 	& 31--33			& 	 		&  \\
				\end{tabular}
				% HDDC = woodford's DC diss; JSR = Marquardt's JS Revelations; TCHC = Vogel HC
			
			\begin{description}
				\item[Online Access] \href{https://www.josephsmithpapers.org/paper-summary/letter-to-edward-partridge-5-december-1833/1}{Here}
				% \item[Analysis] \hyperref[XX]{Here}
			\end{description}
			
			\begin{description}
				\item[Background]
			\end{description}
			
				% It might also be useful to give a two sentence context of the period: e.g., "This speech was given in Nauvoo to a small congregation one month prior to the destruction of the Nauvoo Expositor. These and other tensions may have been on the prophet's mind when he gave the speech."
			
			\begin{description}
				\item[Original Source]
			\end{description}
			
				% brief description of original source material, with reference to JSP or somewhere else for more info
				% Background material: it might be helpful to have a note that says whether a given document was presented as a speech, was written in a journal, etc.
				% Also a comment here about how reliable the source might be, or what shape it is in, if there are large omissions or parts that are hard to understand, and so on.
			
			\begin{description}
				\item[Text]
			\end{description}
			
					\hypertarget{EMS-1833-JS-5-DEC-1833-a}%
				...It
					\marginnote{a}%
				appears brethren that the above statements were mostly from reports and no certainty of their being correct. therefore it is difficult for us to advise and can only say that the destenies of all people are in the hands of a Just God and he will do no injustice to any one and this one thing is sure that they who will live Godly in Christ Jesus shall suffer persecution and before their robes are mad[e] white in the blood of the Lamb it is to be expected they will pass through great tribulation according to John the Revelator,
					\hypertarget{EMS-1833-JS-5-DEC-1833-b}%
				I
					\marginnote{b}%
				wish when you receive this letter that you would collect every particular concerning the Mob from the begining and send us a correct statement of fact as they transpired from time to time that we may be enabled to give the public correct information on the subject and inform us also of the situation of the brethren with respect to their means of sustinance \&c I would inform you that it is not the will of the Lord for you to sell your Lands in Zion if means can possably be procured for their sustenance without, evry exertion should be to maintain the cause you have espoused and to contribute to the necessities of one another as much as possable in this your great calamity and remember not to murmur at the dealings of God with his creature
					\hypertarget{EMS-1833-JS-5-DEC-1833-c}%
				you
					\marginnote{c}%
				are not as yet brought into as trying circumstances as were the ancient Prophets \& apostles Call to mind a Daniel the three Hebrew Children, Jeremiah Paul Stephen and many more too numerous to mention who were stoned sawn asunder tempted slain with the sword and wandered about in sheep skins \& goat skins being destitute afflicted tormented of whom the world was not worthy. they wandered in deserts and in mountains and in dens and caves of the earth yet they all obtained a good report through faith and amidst all their afflictions they rejoiced that they were <counted> worthy to receive persecution for Christs sake
					\hypertarget{EMS-1833-JS-5-DEC-1833-d}%
				we
					\marginnote{d}%
				know not what we shall be called to pass through before Zion is delivered and established therefore we have great need to live near to God and always be in strict obedience to all his commandments that we may have a concience void of offense towards God and man, It is your privelege to use every lawful means in your power to seek redress for your grievances of your enemies and prosecute them to the extent of the Law but it will be impossable for us to render you any assistance in a temporal point of view as our means are already exhausted and are deeply in debt and know no means whereby we shall be able to extricate ourselves;
					\hypertarget{EMS-1833-JS-5-DEC-1833-e}%
				The
					\marginnote{e}%
				inhabitants of this county threaten our distruction and we know not how soon they may be permitted to follow the examples of the Missourians but our trust is in God and we are determined by his grace assisting us to maintain the cause and hold out faithful to the end that we may be crowned with crowns of celestial glory and enter into that rest that is prepared for the children of God,
					\hypertarget{EMS-1833-JS-5-DEC-1833-f}%
				we
					\marginnote{f}%
				are now distributing the tipe and calculate to commence setting to day and issue a paper the Last of this week or beginning of next, we wrote to bro Phelps some time since and also sent by bro Orson for the names of the subscribe[r]s for the star which we have not yet received and until we receive them the most of the brethren will be deprived of them and when you receive this if you have not sent them I wish you to attend to it immediately as much inconvenience will attend a delay,
					\hypertarget{EMS-1833-JS-5-DEC-1833-g}%
				We
					\marginnote{g}%
				expect shortly to publish a political paper weekly in favour the present administration, the influential men of that party have offered a liberal patronage to us and we hope to succeed for thereby we can shew the public the purity of our intention in supporting the government under which we live---...
				
		% -----
		\subsubsection{Letter to Edward Partridge and Others, 10 December 1833}
		% -----
			\label{EMS-1833-JS-10-DEC-1833}
			% \label{EMS-1830-JS-DAY-3LETTERMONTH-YEAR}
		
			\begin{description}
				\item[Print Sources]
			\end{description}

				\begin{tabular}{p{0.5in}p{1in}p{0.5in}p{1in}}
				JSP	 	& D3:375--381		& JSR 		& --- \\
				DC	 	& --- 				& HC 		& 1:453--456 \\
				HDDC 	& ---				& TCHC 		& 1:365--369 \\
				TPJS 	& 33--37			& 	 		&  \\
				\end{tabular}
				% HDDC = woodford's DC diss; JSR = Marquardt's JS Revelations; TCHC = Vogel HC
			
			\begin{description}
				\item[Online Access] \href{https://www.josephsmithpapers.org/paper-summary/letter-to-edward-partridge-and-others-10-december-1833/1}{Here}
				% \item[Analysis] \hyperref[XX]{Here}
			\end{description}
			
			\begin{description}
				\item[Background]
			\end{description}
			
				% It might also be useful to give a two sentence context of the period: e.g., "This speech was given in Nauvoo to a small congregation one month prior to the destruction of the Nauvoo Expositor. These and other tensions may have been on the prophet's mind when he gave the speech."
			
			\begin{description}
				\item[Original Source]
			\end{description}
			
				% brief description of original source material, with reference to JSP or somewhere else for more info
				% Background material: it might be helpful to have a note that says whether a given document was presented as a speech, was written in a journal, etc.
				% Also a comment here about how reliable the source might be, or what shape it is in, if there are large omissions or parts that are hard to understand, and so on.
			
			\begin{description}
				\item[Text]
			\end{description}
			
				\begin{tightright}
						\hypertarget{EMS-1833-JS-10-DEC-1833-a}%
					Kirtland
						\marginnote{a}
					Mills Ohio December 10th 1833
				\end{tightright}
				
				Beloved brethren E[dward] Partridge, W[illiam] W. Phelps J[ohn] Whitmer A S[idney] Gilbert J Corril [John Corrill] I[saac] Morley, and all the saints whom it may concern.
				
				This morning the mail brought bros Partridge \& Corrils letters \& also bro Williams, all mailed at Liberty Nov. 19th which gave us the melancholy inteligence of your flight from the land of your inheritance having been driven before the face of your enemies in that place
				
					\hypertarget{EMS-1833-JS-10-DEC-1833-b}%
				From
					\marginnote{b}%
				previous letters we had learned that a number of our brethren have been slain, but we could not learn from those refered to above as there had been but one, that was bro [Andrew] Barber and bro [Philo] Dibble wounded in the bowels, we were thankful to learn that no more were slain, and our daily prayers are, that the Lord will not suffer his saints who have gone up to his land to keep his commandments, to stain his holy mountain with their blood. I cannot learn from any communication by the spirit to me that Zion has forfeited her claim to a celestial crown notwithstanding the Lord has caused her to be thus afflicted; except it may \sout{it may} be some individuals who have walked in disobedience and forsaken the new covenants; all such will be made manifest by their works in due time.
					\hypertarget{EMS-1833-JS-10-DEC-1833-c}%
				I
					\marginnote{c}%
				have always expected that Zion would suffer sore affliction from what I could learn from the commandments which have been given. but I would remind you of a certain clause in one which says that after \uline{much} tribulation cometh the \uline{blessing}. by this and also others, and also one received of late, I know that Zion, in the own due time of the Lord will be redeemed, but how many will be the days of her purification, tribulation and affliction, the Lord has kept hid from my eyes; and when I enquire concerning this subject the voice of the Lord is, Be still, and know that I am God! all those who suffer for my name shall reign with me, and he that layeth down his life for my sake shall find it again.
					\hypertarget{EMS-1833-JS-10-DEC-1833-d}%
				Now
					\marginnote{d}%
				there are two things of which I am ignorant and the Lord will not show me--- perhaps for a wise purpose in himself. I mean in some respects, and they are these, Why God hath suffered so great calamity to come upon Zion; or what the great moving cause of this great affliction is. \sout{These two things} and again by what means he will return her back to her inheritance with songs of everlasting Joy upon her head. These two things brethren, are in part kept back that they are not plainly <shewn unto me. but there are some things that are plainly> manifest, that has incured <th[e]> displeasure of \sout{displeasure} the Almighty.
					\hypertarget{EMS-1833-JS-10-DEC-1833-e}%
				when
					\marginnote{e}%
				I contemplate upon all things that have been manifested, I am sensable that I aught not to murmer and do not murmer only in this, \sout{but} that those who are innocent are compelled to suffer for the iniquities of the guilty; and I cannot account for this, only on this wise, that the saying of the savior has not been strictly observed: If thy right eye offend thee pluck it out. and cast it from thee <or if thy right arm offend thee pluck it of[f] and cast it from thee>.
					\hypertarget{EMS-1833-JS-10-DEC-1833-f}%
				Now
					\marginnote{f}%
				the fact is, if any of the members of our body are disordered, the rest of our body will be effected with them and then all is brought into bondage together. And yet notwithstanding all this, it is with difficulty that I can restrain my feelings; when I know that you my brethren with whom I have had so many happy hours, sitting as it were in heavenly places in Christ Jesus. and also haveing the witness which I feel, and even have felt, of the purity of your motives--- are cast out, and are as strangers and pilgrims on the earth, exposed to hunger, cold, nakedness peril, sword \&c
					\hypertarget{EMS-1833-JS-10-DEC-1833-g}%
				I
					\marginnote{g}%
				say when I contemplate this, it is with difficulty that I can keep from complaining and murmerings against this dispensation; but I am sensible that this is not right and may God grant that notwithstanding your great afflictions and sufferings there may not any thing sepperate us from the Love of Christ. Brethren, when we learn your sufferings it awakens evry sympathy of our hearts; it weighs \sout{us} us down; we cannot refrain from tears \textit{[illegible]} yet we are not able to realize only in part your sufferings.
					\hypertarget{EMS-1833-JS-10-DEC-1833-h}%
				And
					\marginnote{h}%
				I often hear the brethren saying they wish they were with you that they might bear a part of your sufferings; and I myself should have been with you had not God prevented it in the order of his providence, that the yoke of affliction might be less grievous upon you; God having forewarned me concerning these things for your sakes; and also bro Oliver [Cowdery], could not lighten your afflictions by tarrying longer with you, for his presence would have so much the more enraged your enemies; therefore, God hath \sout{deals} dealt mercifully with us. O brethren, let us be thankful that it is as well with us as it is, and we are yet alive that peradvent[u]re, God hath Laid up in store great good for us in this generation, and grant that we may yet glorify his name, I feel thankful that there have no more denied the faith;
					\hypertarget{EMS-1833-JS-10-DEC-1833-i}%
				I
					\marginnote{i}%
				pray God, in the name of Jesus that you all may be kept in the faith, unto the end, let your sufferings be what they may, it is better that you should die in the ey[e]s of God, then that you should give up the Land of Zion, the \sout{inhabitant} inheritances which you have purchased with your monies; for evry man that giveth not up his inheritances, though he should die yet when the Lord shall come, he shall stand upon it, and with Job in his flesh he shall see God.
					\hypertarget{EMS-1833-JS-10-DEC-1833-j}%
				Therefore
					\marginnote{j}%
				this is my council that you retain your lands even unto the uttermost, and seeking <evry> lawful means to obtain redress of your enemies \&c \&c and pray to God day and night to return you in peace and in safety to the Lands of your inheritance and, \sout{and} when the Judge fails you, appeal unto the Executive, and when the Executive fails you, appeal unto the President, and when the President fails you, and all Laws fail you and the humanity of the people fails you, and all things else fails you but God alone, and you continue to weary him with your importunings, as the poor woman the unjust Judge, he will not fail to exicute Judgment upon your enemies and to avenge his own elect that cry unto him day and night--- Behold he will not fail you! he will come with ten thousand of his saints and all his advisaries shall be distroyed by the breath of his lips!
					\hypertarget{EMS-1833-JS-10-DEC-1833-k}%
				all
					\marginnote{k}%
				those that keep their inheritances notwithstanding they should be pealed and driven shall be likened unto the wise virgins who took oil in their lamps, But all those who are unbelieving and fearful, will be likened unto the foolish virgins, who took no oil in their Lamps; and when they shall return, and say unto the saints, give us of your lands, behold there will be no room found for them. As respects giving \uline{deeds} I would advise to give deeds as far as the brethren have legal and Just claims for \sout{their} them and then let evry man answer to God for the disposal of them. I would suggest some Ideas to bro William P. not knowing as they will be of any real benefit, but suggest them for consideration I would be glad that he were here, but dare not advise, were it possable for him \uline{to} come, not knowing what shall befall us, as we are under very heavy and serious threatening from a great many people in this place.
					\hypertarget{EMS-1833-JS-10-DEC-1833-l}%
				But
					\marginnote{l}%
				purhaps, the people in Liberty may feel willing, God having power to soften the hearts of all men, to have a press established there; and if not, in some other place; any place where it can be the most convenient and it is possable to get to it: God will be willing to have it in any place where it can be practiculer and safe. we must be \uline{wise} as serpents and harmless as doves.
					\hypertarget{EMS-1833-JS-10-DEC-1833-m}%
				Again
					\marginnote{m}%
				I desire that bro William would collect all the information, and give us a true history of the begining and rise of Zion, and her calamities \&c Now hear the prayer of your unworthy Brothe[r] in the bonds of the new and everlasting covenant: O my God! thou who has called and chosen a few through thy weak instrument by commandment and sent them to Missouri a place which thou didst call \uline{Zion} and commanded thy servants to consecrate unto thyself for a place of \sout{a} refuge, and of safety for the gathering of thy saints, to be built up a holy city unto thyself and as thou hast said that none other place should be appointed like unto this
					\hypertarget{EMS-1833-JS-10-DEC-1833-n}%
				therefore
					\marginnote{n}%
				I ask thee in the name of Jesus Christ, to return thy people unto their homes, \& there inheritances, to enjoy the fruit of their Labors; that all the waste places may be built up; that all the enemies of thy people, who will not \sout{return} repent and \sout{re}turn unto thee be distroyed from off the face of that Land; and let an house be built and established unto thy name, and let all the losses that thy people have sustained be rewarded unto them, even more than fourfold;
					\hypertarget{EMS-1833-JS-10-DEC-1833-o}%
				that
					\marginnote{o}%
				the borders of Zion be enlarged forever, and let her be established no more to be thrown down; and let all thy saints when they are scattered as sheep, and are persecuted, \sout{and} flee unto Zion, and be established in the midst of her, and let her be organized according to thy law and let this prayer even \sout{before} be recorded before thy face; give thy holy spirit unto my brethren: unto whom I write: send thy angels to guard them and to deliver them from all evil;
					\hypertarget{EMS-1833-JS-10-DEC-1833-p}%
				and
					\marginnote{p}%
				when they turn there faces towards Zion and bow down before thee and pray may their sins never com[e] up before thy face neithe[r] have place in the book of thy remembrance and may they depart from all their eniquities. provide \sout{bread} <food> for them as thou doest for the ravens, provide clothing to cover there nakedness, and <houses that they may> \sout{cause that they} dwell therein give unto them friends in abundance, and let their names be recorded in the Lambs book of Life eternally before thy face Amen finely [finally], brethren, the grace of our Lord Jesus Christ be with you all unto his coming and Kingdom, Amen
				
				\begin{tightright}
				\uline{Joseph Smith J}
				\end{tightright}
				
		% -----
		\subsubsection{Revelation, 16--17 December 1833 [D\&C 101]}
		% -----
			\label{EMS-1833-JS-16-17-DEC-1833}
			% \label{EMS-1830-JS-DAY-3LETTERMONTH-YEAR}
		
			\begin{description}
				\item[Print Sources]
			\end{description}

				\begin{tabular}{p{0.5in}p{1in}p{0.5in}p{1in}}
				JSP	 	& D3:386--397		& JSR 		& 248--252 \\
				DC	 	& Section 101 		& HC 		& 1:458--464 \\
				HDDC 	& 1290--1307		& TCHC 		& 1:371--374 \\
				\end{tabular}
				% HDDC = woodford's DC diss; JSR = Marquardt's JS Revelations; TCHC = Vogel HC
			
			\begin{description}
				\item[Online Access] \href{https://www.josephsmithpapers.org/paper-summary/revelation-16-17-december-1833-dc-101/1}{Here}
				% \item[Analysis] \hyperref[DC101]{Here}
			\end{description}
			
			\begin{description}
				\item[Background]
			\end{description}
			
				% It might also be useful to give a two sentence context of the period: e.g., "This speech was given in Nauvoo to a small congregation one month prior to the destruction of the Nauvoo Expositor. These and other tensions may have been on the prophet's mind when he gave the speech."
			
			\begin{description}
				\item[Original Source]
			\end{description}
			
				% brief description of original source material, with reference to JSP or somewhere else for more info
				% Background material: it might be helpful to have a note that says whether a given document was presented as a speech, was written in a journal, etc.
				% Also a comment here about how reliable the source might be, or what shape it is in, if there are large omissions or parts that are hard to understand, and so on.
			
			\begin{description}
				\item[Text]
			\end{description}
			
				Verily I say unto you concerning your brethren who have been afflicted and persecuted and cast out from the Land of their inheriten[ces] I the Lord have suffered the affliction to come upon them wherewith they have been afflicted in consequence of their transgressions yet I will own them and they shall be mine in that day when I shall come to make up my jewels Therefore they must needs be chastened and tried even as Abraham who was commanded to offer up his only son for all those who will not endure chastening but deny me cannot be sanctified Behold I say unto you there were jar[r]ings and contentions envyings and strifes and lustful and covetous desires among them Therefore by these things they poluted their inheritances they were \sout{also} slow to hearken unto the voice of the Lord their God Therefore the Lord their God is slow to hearken unto their prayers to answer them in the day of their trouble In the day of their peace they esteemed lightly my council but in the day of their trouble of necessity they feel after me Verely I say unto you notwithstanding their sins my bowels are filled with compassion toward\sout{s} them I will not utterly cast them off and in the day of wrath I will remember mercy I have sworn and the decree hath gone forth by a former commandment which I have given unto you that I would let fall the sword of mine indignation in the behalf of my people and even as I have said it shall come to pass mine indignation is soon to be poured without measure upon all nations and this will I do when the cup of their eniquity is full and in that day all who are found upon the watch tower or in other words all mine Israel shall be saved and they that have been scattered shall be gathered and all they who have mourned shall be comforted and all they who have given their lives for my name shall be crowned Therefore let your hearts be comforted concerning Zion for all flesh is in mine hands be still and know that that I am God Zion shall not be moved out of her place notwithstanding her children are scattered they that remain and are pure in heart shall return and come to their inheritances they and their children with songs of everlasting joy to build up the waste places of Zion and all these things that the prophets might \sout{might} be fulfilled and behold there [is] none other place appointed than that which I have appointed neither shall there be any other place appointd that [than] that which I have appointed for the work of the gathering of my saints until the day cometh when there is found no more room for them and they shall be called stakes for the curtains of Zion or strength of Zion
				
				Behold it is my will that all they who call on my name and worship me according to mine everlasting gospel should gather to gether and stand in holy places and prepare for the revelation which is to come when the veil of the covering of my temple in my tabernacle which hideth the earth shall be taken off and all flesh shall see me together and evry coruptable thing both of man or of \sout{of} the beasts of the field or of the fowls of heaven or of the fish of the sea that dwell upon all the face \sout{of} of the earth shall be consumed And also that of element shall melt with fervant heat and all things shall become new that my knowledge and glory may dwell upon all the earth and in that day the enmity of man and the enmity of beasts yea the enmity of all flesh shall cease from before my face And in that day whatsoever any man \sout{ask} shall ask it shall be given unto him and in that day \sout{shall} satan shall not have power to tempt any man and there shall be no sorrow because there is no death In that day an infant shall not [die] until he is old and his life shall be as the age of a tree and when he dies he shall not sleep that is to say in the earth but shall be changed in the twinkling of an eye and shall be caught up and his rest shall be glorious yea verely I say unto you in that day when the Lord shall come he shall reveal all things things which have passed and hidden things which no man know things of the earth by which it was made and the purpose and the end thereof things most precious things that are above and things that are benieth things that are in the earth and upon the earth and in heaven All they that suffer persecution for my name and endure in faith though they are called to lay down their lives for my sake yet shall they partake of all this glory wherefore fear not even unto death for in this world your joy is not full but in me your joy is full therefore care not for the body neither for the Life of the body but care for the soul and for the Life of the soul and seek the face of the Lord always that in patience ye may possess your souls and ye shall have eternal life when men are called unto mine everlasting gospel and covenant with an everlasting covenant they are accounted as the salt of the earth and the savor of men they are called to be the savor of men there fore if that salt of the earth hath Lost its savor behold it is thenceforth good for nothing only to be cast out and troden under the feet of men Behold here is wisdom concerning the children of Zion even many but not all they were found transgressors therefore they must needs be chastened he that exalteth himself shall be abased and he that abaseth himself shall be exalted And now I will shew unto you a parable that you may know my will concerning the redemption of Zion A certain noblemen had a spot of Land very choice and he said unto his servants go ye into my vineyard even upon this very choice piecce of land and plant twelve olive trees and set watchmen round about them and build a tower that one may overlook the Land round about to be a watchman upon the tower that mine olive trees may not be broken down when the enemy shall come to spoil and take unto themselves the fruit of my vineyard Now the servants of this nobleman went and did as their Lord commanded them and planted the olive trees and built a hedge round about and set watchmen and began to build the tower and while they were yet Laying the foundation thereof they began to say among themselves and what need hath my Lord of this tower and consulted for a Long time saying among themselves what need hath my Lord of this tower seeing this is a time of peace might not this money be given to the exchangers for there is no need of these things and while they were at varience one with another they become very slothful and they harkened not unto the commandment of their Lord and the enemy came by night and broke down the hedge \sout{and} and the servants of the nobleman arose and were affrighted and fled and the enemy distroyed their works and broke down the Olive trees
				
				Now behold the nobleman the Lord of the vineyard called upon his servants and said unto them why what is the cause of this great evil ought ye not to have done even as I commanded you and after you had planted the vineyard and built the hedge round about and set watchmen upon the walls thereof built the tower also and set a watchman upon the tower and watched for my vineyard and not have fallen asleep lest the enemy should come upon you And behold the watchman upon the tower would have seen the enemy while he was yet afar off and then ye could have made ready and kept the enemy from breaking down the hedge thereof and saved my vineyard from the hands of the distroyer And the Lord of the vineyard said unto one of his servants, go and gather togethe[r] the residue of my servants and take all the strength of mine house which are my wariors my young men and they that are of middle age also among all my servants who are the strength of mine house save th$\diamond\diamond$ only whom I have appointed to tarry and go ye straitway unto the Land of my vineyard and redeem my vineyard for it is mine I have bought it with money therefore get ye straitway unto my Land break down the walls of mine enemies th[r]ow down their tower and scatte[r] their watchmen and inasmuch as they gather to gether against you avenge me of mine enemies that by and by I may come with the residue of mine house and possess the Land And the servant said unto his Lord when shall these things be And he said unto his servant when I will go ye strait way and do all things whatsoever I have commanded you \sout{and} and this shall be my seal and blessing upon you a faithful and wise steward in the midst of mine house a ruler in my kingdom--- And his servant went straitway and done all things whatsoever his Lord commanded him and afte[r] many days all things were fulfilled
				
				And again verely I say unto you I will shew unto you wisdom in me concerning all the churches in asmuch as they are willing to be guided in a right and proper way for their salvation that the work of the gathering to gether of my saints may continue that I may build them up unto my name upon holy places for the time of harvist is come and my word must needs be fulfilled therefore I must gather to gether my people according to the parable of the wheat and the tares that the wheat may be secured in the garner to possess eternal life and be crowned with celestial glory when I come in the Kingdom of my fathe[r] to reward evry man according as his work shall be whilst the tares shall be bound in bundles and their bands made strong that they may be burned with unquenchable fire, therefore a commandment I give unto all the churches that they shall continue to gather to gether into the places which I have appointed, nevertheless as I have said unto you in a former commandment let not your gathering be in haste nor by flight but let all things be prepared before you and in order that all things be prepared before you observe the commandments which I have given concerning these things which saith or teacheth to purchace all the Land by mony which can be purchaced for mony in the regions round about the Land which I have appointed to be the Land of Zion for the begining of the gathering of my saints all the Land which can be purchaced in Jackson County and the counties round about and leave the residue in mine hand Now verily I say unto you let all the churches gather to gether all their monies let these things be done in their time lo not in haste and observe to have all things prepared before you and let honorable men be appointed even wise men and send them to purchace the lands and every church in the eastern countries when they are built up if they will harken unto the council they may buy Lands and gather together upon them and in this way they may establish Zion there is even now already \sout{enough} in store a sufficient yea even abundence to redeem Zion and establish her waste places no more to be thrown down were the churches who call themselves after my name willing to \sout{hear} harken to my voice and again I say unto you those who have been scattered by their enemies it is my will that they should continue to importune for redress and redemption by the hand of those who are placed as rulers and are in authority over you according to the Law and constitution of the people which I have suffered to be established and should be maintained for the rights and protection of all flesh according to just and holy principles, that evry man may act in doctrine and principle pertaining to futurity according to the moral agency which I have given unto them that evry man may be accountable for his own \sout{sin in} sins [in] the day of judgment therefor it is not right that any man should be in bondage one to another and for this purpose have I established the constitution of this Land by the hands of wise men whom I raised up unto this very purpose and redeemed the Land by the shedding of blood, Now unto what shall I liken the children of Zion I will liken them \sout{to} unto the parable of the woman and the unjust judge (for men ought always to pray and not faint) which saith there was in a city a judge which feared not God neither regarded man and there was a widow in that city and she came unto him saying avenge me of mine adver[s]ary and he would not for a while but afterward he said within himself though I fear not God nor regard man yet because this widow troubleth me I will avenge her lest by her continual coming she weary me thus will I liken the children of Zion let them impertune at the feet of the judge and if he heed them not let them impertune at the feet of the Govoner and if the Govoner heed them not let them importune at the feet of the President and if the President heed them not then will the Lord arise and come forth out of his <hiding> place \& in his fury vex the nation and in his hot displeasure and in his fierce ander [anger] in his time will cut off these wicked unfaithful and unjust stewards and appoint them their portion among hypocrits and unbelievers even in outer darkness where there is weeping and wailing and gnashing of teeth pray ye therefore that their ears may be opened \sout{to} unto your cries that I may be merciful unto them that these things come not upon them what I have said unto you must needs be that all men may be left without excuse that wise men and rulers may hear and know that which \sout{I} <they> have never considered that I may procede to bring to pass my act my strange act and perform my work my strange work that men may desern between the righteous and the wicked saith your God and again I say unto you it is conterary to my commandment and my will that my servant Alge[r]non Sidney Gilbert should sell my store house which I have appointed unto my people \sout{unto} <into> the hands of mine enemies let not that which I have appointed be poluted by mine enemies by the consent of those who call themselves afte[r] my name for this is a very soar and grievous sin against me and against my people in consequence of those things which I have decreed and are soon to <be>fall the nations therefore it is my will that my people should claim and hold claim upon that which I have appointed unto them though they should not be permited to dwell thereon nevertheless I do not say they shall not dwell thereon for in as much as they bring forth fruit and works meet for my kingdom they shall dwell thereon they shall build and anothe[r] shall not inherit it they shall plant vineyards and they shall eat the fruit thereof even so amen
				
		% -----
		\subsubsection{Letter to the Church, not after 18 December 1833}
		% -----
			\label{EMS-1833-JS-18-DEC-1833-BEFORE}
			% \label{EMS-1830-JS-DAY-3LETTERMONTH-YEAR}
		
			\begin{description}
				\item[Print Sources]
			\end{description}

				\begin{tabular}{p{0.5in}p{1in}p{0.5in}p{1in}}
				JSP	 	& D3:397--401		& JSR 		& --- \\
				DC	 	& --- 		& HC 		& --- \\
				HDDC 	& ---		& TCHC 		& --- \\
				TPJS 	& 42--43			& 	 		&  \\
				\end{tabular}
				% HDDC = woodford's DC diss; JSR = Marquardt's JS Revelations; TCHC = Vogel HC
			
			\begin{description}
				\item[Online Access] \href{https://www.josephsmithpapers.org/paper-summary/letter-to-the-church-not-after-18-december-1833/1}{Here}
				% \item[Analysis] \hyperref[XX]{Here}
			\end{description}
			
			\begin{description}
				\item[Background]
			\end{description}
			
				% It might also be useful to give a two sentence context of the period: e.g., "This speech was given in Nauvoo to a small congregation one month prior to the destruction of the Nauvoo Expositor. These and other tensions may have been on the prophet's mind when he gave the speech."
			
			\begin{description}
				\item[Original Source]
			\end{description}
			
				% brief description of original source material, with reference to JSP or somewhere else for more info
				% Background material: it might be helpful to have a note that says whether a given document was presented as a speech, was written in a journal, etc.
				% Also a comment here about how reliable the source might be, or what shape it is in, if there are large omissions or parts that are hard to understand, and so on.
			
			\begin{description}
				\item[Text]
			\end{description}
			
				\begin{tightcenter}
						\hypertarget{EMS-1833-JS-18-DEC-1833-BEFORE-a}%
					THE
						\marginnote{a}
					ELDERS IN KIRTLAND, TO THEIR BRETHREN ABROAD.
				\end{tightcenter}
				
				\textit{Dear Brethren in Christ, and companions in tribulation}:
				
				It seemeth good unto us, to drop a few lines to you, giving you some instruction relative to conducting the affairs of the kingdom of God, which has been committed unto us in these later times, by the will and testament of our Mediator, whose intersessions in our behalf, are lodged in the bosom of the Eternal Father, and ere long will burst with blessings upon the heads of all the faithful:
				
					\hypertarget{EMS-1833-JS-18-DEC-1833-BEFORE-b}%
				We
					\marginnote{b}%
				have all been children, and are too mutch so at the present time; but we hope in the Lord, that we may grow in grace and be prepared for all things which the bosom of futurity may disclose unto us. Time is rapidly rolling on, and the prophecies must be fulfilled. The days of tribulation are fast approaching, and the time to test the fidelity of the Saints, has come.--- Rumor with her ten thousand tongues is diffusing her uncertain sounds in almost every ear: but in these times of sore trial, let the saints be patient and see the salvation of God. Those who cannot endure persecution and stand in the day of affliction, cannot stand in the day when the Son of God shall burst the veil, and appear in all the glory of his Father with the holy angels.
				
					\hypertarget{EMS-1833-JS-18-DEC-1833-BEFORE-c}%
				On
					\marginnote{c}%
				the subject of ordination, a few words are necessary: In many instances there has been too much haste in this thing, and the admonition of Paul has been too slightingly passed over, which says, ``\textit{Lay hands suddenly upon no man}.'' Some have been ordained to the ministry, and have never acted in that capacity, or magnified their calling, at all: Such may expect to lose their calling, except they awake and magnify their office. Let the elders abroad be exceedingly careful upon this subject, and when they ordain a man to the holy ministry, let it be a \textit{faithful man}, who is able to teach others also; that the cause of Christ suffer not. It is not the multitude of preachers that is to bring about the glorious Millenium! but it is those who are ``\textit{called, and chosen, and faithful}.''
				
					\hypertarget{EMS-1833-JS-18-DEC-1833-BEFORE-d}%
				Let
					\marginnote{d}%
				the elders be exceedingly careful about \textit{unnecessarily} disturbing and harrowing up the feelings of the people. Remember, that your business is, to preach the gospel in all humility and meekness, and warn sinners to repent and come to Christ. Avoid contentions and vain disputes with men of corrupt minds, who do not desire to know the truth. Remember that ``\textit{it is a day of warning, and not a day of many words}.'' If they receive not your testimony in one place, flee to another, remembering, to cast no reflections, nor throw out any bitter sayings. If you do your duty, it will be just as well with you, as though all men embraced the gospel.
				
					\hypertarget{EMS-1833-JS-18-DEC-1833-BEFORE-e}%
				Be
					\marginnote{e}%
				careful about sending boys to preach the gospel to the world; if they go, let them be accompanied by some one who is able to guide them in the proper channel, lest they become puffed up, and fall under condemnation and into the snare of the devil: finally, in these critical times, be careful; call on the Lord day and night. Beware of pride: Beware of \textit{false brethren}, who will creep in among you to spy out your liberties, \&c. Awake to righteousness and sin not; let your light shine, and show yourselves workmen that need not be ashamed, rightly dividing the word of truth. Apply yourselves diligently to study, that your minds may be stored with all necessary information.
				
					\hypertarget{EMS-1833-JS-18-DEC-1833-BEFORE-f}%
				We
					\marginnote{f}%
				remain your brethren in Christ, anxiously praying for the day of redemption to come, when iniquity shall be swept from the earth; and everlasting righteousness brought in:
				
				Farewell.
				
		% -----
		\subsubsection{Record, 18 December 1833}
		% -----
			\label{EMS-1833-JS-18-DEC-1833}
			% \label{EMS-1830-JS-DAY-3LETTERMONTH-YEAR}
		
			\begin{description}
				\item[Print Sources]
			\end{description}

				\begin{tabular}{p{0.5in}p{1in}p{0.5in}p{1in}}
				JSP	 	& J1:21--24	& JSR 		& --- \\
				DC	 	& ---	& HC 		& --- \\
				HDDC 	& ---	& TCHC 		& --- \\
				TPJS 	& 38--41 &  &  \\
				\end{tabular}
				% HDDC = woodford's DC diss; JSR = Marquardt's JS Revelations; TCHC = Vogel HC
			
			\begin{description}
				\item[Online Access] \href{https://www.josephsmithpapers.org/paper-summary/journal-1832-1834/31}{Here}
				% \item[Analysis] \hyperref[XX]{Here}
			\end{description}
			
			\begin{description}
				\item[Background]
			\end{description}
			
				% It might also be useful to give a two sentence context of the period: e.g., "This speech was given in Nauvoo to a small congregation one month prior to the destruction of the Nauvoo Expositor. These and other tensions may have been on the prophet's mind when he gave the speech."
			
			\begin{description}
				\item[Original Source]
			\end{description}
			
				% need a discussion here of how the actual text of these blessings was recorded later; see https://www.josephsmithpapers.org/paper-summary/appendix-5-document-1-blessing-to-joseph-smith-sr-and-lucy-mack-smith-between-circa-15-and-28-september-1835/1#historical-intro , and Marquardt's 2007 Early Patriarchal Blessings book (the first set of blessings, they start on page 3). It would probably be a good idea to create a new entry on the same date for all that stuff in early patriarchal blessings.
			
				% brief description of original source material, with reference to JSP or somewhere else for more info
				% Background material: it might be helpful to have a note that says whether a given document was presented as a speech, was written in a journal, etc.
				% Also a comment here about how reliable the source might be, or what shape it is in, if there are large omissions or parts that are hard to understand, and so on.
			
			\begin{description}
				\item[Text]
			\end{description}
			
					\hypertarget{EMS-1833-JS-18-DEC-1833-a}%
				1833
					\marginnote{a}%
				Dec. 18 This day the Elders assembled togeth[er] in the printing office and then proceded to bow down before the Lord and dedicate the printing press and all that pertains thereunto to God by mine own hand and confirmed by bro Sidney Rigdon and Hyrum Smith and then proceded to take the first proof sheet of the star edited by Bro Oliv Cowdy [Oliver Cowdery] blessed of the Lord is bro Oliver nevertheless there are are two evils in him that he must needs forsake or he cannot altogeth[er] escape the buffettings of the advers[ar]y
					\hypertarget{EMS-1833-JS-18-DEC-1833-b}%
				if
					\marginnote{b}%
				he shall forsak these evils he shall be forgiven and shall be made like unto the bow which the Lord hath set in the heavens he shall be a sign and an ensign unto the nations behold he is blessed of the Lord for his constancy and steadfastness in the work of the Lord wherefore he shall be blessed in his generation and they shall never be cut off and he shall be helped out of many troubles and if he keep the \sout{commandmend} <commandments> and harken unto the <council of the> Lord \sout{his and} his rest shall be glorious and again blessed of the Lord is my father and also my mother and my brothers and my sisters for they shall yet find redemption in the house of the Lord and their of[f]springs shall be a blessing a Joy and a comfort unto them
					\hypertarget{EMS-1833-JS-18-DEC-1833-c}%
				blessed
					\marginnote{c}%
				is my mother for her soul is ever fill[ed] with benevolence and phylanthropy and notwithstanding her age yet she shall receive strength and shall be comferted in the midst of her house and she shall have eternal life and blessed is my father for the hand of the Lord shall be over him for he shall see the affliction <of his children> pass away and when his head is fully ripe he shall behold himself as an olive tree whose branches are bowed down with much fruit he shall also possess a mansion on high blessed of the Lord is my brothe[r] Hyram for the integrity of his heart he shall be girt about with truth and faithfulness shall be the strength of his loins from generation to generation he shall be a shaft in the hand of his God to exicute Judgment upon his enemies and he shall be hid by the hand of the Lord that none of his secret parts shall be discovered unto his hu[r]t his name shall be accounted a blessing among men
					\hypertarget{EMS-1833-JS-18-DEC-1833-d}%
				and
					\marginnote{d}%
				when he is in trouble and great tribulation hath come upon him he shall remember the God of Jacob and he will shield him from the power of satan and he shall receive \sout{councl} <councel> in the house of the most high that he may be streng[t]hened in hope that the goings <of his feet> may be established for ever blessed of the Lord is bro Samuel [Smith] because the Lord shall say unto him Sam\supers{l.}, Sam\supers{l.}, therefore he shall be made a teache[r] in the hous of the Lord and the Lord shall mature his mind in Judgment and thereby he shall obtain the esteem and fellowship of his brethren and his soul shall be established and he shall benefit th[e] house of the Lord because he shall obtain answer to prayer in his faithfulness---
					\hypertarget{EMS-1833-JS-18-DEC-1833-e}%
				Bro
					\marginnote{e}%
				William [Smith] is as the firce Lion who devideth not the spoil because of his strength and in the pride of his heart he will neglect the more \sout{weightier} weighty matters until his soul is bowed down in sorrow and then he shall return and call on th[e] name of his God and shall find forgivness and shall wax valient
					\hypertarget{EMS-1833-JS-18-DEC-1833-f}%
				therefor
					\marginnote{f}%
				he shall be saved unto the utter most and as the roaring Lion of the forest in the midst of his prey so shall the hand of his generation be lifted up against those who are set on high that fight against the God of Israel fearless and unda[u]nted shall they be in battle in avenging the rongs of th[e] innocent and relieving the oppressed therfor the blessings of the God of Jacob shall be in the midst of his house notwithstanding his rebelious heart 
					\hypertarget{EMS-1833-JS-18-DEC-1833-g}%
				and
					\marginnote{g}%
				<now> O God let the residue of my fathers house ever come up in remembrance before thee that thou mayest save them from the hand of the oppressor and establish their feet upon the rock of ages that they may have place in thy house and be saved in thy Kingdom and let all these things be even as I have said for Christs sake Amen
				
		% -----
		\subsubsection{Minutes, 26 December 1833}
		% -----
			\label{EMS-1833-JS-26-DEC-1833}
			% \label{EMS-1830-JS-DAY-3LETTERMONTH-YEAR}
		
			\begin{description}
				\item[Print Sources]
			\end{description}

				\begin{tabular}{p{0.5in}p{1in}p{0.5in}p{1in}}
				JSP	 	& D3:401--402		& JSR 		& --- \\
				DC	 	& --- 				& HC 		& 1:469--470 \\
				HDDC 	& ---				& TCHC 		& 1:378--379 \\
				\end{tabular}
				% HDDC = woodford's DC diss; JSR = Marquardt's JS Revelations; TCHC = Vogel HC
			
			\begin{description}
				\item[Online Access] \href{https://www.josephsmithpapers.org/paper-summary/minutes-26-december-1833/1}{Here}
				% \item[Analysis] \hyperref[XX]{Here}
			\end{description}
			
			\begin{description}
				\item[Background]
			\end{description}
			
				% It might also be useful to give a two sentence context of the period: e.g., "This speech was given in Nauvoo to a small congregation one month prior to the destruction of the Nauvoo Expositor. These and other tensions may have been on the prophet's mind when he gave the speech."
			
			\begin{description}
				\item[Original Source]
			\end{description}
			
				% brief description of original source material, with reference to JSP or somewhere else for more info
				% Background material: it might be helpful to have a note that says whether a given document was presented as a speech, was written in a journal, etc.
				% Also a comment here about how reliable the source might be, or what shape it is in, if there are large omissions or parts that are hard to understand, and so on.
			
			\begin{description}
				\item[Text]
			\end{description}
			
					\hypertarget{EMS-1833-JS-26-DEC-1833-a}%
				This
					\marginnote{a}%
				day at evening, a Bishops Court was called to take into consideration the case of Bro. Ezekiel Rider an Elder of the Church, who had \sout{brought} said many hard things against Bro [Newel K.] Whitney, the Bishop of the church--- he said that Bro. Whitney was not fit for a Bishop and that he treated the Brethren, who came into the Store, with disrespect that he was overbearing and fain would walk on the necks of the Bretheren \&c------
				
					\hypertarget{EMS-1833-JS-26-DEC-1833-b}%
				Bro.
					\marginnote{b}%
				Story was also in a Similar transg[r]ession. They were both rebuked sharply by bro. Sidney [Rigdon] \& Bro Joseph who told them that this church must feel the wrath of God except they repent of their Sins, cast away their murmurings and complainings one of another. \&c. \&c. Bro. Rider \& Bro. Story confessed their wrongs. and all forgave one another \& closed by praying to the Lord for his blessings to rest upon us
				
				Kirtland Dec 26. 1833
				
				\begin{tightright}
				Orson Hyde, Clk
				\end{tightright}